%\documentclass{acm_proc_article-sp}
\documentclass{sig-alternate}
\usepackage{algorithm,algpseudocode,cite,amsmath,amssymb,amsfonts,dsfont,multirow,multicol,enumitem,epsfig,url,array,makecell,balance}
%\usepackage{algorithm,algpseudocode,cite,amsmath,amssymb,amsfonts,dsfont,multirow,multicol,enumitem,epsfig,array,makecell,balance,hyperref,breakurl}
\usepackage{times}
%\usepackage[dvipdfm]{hyperref}

\newtheorem{definition} {Definition}
\newtheorem{example} {Example}
\newtheorem{theorem} {Theorem}
\newtheorem{lemma} {Lemma}
\newtheorem{problem} {Problem}
\newtheorem{observation} {Observation}
\newtheorem{proposition} {Proposition}
\newtheorem{assumption} {Assumption}
\newtheorem{property} {Property}
\newcommand{\proofsketch}{\noindent{\em Proof (Sketch): }}
\def\done{\hspace*{\fill} {$\square$}}

\def\header{\vspace{2mm} \noindent}
%\def\header{\noindent}
\def\extraspacing{\vspace{2mm}}
\def\myspacing{\vspace{2mm}}
\def\figcapup{\vspace{-2mm}}
\def\figcapdown{\vspace{-3mm}}
\def\tblcapup{\vspace{-0mm}}
\def\tblcapdown{\vspace{-3mm}}
\def\tbldown{\vspace{-0mm}}
\def\tblup{\vspace{2mm}}
\def\codecapup{\vspace{-3mm}}
\def\codecapdown{\vspace{-2mm}}

\def\vgap{\vspace{0mm}}

%\itemsep = 0mm
%\topsep = 0mm
%\partopsep = 0mm
%\parsep = 0mm
%\parskip = -4mm

%\setlength{\parskip}{0cm}

%\topsep, \partopsep, \itemsep, \parsep, \itemindent, \labelsep, \labelwidth, \makelabel, \jot, \arraycolsep, \tabcolsep, \arrayrulewidth, \doublerulesep, \arraystretch.


% ----------------------------------
% HELPFUL MACROS
\def\r{\alpha}
\def\H{H^*}
\def\G{G^*}
\def\E{E^*}
\def\GC{G_c}
\def\VV{V^c}
\def\EE{E^c}
\def\d{dist}
\def\a{\lambda}
\def\v{v^*}
\def\GF{G_f}
\def\GB{G_b}
%\def\GR{G_r}
\def\FF{F_f}
\def\FB{F_b}
%\def\FR{F_r}
\def\HF{H_f}
\def\HB{H_b}
\def\QF{Q_f}
\def\QB{Q_b}
\def\kf{\kappa_f}
\def\kb{\kappa_b}
%\def\QFF{Q^*_f}
%\def\QBB{Q^*_b}
%\def\dm{d_{min}}
\def\QC{Q_c}
\def\P{P^*}
\def\PA{P^\circ}

%\def\PAR{7pt}

% ----------------------------------

\begin{document}
\begin{sloppy}

\title{Efficient Single-Source Shortest Path and Distance Queries on Large Graphs}

\numberofauthors{1}

\author{
\alignauthor
Andy Diwen Zhu $\qquad$ Xiaokui Xiao $\qquad$ Sibo Wang $\qquad$ Wenqing Lin \\
\affaddr{
%\begin{tabular}{c}
%$^\dag$School of Computer Engineering & $\qquad$ & $^\S$School of
%Computer Science \\
%Nanyang Technological University & & Fudan University  \\
%Singapore & & China
School of Computer Engineering \\
Nanyang Technological University\\
Singapore
}\\
\email{
\{dwzhu, xkxiao, swang, wlin\}@ntu.edu.sg
}
}

\maketitle
\begin{abstract}
%This paper investigates two types of graph queries: single source distance (SSD) queries and single source shortest path (SSSP) queries. Given a node v in a graph G, an SSD query from v asks for the distance from v to every other node in G, while an SSSP query retrieves the shortest path from v to any other node. These two types of queries find important applications in graph analysis, especially in the computation of graph measures. Most of the existing solutions for SSD and SSSP queries, however, require that the input graph fits in the main memory, which renders them inapplicable for the massive disk-resident graphs commonly used in web and social applications. There are several techniques that are designed to be I/O efficient, but they all focus on undirected and/or unweighted graphs, and they provide rather sub-optimal query efficiency.
%
%To address the deficiency of existing work, this paper presents Highways-on-Disk (HoD), a disk-based index that supports both SSD and SSSP queries on directed and weighted graphs. The key idea of HoD is to augment the input graph with a set of auxiliary edges, and exploit them during query processing to reduce I/O and computation costs. We experimentally evaluate HoD on a variety of real-world graphs with up to billions of nodes and edges, and we demonstrate that HoD significantly outperforms alternative solutions in terms of query efficiency.


This paper investigates two types of graph queries: {\em single source distance (SSD)} queries and {\em single source shortest path (SSSP)} queries. Given a node $v$ in a graph $G$, an SSD query from $v$ asks for the distance from $v$ to any other node in $G$, while an SSSP query retrieves the shortest path from $v$ to any other node. These two types of queries are fundamental building blocks of numerous graph algorithms, and they find important applications in graph analysis, especially in the computation of graph measures. Most of the existing solutions for SSD and SSSP queries, however, require that the input graph fits in the main memory, which renders them inapplicable for the massive disk-resident graphs commonly used in web and social applications. The only exceptions are a few techniques that are designed to be I/O efficient, but they all focus on undirected and/or unweighted graphs, and they only offer sub-optimal query efficiency.

To address the deficiency of existing work, this paper presents {\em Highways-on-Disk (HoD)}, a disk-based index that supports both SSD and SSSP queries on directed and weighted graphs. The key idea of HoD is to augment the input graph with a set of auxiliary edges, and exploit them during query processing to reduce I/O and computation costs. We experimentally evaluate HoD on both directed and undirected real-world graphs with up to billions of nodes and edges, and we demonstrate that HoD significantly outperforms alternative solutions in terms of query efficiency.


%Given two locations $s$ and $t$ in a road network, a distance query returns the minimum network distance from $s$ to $t$, while a shortest path query computes the actual route that achieves the minimum distance. These two types of queries find important applications in practice, and a plethora of solutions have been proposed in past few decades. The existing solutions, however, are optimized for either practical or asymptotic performance, but not both. In particular, the techniques with enhanced practical efficiency are mostly heuristic-based, and they offer unattractive worst-case guarantees in terms of space and time. On the other hand, the methods that are worst-case efficient often entail prohibitive preprocessing or space overheads, which render them inapplicable for the large road networks (with millions of nodes) commonly used in modern map applications.
%
%This paper presents {\em Arterial Hierarchy (AH)}, an index structure that narrows the gap between theory and practice in answering shortest path and distance queries on road networks. On the theoretical side, we show that, under a realistic assumption, AH answers any distance query in $\tilde{O}(\log \r)$ time, where $\r = d_{max}/d_{min}$, and $d_{max}$ (resp.\ $d_{min}$) is the largest (resp.\ smallest) $L_\infty$ distance between any two nodes in the road network. In addition, any shortest path query can be answered in $\tilde{O}(k + \log \r)$ time, where $k$ is the number of nodes on the shortest path. On the practical side, we experimentally evaluate AH on a large set of real road networks with up to twenty million nodes, and we demonstrate that (i) AH outperforms the state of the art in terms of query time, and (ii) its space and pre-computation overheads are moderate.
\end{abstract}

%In addition, AH requires only $O(h n)$ space, and it incurs $O(h n^2)$ pre-computation cost.

%============================================

%% [inlined 1-Introduction/introduction.tex]
\section{Introduction} \label{sec:intro}

%This paper considers two types of graph queries, namely, {\em single source distance (SSD)} queries and {\em single source shortest path (SSSP)} queries.

Given a graph $G$, a {\em single source distance (SSD)} query from a node $v \in G$ asks for the distance from $v$ to any other node in $G$. Meanwhile, a {\em single source shortest path (SSSP)} query retrieves the shortest path from $v$ to any other node. These two types of queries find important applications in graph analysis \cite{ckcc12}, especially in the computation of graph measures \cite{acim99,bkmm07,ew04,tk11}. For example, the estimation of {\em closeness} measures \cite{ew04} on a graph $G$ requires performing SSD queries from a large number of nodes in $G$, while the approximation of {\em betweenness} measures \cite{bkmm07} requires executing numerous SSSP queries.

%These two types of queries find numerous applications in practice, including map services \cite{rt10}, graph analysis \cite{ut05}, web search \cite{vfdg07}, etc. For instance, SSD queries are a primitive in the existing solutions for computing various important graph measures, such as {\em diameter} \cite{acim99,dflo06}, {\em centrality} \cite{ew04}, and {\em eccentricity} \cite{tk11}. As another example, PPD queries play a significant role in social search, where the ranking of search results highly depends on the distance between users in the social graph \cite{pbcg09}.

%\parskip = \PAR
\vgap

%\footnote{On unweighted graph, SSD and SSSP queries can also be answered using breadth first search.}

The classic solution for SSD and SSSP queries is Dijkstra's algorithm \cite{d59}. Given a SSD or SSSP query from a node $s$, Dijkstra's algorithm traverses the graph starting from $s$, such that the nodes in $G$ are visited in ascending order of their distances from $s$. Once a node $v$ is visited, the algorithm returns the distance from $s$ to $v$ based on the information maintained during the traversal; the shortest path from $s$ to $v$ can also be efficiently derived if required.

\vgap

A plethora of techniques have been proposed to improve over Dijkstra's algorithm for higher query efficiency (see \cite{dss09,s12} for surveys). Although those techniques all require pre-processing the given graph (which incurs extra overhead compared with Dijkstra's algorithm), the pre-computation pays off when the number of queries to be processed is large, as is often the case in graph analysis. Nevertheless, most of the existing techniques assume that the given graph fits in the main memory (for pre-computation and/or query processing), which renders them inapplicable for the massive disk-resident graphs commonly used in web and social applications. There are a few methods \cite{mo09,mz03,mz12,m09,ks96} that address this issue by incorporating Dijkstra's algorithm with I/O-efficient data structures, but the performance of those methods are shown to be insufficient for practical applications \cite{ckcc12}. The main reason is that, when Dijkstra's algorithm traverses the graph, the order in which it visits nodes can be drastically different from the order in which the nodes are arranged on the disk. This leads to a significant number of random disk accesses, which results in poor query performance.

\vgap

In contrast to the aforementioned techniques, Cheng et al.\ \cite{ckcc12} propose the first practically efficient index (named {\em VC-Index}) for SSD and SSSP queries on disk-resident graphs. The basic idea of VC-Index is to pre-compute a number of {\em reduced} versions of the input graph $G$. Each reduced graph contains some relatively important nodes in $G$, as well as the distances between some pairs of those nodes. During query processing, VC-Index scans a selected subset of reduced graphs, and then derives query results based on the pre-computed distances. Compared with those methods based on Dijkstra's algorithm \cite{mo09,mz03,mz12,m09,ks96}, VC-Index is more efficient as it only performs sequential reads on disk-resident data.

\vgap

\header
{\bf Motivation and Contribution.} All existing disk-based solutions for SSD and SSSP queries \cite{mo09,mz03,mz12,m09,ks96} require that the input graph is undirected, which renders them inapplicable for any application built upon directed graphs. This is rather restrictive as numerous important types of graphs (e.g., road networks, web graphs, social graphs) are directed in nature. Furthermore, even when the input graph is undirected, the query efficiency of the existing solutions is less than satisfactory. In particular, our experiments (in Section~\ref{sec:exp}) show that VC-Index, albeit being the state of the art, requires tens of seconds to answer a single SSD query on a graph with less than $100$ million edges, and needs more than two days to estimate the closeness measures on the same graph.

%though VC-Index is more efficient than other existing solutions, it has several deficiencies that restrict its applications in practice. First, VC-Index is only applicable on undirected graphs, while many important types of graph data (e.g., road networks and web graphs) are directed in nature. Second, VC-Index cannot efficiently support PPD queries. Specifically, if we are to use VC-Index to compute the distance between two nodes $s$ and $t$, we need to first perform an SSD query from $s$ (or $t$), and then retrieve the distance between $s$ and $t$ from the result of the SSD query. This is inefficient since, intuitively, processing the SSD query requires visiting a large number of nodes that are irrelevant to the answer of the PPD query, which entails unnecessary query cost. Third, the query efficiency of VC-Index is rather sub-optimal: Our experiments show that it takes tens of seconds to answer SSD queries even on moderate-size graphs.

%VC-Index incurs significant space overhead, and its query efficiency leaves much room for improvement. In particular, our experiments show that VC-Index may require ten times more space than the input data, and it takes tens of seconds to answer SSD queries even on moderate-size graphs.

%\begin{figure}[t]
%\centering
%\begin{small}
%\begin{tabular}{cc}
%\multicolumn{2}{c}{\hspace{0mm} \includegraphics[height=3mm]{./figures/legend1.eps}} \vspace{1mm} \\
%\hspace{0mm} \includegraphics[width=32mm]{./figures/original.eps} & \hspace{5mm}
%\includegraphics[width=32mm]{./figures/hod.eps} \\
%\hspace{0mm}  (a) Graph $G$ & \hspace{3mm} (b) Augmented Graph $\G$
%\end{tabular}
%\figcapup  \caption{Illustration of HoD.} \figcapdown
%\label{fig:intro-hod}
%\end{small}
%\end{figure}



\vgap

To address the deficiency of existing work, this paper proposes {\em Highways-on-Disk (HoD)}, a disk-based index that supports both SSD and SSSP queries on directed and weighted graphs. The key idea of HoD is to augment the input graph with a set of auxiliary edges (referred to as {\em shortcuts} \cite{ss05}), and exploit them during query processing to reduce I/O and computation costs. For example, Figure~\ref{fig:over-steps}a illustrates a graph $G$, and Figure~\ref{fig:over-steps}e shows an augmented graph $\G$ constructed from $G$. $\G$ contains three shortcuts: $\langle v_8, v_9\rangle$, $\langle v_9, v_7\rangle$, and $\langle v_9, v_{10}\rangle$. Each shortcut has the same length with the shortest path connecting the endpoints of the shortcut. For example, the length of $\langle v_8, v_9\rangle$ equals $2$, which is identical to the length of the shortest path from $v_8$ to $v_9$. Intuitively, the shortcuts in $\G$ enable HoD to efficiently traverse from one node to another (in a manner similar to how highways facilitate traversal between distant locations). For instance, if we are to traverse from $v_1$ to $v_{10}$ in $\G$, we may follow the path $\langle v_1, v_9, v_{10}\rangle$, which consists of only three nodes; in contrast, a traversal from $v_1$ to $v_{10}$ in $G$ would require visiting five nodes: $v_1$, $v_9$, $v_6$, $v_7$, and $v_{10}$.

\vgap

In general, when HoD answers an SSD or SSSP query, it often traverses the augmented graph via shortcuts (instead of the original edges in $G$). We show that, with proper shortcut construction and index organization, the query algorithm of HoD always traverses nodes in the same order as they are arranged in the index file. Consequently, HoD can answer any SSD or SSSP query with a linear scan of the index file, and its CPU cost is linear to the number of edges in the augmented graph. %Furthermore, for PPD queries, HoD can optimize the linear scan to avoid visiting irrelevant nodes in the augmented graph, which considerably reduces query costs.
We experimentally evaluate HoD on a variety of real-world graphs with up to $100$ million nodes and $3$ billion edges, and we demonstrate that HoD significantly outperforms VC-Index in terms of query efficiency. In particular, the query time of HoD is smaller than that of VC-Index by up to two orders of magnitude. Furthermore, HoD requires a smaller space and pre-computation time than VC-Index in most cases.

\vgap



%The remainder of the paper is organized as follows.

%nodes in the graph in a {\em best-first} manner, i.e., it always traverse
%
%Define SSSP and P2P queries. Talk about their applications.
%
%Talk about the existing approaches and say that they are main memory based.
%
%Talk about James' method and point out its problems: (i) restricted to undirected graphs, (ii) sub-optimal query efficiency, and (iii) does not handle graphs with relatively large degrees.
%
%Brag about HoD. Talk about ordering and shortcuts. Draw the similarity between shortcuts and highways. Emphasize the selling points of HoD: (i) applicable for both directed and undirected graphs, (ii) finishes any ss-dist query with two scans of the index structure, and (iii) efficiently handles the case when the graph has vertices with large degrees.
%
%To the best of our knowledge, this is the first index that demonstrates efficiency on graphs with size over XXX vertices.




%Given two locations $s$ and $t$ in a road network, a distance query returns the network distance from $s$ to $t$, while a shortest path query computes the actual route that achieves the minimum distance. These two types of queries find important applications in map, navigation, and location-based services. To illustrate, consider that a user of a map service is looking for a nearby Italian restaurant for dinner. In response to the user's query, the service provider can first retrieve the list of Italian restaurants in the region close to the user's current location $u$. After that, the network distance from $u$ to each restaurant is computed (using a distance query), and those distances are returned to the user along with the list of restaurants. Then, if the user chooses a preferred restaurant $r$ from the list, the service provider can employ a shortest path query to provide the user with driving directions from $u$ to $r$.
%
%The classic solution for shortest path and distance queries is Dijkstra's algorithm \cite{d59}. It traverses the road network nodes in ascending order of their distances from $s$; once it reaches $t$ during the traversal, it can compute the distance from $s$ to $t$ and can retrieve the shortest path based on the information recorded before $t$ is visited. With proper data structures, Dijkstra's algorithm runs in $O(n \log n + m)$ time for any shortest path or distance query, where $n$ (resp.\ $m$) is the number of nodes (resp.\ edges) in the road network. Albeit simple and elegant, Dijkstra's algorithm is inefficient for sizable road networks, as it requires traversing all network nodes that are closer to $s$ than $t$, which incurs a significant overhead when $s$ and $t$ are far part.
%
%A plethora of techniques \cite{ssa08,ssa09,ss10,gss08,bfm06,tsp11,rt10,gh05,gkw06,kmw10,hkm06,dss09,fr06,afg10,bfs07,jp02,jhr98,gkr04,ms12} have been proposed to improve over Dijkstra's algorithm in terms of either practical efficiency or asymptotic bounds. The existing methods that focus on practical performance mostly rely on heuristics, and hence, their asymptotic bounds are unattractive in general. For instance, the best heuristic approach by Geisberger et al.\ \cite{gss08} answers shortest path or distance queries in at most a few milliseconds even on road networks with millions of nodes, but its space and time complexities are both $O(n^2)$, i.e., its asymptotic performance is even worse than that of Dijkstra's algorithm. On the other hand, the solutions that offer favorable query time complexities often entail prohibitive preprocessing cost or space overhead, rendering them only applicable for small datasets. For example, the state-of-the-art approaches by Samet et al.\ \cite{ssa08} and Abraham et al.\ \cite{afg10} provide superior bounds on query time, but they require pre-computing the shortest path between {\em any} pair of nodes, which is impractical for the large road networks commonly used in modern map applications.
%
%The aforementioned time complexities of AH rely on an assumption on road networks (to be clarified in Section~\ref{sec:def}). We provide detailed discussion on the assumption, and we demonstrate its applicability on practical road networks with extensive experiments on a large collection of real datasets. These experimental findings not only form a basis for our theoretical claims but also shed light on the characteristics of real road networks, which paves the path for future research on shortest path and distance queries.
%
%
%
%\begin{figure}[t]
%%\begin{small}
%\begin{minipage}[t]{1.63 in}
%\centering
%\includegraphics[width = 1.62in]{./figures/original.eps}
%\vspace{-3.52mm}
%\figcapup \caption{Road network $G$ $\quad$ (with bidirectional edges).} \figcapdown \label{fig:intro-original}
%\end{minipage}
%\begin{minipage}[t]{1.85 in}
%\centering
%%\hspace{3mm}
%\includegraphics[width = 1.70in]{./figures/levelhierarchy.eps}
%\figcapup \caption{Node hierarchy $\G$. $\quad$} \figcapdown \label{fig:intro-hierarchy}
%\end{minipage}
%%\end{small}
%\vspace{-1mm}
%\end{figure}





%% [inlined 2-Definition/definition.tex]


\section{Problem Definition}\label{sec:def}

%and let $n = |V|$ and $m = |E|$.

Let $G$ be a weighted and directed graph with a set $V$ of nodes and a set $E$ of edges. Each edge $e$ in $E$ is associated with a positive weight $l(e)$, which is referred to as the {\em length} of $e$. A path $P$ in $G$ is a sequence of nodes $\langle v_1, v_2, \ldots, v_k \rangle$, such that $\langle v_i, v_{i+1}\rangle$ ($i \in [1, k-1]$) is a directed edge in $G$. The length of $P$ is defined as the sum of the length of each edge on $P$. We use $l(e)$ and $l(P)$ to denote the length of an edge $e$ and a path $P$, respectively.

%\parskip = \PAR

\vgap

%\begin{table}[t]
%\renewcommand{\arraystretch}{1.2}
%\centering
%\begin{small}
%%\begin{tabular}{|@{\hspace{1mm}}c@{\hspace{1mm}}|l|}
%\begin{tabular}{|c|l|}
%  \hline
%  {\bf Symbol} & \makebox[6cm][c]{\bf Description}  \\ \hline
%  $G = (V, E)$ & a graph with a node set $V$ and an edge set $E$ \\ \hline
%  $e = \langle v_i, v_j\rangle$ & an edge from a node $v_i$ to a node $v_j$ \\ \hline
%  $l(e)$ & the length of an edge $e$ \\ \hline
%%  $P = \langle v_1, v_2, \ldots, v_k\rangle$ & a path  \\ \hline
%  $l(P)$ & the length of a path $P$ \\ \hline
%  $\d(v_i, v_j)$ & the distance from $v_i$ to $v_j$ in $G$ \\ \hline
%%  $n$ & the number of nodes in $V$ \\ \hline
%%  $m$ & the number of edges in $E$ \\ \hline
%%  $M$ & the size of the main memory available \\ \hline
%  $B$ & the size of a disk block \\ \hline
%  $\G$ & the augmented graph maintained by HoD \\ \hline
%  $\GC$ & the core (sub)graph of $\G$ (see Section~\ref{sec:overview-pre}) \\ \hline
%  $r(v)$ & the rank of a node $v$ (see Section~\ref{sec:overview-pre}) \\ \hline
%  $s(v)$ & the score of a node $v$ (see Section~\ref{sec:pre-select}) \\ \hline
%  $\GF$ & the forward graph (see Section~\ref{sec:pre-terminate}) \\ \hline
%  $\GB$ & the backward graph (see Section~\ref{sec:pre-terminate}) \\ \hline
%  $\GR$ & the reversed version of $\GB$ (see Section~\ref{sec:pre-terminate}) \\ \hline
%  $\theta(v)$ & the index value of a node $v$ (see Section~\ref{sec:p2p}) \\ \hline
%\end{tabular}
%\end{small}
%\tblcapup \caption{Frequently Used Notations.} \label{tbl:def-notation}
%\tblcapdown
%\end{table}

For any two nodes $s$ and $t$ in $G$, we define the {\em distance} from $s$ to $t$, denoted as $\d(s, t)$, as the length of the shortest path from $s$ to $t$. Given a {\em source node} $s$ in $G$, a {\em single-source distance (SSD)} query asks for the distance from $s$ to any other node in $G$. Meanwhile, a {\em single-source shortest path (SSSP)} query from $s$ retrieves not only the distance from $s$ to any other node $v$, but also the {\em predecessor} of $v$, i.e., the node that immediately precedes $v$ in the shortest path from $s$ to $v$. Note that, given the predecessor of each node, we can easily reconstruct the shortest path from $s$ to any node $v$ by backtracking from $v$ following the predecessors. One may also consider an alternative formulation of SSD (resp.\ SSSP) query that, given only a {\em destination node} $t$, asks for the distance (resp.\ shortest path) from any other node to $t$. For simplicity, we will focus on SSD and SSSP queries from a source node $s$, but our solution can be easily extended to handle queries under the alternative formulation.

%\begin{figure*}[t]
%%\centering
%\begin{small}
%\centering
%\begin{tabular}{ccccc}
%\multicolumn{5}{c}{\hspace{-4mm} \includegraphics[height=3mm]{./figures/legend1.eps}} \vspace{1mm} \\
%\hspace{0mm} \includegraphics[width=32mm]{./figures/original.eps}
%&
%\hspace{2mm} \includegraphics[width=24mm]{./figures/pre1.eps}
%&
%\hspace{2mm} \includegraphics[width=24mm]{./figures/pre2.eps}
%&
%\hspace{2mm} \includegraphics[width=24mm]{./figures/pre3.eps}
%&
%\hspace{2mm} \includegraphics[width=24mm]{./figures/pre4.eps} \\
%%\hspace{0mm} (a) Graph $G$. & \hspace{0mm} (b) $v_1, v_2, v_3$ removed. & \hspace{0mm} (c) $v_4$ removed. &
%\hspace{0mm} (a) Graph $G$. & \hspace{2mm} (b) Reduced graph $G_1$. & \hspace{2mm} (c) Removal of $v_4$. &
%\hspace{2mm} (d) Reduced graph $G_2$. & \hspace{2mm} (e) Reduced graph $G_3$.
%\end{tabular}
%\figcapup \caption{Graph reduction and shortcut construction.} \figcapdown
%\label{fig:over-steps}
%\end{small}
%\end{figure*}


\vgap

Let $M$ be the size of the main memory available, and $B$ be the size of a disk block, both measured in the number of words. We assume that $B \le |V| \le M \le |E|$, i.e., the main memory can accommodate all nodes but not all edges in $G$. This is a realistic assumption since modern machines (even the commodity ones) have gigabytes of main memory, which is sufficient to store the node set of a graph with up to a few billion nodes. On the other hand, the number of edges in a real graph is often over an order of magnitude larger than the number of nodes, due to which $E$ can be enormous in size and does not fit in the main memory.

\vgap




Our objective is to devise an index structure on $G$ that answers any SSD or SSSP query with small I/O and CPU costs, such that the index requires at most $M$ main memory in pre-computation and query processing. In what follows, we will first focus on SSD queries in Sections \ref{sec:overview}-\ref{sec:sssp}, and will extend our solution for SSSP queries in Section~\ref{sec:exten}.


%For ease of reference, Table~\ref{tbl:def-notation} lists the notations that are frequently used in this paper.




%Let $G$ be a road network. We assume that $G$ is a directed, degree-bounded, and connected graph with a node set $V$ and an edge set $E$, such that (i) $|V| = n$, (ii) each node in $V$ locates in a two-dimensional space, and (iii) each edge $e \in E$ is associated with a positive weight $l(e)$. Without loss of generality, we consider that $l(e)$ equals the distance between $e$'s two endpoints.
%%, and we assign to each edge in $E$ a unique ID.
%For any path $P$ in $G$, we define its length $l(P)$ as the total length of the edges in $P$.
%
%%For simplicity, in all of our figures, we use an undirected line between two nodes to represent a bidirectional edge connecting the nodes.
%
%We study two types of queries on $G$, namely, shortest path queries and distance queries. Given an ordered pair of nodes $(s, t) \in V \times V$, a shortest path query asks for a sequence of edges $e_1, e_2, \ldots, e_k$ that form a path from $s$ to $t$, such that $\sum_{i = 1}^k l(e_i)$ is minimized. On the other hand, a distance query from $s$ to $t$ asks only for the value of $\sum_{i = 1}^k l(e_i)$ instead of the actual shortest path. For convenience, we define the {\em distance} from $s$ to $t$ as $dist(s, t) = \sum_{i = 1}^k l(e_i)$.
%
%Our solution for shortest path and distance queries is developed based on an observation on the properties of real road networks, as explained in the following.
%
%\header {\bf Observation.} Assume that we impose a square grid $R$ on $G$. Let $B$ be a region containing $4 \times 4$ cells in the grid. We define the left-most (resp.\ right-most) column of cells in $B$ as the {\em west strip} (resp.\ {\em east strip}) of $B$, and we refer to the vertical line that evenly divides $B$ as the {\em vertical bisector} of $B$. We also define $B$'s {\em north strip}, {\em south strip}, and {\em horizontal bisector} in a similar manner. For example, Figure~\ref{fig:def-grid} illustrates (i) a square grid imposed on the road network in Figure~\ref{fig:intro-original}, (ii) a region $B$ covering $4 \times 4$ grid cells, and (iii) the strips and bisectors of $B$.
%
%
%
%%For each grid cell $C$ in $R$, we define the {\em inner shell} ({\em outer shell}) of $C$ as the boundary of the region comprising the $3 {\times} 3$ ($5 {\times} 5$) cells centered at $C$. For example, Figure~\ref{fig:def-grid} illustrates a grid imposed on the road network in Figure~\ref{fig:intro-original}. For the cell $C_0$ that contains $v_9$, its inner and outer shells are highlighted with dashed-lines and dotted-lines, respectively. Intuitively, the inner shell (outer shell) of a cell encloses a region that is within small (medium) distance to the cell. We define the {\em vicinity graph} of $C$ as the subgraph of $G$ comprising any edge that has at least one endpoint inside of $C$'s outer shell. For instance, in Figure~\ref{fig:def-grid}, the vicinity graph of $C$ consists of all edges in the road network except $\langle v_4, v_7 \rangle$, $\langle v_7, v_8 \rangle$, and $\langle v_8, v_3 \rangle$, as the endpoints of those three edges all fall in the outside of $C$'s outer shell.
%
%
%
%%Now let us consider the set $S_P$ of shortest paths that connect any node in $C$ with any node in the exterior of $C$'s outer shell, i.e., the set of paths that should be followed when traveling between $C$ and distant locations. We observe that, in practice, {\em there should exist a small set $S_e$ of edges in the vicinity graph of $C$ that covers all paths in $S_P$} (i.e., any path in $S_p$ should pass through at least one edge in $S_e$). To explain, assume that $C$ covers the area of a city. In that case, any node $v$ beyond the outer shell of $C$ should correspond to a location in a distant city, and hence, any shortest path between $v$ and a location in $C$ should go through some inter-city highway. This indicates that, if we select an edge from every inter-city highway adjacent to $C$, then the resulting set of edges should cover any shortest path between $C$ and any location out of $C$'s outer shell. Furthermore, the cardinality of the edge set should be bounded by a constant, as there should only exist a handful of inter-city highways in the vicinity of a city. Similar observations can be made when $C$ corresponds to a larger region (e.g., a state) or a smaller one (e.g., a neighborhood). Intuitively, if we can pre-compute a set of edges that cover any shortest path between $C$ and any distant location, then we may efficiently answer the shortest path and distance queries that concern $C$. In the following, we will show that such a set of edges for a cell $C$ can be computed by identifying the {\em arterial edges} of $C$.
%
%We observe that, in practice, the shortest paths between the west and east strips of $B$ can often be {\em covered} by a small set $S_{we}$ of road network edges located on $B$'s vertical bisector. That is, given any two points in $B$'s west and east strips, respectively, the shortest path between the two points should pass through at least one edge in $S_{we}$. For instance, suppose that $B$ covers the area of a state. In that case, any shortest path $P$ between the west and east strips of $B$ corresponds to a route that connects the west and east ends of the state. Intuitively, $P$ would have to pass through some major intra-state highways. Therefore, if $S_{we}$ contains the road network edges on the intra-state highways that intersect $B$'s vertical bisector, then $S_{we}$ should cover any aforementioned shortest path $P$. Furthermore, the cardinality of $S_{we}$ should be small, as there should exist only a handful of major highways in the state that go across the vertical bisector. Similar statements can be made even when $B$ corresponds to a larger region (e.g., a continent) or a smaller one (e.g., a city). In addition, we also observe that all shortest paths between the north and south strips of $B$ can be covered by a few edges intersecting $B$'s horizontal bisector.
%
%
%
%The above observations are similar in spirit to those made in previous work \cite{afg10,bfm06,ssa09}, which all illustrate that there exists a small set of {\em important} road network edges or nodes that cover all shortest paths connecting distant regions (see Section~\ref{sec:related} for a survey of related work). In what follows, we will formalize our observations and provide empirical evidence, so as to form a basis for further discussions in Sections \ref{sec:fc} and \ref{sec:ah}.
%
%\header {\bf Formalization.} Give a region $B$ of $4\times4$ grid cells, we say that a road network path $P$ is a {\em local path} in $B$, if at most one edge in $P$ intersects the boundary of $B$. For instance, in Figure~\ref{fig:def-grid}, the paths $\langle v_9, v_5, v_8 \rangle$ and $\langle v_{11}, v_7, v_4 \rangle$ are both local paths in $B$. A local path in $B$ is the {\em shortest}, if it is shorter than any other local path in $B$ with the same endpoints. For simplicity, we assume that no two local paths in $B$ have the same length -- This assumption can be enforced by adding a small perturbation to each edge in $G$. (Interested readers are referred to our online and anonymous technical report \cite{AhReport} for a thorough treatment of this issue.)
%
%We are interested in the local shortest paths between opposite strips of $B$, and a set of edges on $B$'s bisectors that cover all such paths, as defined in the following.
%
%\begin{figure}[t]
%%\begin{small}
%\centering
%\includegraphics[width = 3.3in]{figures/grid.eps}
%\vspace{-3.5mm}
%\figcapup \caption{Strips and bisectors of a region with $4 \times 4$ cells.} \label{fig:def-grid} \figcapdown \vspace{-1.5mm}
%\end{figure}
%
%\begin{definition}[Spanning Paths \& Arterial Edges] \label{def:def-ae}
%A local shortest path $P$ in $B$ is a {\bf spanning path} of $B$, if (i) the two endpoints of $P$ are on different sides of a bisector of $B$ (denoted as $l_b$), and (ii) neither of the endpoints is contained in a grid cell adjacent to the bisector $l_b$. Any edge on $P$ that intersect $l_b$ is an {\bf arterial edge} of $B$.
%\end{definition}
%
%By Definition~\ref{def:def-ae}, the path $P = \langle v_9, v_6, v_{10}, v_8\rangle$ in Figure~\ref{fig:def-grid} is a spanning path of $B$, since (i) $P$ is a local shortest path of $B$, (ii) $v_9$ and $v_8$ are on different sides of $B$'s vertical bisector, and (iii) neither $v_9$ nor $v_8$ is in a grid cell adjacent to the bisector. Accordingly, the edge $\langle v_6, v_{10}\rangle$ is an arterial edge of $B$, as it is the only edge in $P$ that intersects $B$'s vertical bisector. Likewise, $\langle v_{11}, v_7, v_4\rangle$ is also a spanning path of $B$, and $\langle v_{11}, v_7\rangle$ is an arterial edge of $B$.
%
%As explained previously, the number of arterial edges of $B$ is small in practice since there usually exist only a few major connections between opposite strips of $B$. We formalize this observation as follows.
%
%%Now let us consider the set $S_P$ of shortest paths that connect any node in $C$ with any node in the exterior of $C$'s outer shell, i.e., the set of paths that should be followed when traveling between $C$ and distant locations. We observe that, in practice, {\em there should exist a small set $S_e$ of edges in the vicinity graph of $C$ that covers all paths in $S_P$} (i.e., any path in $S_p$ should pass through at least one edge in $S_e$). To explain, assume that $C$ covers the area of a city. In that case, any node $v$ beyond the outer shell of $C$ should correspond to a location in a distant city, and hence, any shortest path between $v$ and a location in $C$ should go through some inter-city highway. This indicates that, if we select an edge from every inter-city highway adjacent to $C$, then the resulting set of edges should cover any shortest path between $C$ and any location out of $C$'s outer shell. Furthermore, the cardinality of the edge set should be bounded by a constant, as there should only exist a handful of inter-city highways in the vicinity of a city. Similar observations can be made when $C$ corresponds to a larger region (e.g., a state) or a smaller one (e.g., a neighborhood). Intuitively, if we can pre-compute a set of edges that cover any shortest path between $C$ and any distant location, then we may efficiently answer the shortest path and distance queries that concern $C$. In the following, we will show that such a set of edges for a cell $C$ can be computed by identifying the {\em arterial edges} of $C$.
%
%%\begin{definition}[Arterial Edge] \label{def:def-ae}
%%For any cell $C$ of a square grid on $G$, an edge $e \in E$ is an {\bf arterial edge} for $C$ iff it satisfies all of the following conditions:
%%\begin{enumerate} %[topsep = 0.2 \baselineskip, itemsep = 0pt, parsep = 0.2 \baselineskip, partopsep = 0 \baselineskip]
%%\item $e$ goes across $C$'s inner shell, i.e., the two endpoints of $e$ are inside and outside of $C$'s inner shell, respectively.
%%
%%\item $e$ is on a shortest path $P$ in the vicinity graph of $C$ (i.e., all edges in $P$ are in the vicinity graph), such that the two endpoints of $P$ are in $C$ and out of $C$'s outer shell, respectively.
%%
%%\item Among the edges in $P$ that go across $C$'s inner shell, $e$ has the smallest ID. \done
%%\end{enumerate}
%%\end{definition}
%
%%By Definition~\ref{def:def-ae}, for any shortest path $P$ (in the vicinity graph of $C$) that connects $C$ with the exterior of $C$'s outer shell, the edge in $P$ that goes through $C$'s inner shell is an arterial edge of $C$ (conditions 1 and 2). In case that multiple edges in $P$ intersect $C$'s inner shell, we will select the edge with the smallest ID as the arterial edge (condition 3). For example, if we assume that the thin (thick) lines in Figure~\ref{fig:def-grid} represent edges with weight $1$ ($2$), then $\langle v_6, v_{10} \rangle$ is an arterial edge of cell $C_0$. This is because (i) $\langle v_6, v_{10} \rangle$ goes across $C_0$'s inner shell, (ii) in the vicinity graph of $C_0$, the shortest path from $v_9$ to $v_8$ goes through $\langle v_6, v_{10} \rangle$, and (iii) $v_9$ and $v_8$ are inside of $C_0$ and outside of $C_0$'s outer shell, respectively.
%%
%%Notice that the arterial edges of a cell $C$ are defined based on the shortest paths in the vicinity graph of $C$ (instead of the shortest paths in the road network $G$). Nevertheless, it can proved that the arterial edges of $C$ still cover all shortest paths in $G$ that connect $C$ with the exterior of $C$'s outer shell.
%%
%%\begin{proposition} \label{prop:def-ae}
%%Any shortest path in $G$ that connects a cell $C$ with the exterior of its outer shell must pass through an arterial edge of $C$.
%%\end{proposition}
%%
%%The proofs for the propositions, lemmas, theorems in this paper are included in Appendix~\ref{sec:proofs}. The following assumption formalizes our observation that all shortest paths between a cell $C$ and the locations far from $C$ can be covered with a small set of edges.
%
%
%\begin{assumption}[Arterial Dimension] \label{assu:def-ae}
%For any square grid on $G$ and any region $B$ with $4 \times 4$ grid cells, the number of arterial edges of $B$ is at most a constant $\lambda$, referred to as the {\bf arterial dimension} of $G$.
%\end{assumption}
%
%To demonstrate the applicability of Assumption~\ref{assu:def-ae}, we have conducted an experiment on eight real datasets that represent various parts of the road network in the United States (see Section~\ref{sec:exp} for details of the datasets). On each dataset, we impose a $2^r \times 2^r$ square grid ($r \in [3, 17]$), and we compute the number of arterial edges for each region $B$ of $4 \times 4$ cells (ignoring the regions that do not contain any road network node). After that, we compute the maximum number of arterial edges for the regions, as well as the mean, $90\%$ quantile, and $99\%$ quantile of the numbers. Figure~\ref{fig:def-ae-num} plots the results as functions of the grid resolution $r$. Regardless of the grid resolution and the size of the dataset, the maximum number of arterial edges for a ($4 {\times} 4$)-cell region is at most $97$, and it is below $60$ in most cases. Furthermore, the $90\%$ quantile and $99\%$ quantile are at most $60$, while the mean is never above $22$. This indicates that practical road networks have fairly small arterial dimensions. In Sections \ref{sec:fc} and \ref{sec:ah}, we will exploit this property of road networks to construct efficient indices for shortest path and distance queries.
%
%
%
%\begin{figure*}[!t]
%\begin{small}
%\begin{tabular}{cccc}
%\multicolumn{4}{c}{\hspace{-4mm} \includegraphics[height=3.6mm]{./figures/tn-legends.eps}} \\ \vspace{-2.5mm} \\
%\hspace{-9.5mm}\includegraphics[height=38mm]{./figures/2ME.eps}
%&
%\hspace{-12.5mm}\includegraphics[height=38mm]{./figures/3COL.eps}
%&
%\hspace{-12.5mm}\includegraphics[height=38mm]{./figures/4FLA.eps}
%&
%\hspace{-12.5mm}\includegraphics[height=38mm]{./figures/5CAL.eps} \\
%
%\hspace{-5.5mm}(a) ME ($n = 187{,}315$) & \hspace{-8mm}(b) CO ($n = 435{,}666$) & \hspace{-8mm}(c) FL ($n = 1{,}070{,}376$) & \hspace{-8mm}(d) CA ($n = 1{,}890{,}815$) \\ \vspace{-1.5mm}\\
%
%\hspace{-9.5mm}\includegraphics[height=38mm]{./figures/6E.eps}
%&
%\hspace{-12.5mm}\includegraphics[height=38mm]{./figures/7W.eps}
%&
%\hspace{-12.5mm}\includegraphics[height=38mm]{./figures/8CTR.eps}
%&
%\hspace{-12.5mm}\includegraphics[height=38mm]{./figures/9USA.eps} \\
%
%\hspace{-5.5mm}(e) E-US ($n = 3{,}598{,}623$) & \hspace{-8mm}(f) W-US ($n = 6{,}262{,}104$) & \hspace{-8mm}(g) C-US ($n = 14{,}081{,}816$) & \hspace{-8mm}(h) US ($n = 23{,}947{,}347$)
%\end{tabular}
%\end{small}
%\figcapup \caption{Arterial dimensions of real road networks.} \figcapdown
%\label{fig:def-ae-num}
%\end{figure*}
%
%%To demonstrate the applicability of Assumption~\ref{assu:def-ae}, we have conducted an experiment on ten real datasets that represent various parts of the road network in the United States (see Section~\ref{sec:exp} for details of the datasets). On each dataset, we impose a $2^r \times 2^r$ square grid ($r \in [3, 17]$), and we compute the number of arterial edges for each region $B$ of $4 \times 4$ cells (ignoring the regions that do not contain any road network node). After that, we compute the maximum number of arterial edges for the non-empty cells, as well as the mean, $90\%$ quantile, and $99\%$ quantile of the numbers. Figure~\ref{fig:def-ae-num} plots the results (on the eight largest datasets) as functions of the grid resolution $r$. Regardless of the grid resolution and the size of the dataset, the maximum number of arterial edges for a cell is at most $83$, and it is well below $50$ in most cases. Furthermore, the $90\%$ quantile and $99\%$ quantile are always below $40$, while the mean is never above $20$. This indicates that practical road networks have fairly small arterial dimensions. In Sections \ref{sec:ah} and \ref{sec:exten}, we will show how this property of road networks can be exploited to accelerate shortest path and distance queries.
%
%%In addition to Assumption~\ref{assu:def-ae}, we also formulate a hypothesis about the {\em stabbing number} of the road network $G$, which is defined as the maximum number of edges in $G$ that simultaneously intersect a horizontal (or vertical) line.
%%
%%\begin{assumption}[Stabbing Number] \label{assu:def-stab}
%%The stabbing number of $G$ is $O(\sqrt{n})$, where $n$ is the number of nodes in $G$.
%%\end{assumption}
%%
%%Intuitively, Assumption~\ref{assu:def-stab} states that the edges in $G$ does not follow too skewed a distribution along the horizontal or vertical dimension. To test this assumption, we compute the stabbing numbers of our datasets, and we plot them in Figure~\ref{fig:def-stabbing} as a function of the number $n$ of nodes in the road network. The dotted line in Figure~\ref{fig:def-stabbing} is a linear regression model constructed from the observed stabbing numbers under the log-log scale. Observe that the stabbing number grows at a rate smaller than $\sqrt{n}$, which is consistent with Assumption~\ref{assu:def-stab}.
%%
%%\begin{figure}[t]
%%%\begin{small}
%%\centering
%%\includegraphics[height = 40mm]{figures/tn-varygraph-cut.eps}
%%\figcapup \caption{Stabbing Number vs.\ Number of Nodes $n$} \label{fig:def-stabbing} \figcapdown
%%\end{figure}
%
%%Finally, we make an assumption about the $L_\infty$ distance between any two nodes in $G$.
%%
%%\begin{assumption}[$L_\infty$ Distance] \label{assu:def-log}
%%Let $d_{max}$ and $d_{min}$ be the maximum and minimum $L_\infty$ distances between any two nodes in $G$. Then, $\log_2(d_{max} / d_{min})$ is at most a constant.
%%\end{assumption}
%%
%%To understand why Assumption~\ref{assu:def-log} holds, observe that, in practice, the $L_\infty$ distance between any two nodes is lower-bounded by a small constant, say, $1$ foot. On the other hand, $d_{max}$ should be upper-bounded by a large constant, e.g., the length of the Equator ($\approx 1.3 \times 10^8$ feet). In that case, $\log_2(d_{max} / d_{min}) < 27$.
%
%%\header {\bf Remark.} As will be clarified in Section~\ref{sec:ah}, Assumption~\ref{assu:def-stab} affects only the preprocessing time complexity of our solution. In particular, given Assumptions \ref{assu:def-ae} and \ref{assu:def-log} but not Assumption~\ref{assu:def-stab}, our approach requires $O(n^2)$ preprocessing time, but its space and time complexities remain the same.
%
%
%%Throughout this paper, a road network is a degree-bounded directed graph $G=(V,E)$, where $V$ is a set of nodes (e.g. road intersections) and $E$ a set of edges (e.g. roads, highways etc.). %Due to the nature of road network, we only consider the undirected graph, i.e. if $(u,v)\in E$, then $(v,u)\in E$.
%%Let each edge $e\in E$ be associated with a positive real-valued weight $weight(e)$. $weight(e)$ could be measured by the length of the road or the travelling time spent on that road. Given two nodes $s,t\in V$, a shortest path from $s$ to $t$, denoted as $SP(s,t)$, is a consecutive sequence of edges $(e_1, e_2, \ldots, e_k)$ connecting $s$ to $t$, such that the length of the path $\sum_{i=1}^k weight(e_i)$ is minimal. In a clear context we also use a sequence of nodes to represent a shortest path, i.e. $SP(s,t)=(v_0,v_1\ldots v_k)$ such that $v_0=s$, $v_k=t$ and $<v_{i-1},v_i>=e_i$ for $1\le i\le k$. For convenience, we use $dist(s,t)$ to denote the shortest-path distance from $s$ to $t$. A shortest-path query returns the edges sequence $SP(s,t)$, while a distance query only returns the value of $dist(s,t)$.
%
%
%
%
%
%%\begin{figure}[t]
%%%\begin{small}
%%\hspace{-4mm}
%%\begin{minipage}[t]{1.67 in}
%%\centering
%%\includegraphics[width = 1.8in]{figures/grid2.eps}
%%\vspace{-3.49mm}
%%\figcapup \hspace{-3mm} \caption{Transit Nodes} \label{fig:def-grid} \figcapdown
%%\end{minipage}
%%\hspace{2mm}
%%\begin{minipage}[t]{1.73 in}
%%\centering
%%\hspace{2.5mm}
%%\includegraphics[width = 1.6in]{./figures/qu-de-sd2.eps}
%%\vspace{-3.49mm}
%%\figcapup \caption{Stabbing Number vs.\ $n$} \figcapdown \label{fig:def-stabbing}
%%\end{minipage}
%%%\end{small}
%%\end{figure}



%% [inlined 2.5-Overview/overview.tex]


\section{Solution Overview}\label{sec:overview}

As mentioned in Section~\ref{sec:intro}, the main structure of HoD is a graph $\G$ that augments the input graph $G$ with shortcuts. In this section, we present the overall idea of how the shortcuts in $\G$ are constructed and how they can be utilized for query processing, so as to form a basis for the detailed discussions in subsequent sections.


\begin{figure*}[t]
%\centering
\begin{small}
\centering
\begin{tabular}{ccccc}
\multicolumn{5}{c}{\hspace{-4mm} \includegraphics[height=3mm]{./figures/legend1.eps}} \vspace{0mm} \\
\hspace{-2mm} \includegraphics[width=32mm]{./figures/original.eps}
&
\hspace{0mm} \includegraphics[width=24mm]{./figures/pre1.eps}
&
\hspace{0mm} \includegraphics[width=24mm]{./figures/pre3.eps}
&
\hspace{-2mm} \includegraphics[width=24mm]{./figures/pre4.eps}
&
\hspace{-2mm} \includegraphics[width = 47mm]{./figures/ranks2.eps}\\
\hspace{-2mm} (a) Graph $G$. & \hspace{0mm} (b) Reduced graph $G_1$. &
\hspace{0mm} (c) Reduced graph $G_2$. & \hspace{-2mm} (d) Reduced graph $G_3$ & \hspace{-5mm} (e) Augmented Graph $\G$. \\
(equivalently, $G_0$) & & & (i.e., the core graph $\GC$). &
\end{tabular}
\vspace{-4mm} \caption{Graph reduction and shortcut construction.} \figcapdown
\label{fig:over-steps}
\end{small}
\end{figure*}

%\parskip = \PAR

\subsection{Shortcut Construction} \label{sec:overview-pre}

In a nut shell, HoD constructs shortcuts with an iterative procedure, which takes as input a copy of the graph $G$ (denoted as $G_0$). In the $i$-th ($i \ge 1$) iteration of the procedure, HoD first {\em reduces} $G_{i-1}$ by removing a selected set of {\em less important} nodes in $G_{i-1}$, and then, it constructs shortcuts in the reduced graph to ensure that the distance between any two remaining nodes is not affected by the node removal. The resulting graph (with shortcuts added) is denoted as $G_i$, and it is fed as the input of the ($i{+}1$)-th iteration of procedure. This procedure terminates only when the reduced graph $G_i$ is {\em sufficiently small}. All shortcuts created during the procedure are inserted into the original graph $G$, leading to an augmented graph $\G$ that would be used by HoD for query processing. We illustrate the iterative procedure with an example as follows.

%\parskip = \PAR

% such that each iteration (i) {\em reduces} $G$ by removing a selected set of {\em less important} nodes, and (ii) inserts shortcuts into the reduced graph to ensures that the distance between any two remaining nodes is not affected by the removal of nodes. This procedure terminates when the reduced graph is {\em sufficiently small}. All shortcuts created during the procedure are added into the original graph $G$ to facilitate query processing. We illustrates the procedure with an example.

\vgap

\begin{example} \label{example:over-steps}
\em Assume that the input to the iterative procedure is the graph $G_0$ in Figure~\ref{fig:over-steps}a. Further assume that the reduced graph is sufficiently small if it contains at most two nodes and two edges. In the first iteration of the procedure, HoD inspects $G_0$ and identifies $v_1$, $v_2$, and $v_3$ as less important nodes. To explain, observe that the node $v_1$ in $G_0$ does not have any incoming edge, while $v_2$ and $v_3$ have no outgoing edges. As a consequence, $v_1$, $v_2$, and $v_3$ do not lie on the shortest path between any two other nodes. That is, even if we remove $v_1$, $v_2$, and $v_3$ from $G_0$, the distance between any two remaining nodes is not affected. Intuitively, this indicates that $v_1$, $v_2$, $v_3$ are of little importance for SSD queries. Therefore, HoD eliminates $v_1$, $v_2$, and $v_3$ from $G_0$, which results in the reduced graph $G_1$ in Figure~\ref{fig:over-steps}b.

\vgap

In the second iteration, HoD selects $v_4$, $v_5$, and $v_6$ as the less important nodes in $G_1$, and removes them from $G_1$. The removal of $v_4$ changes the distance from $v_8$ to $v_9$ to $+\infty$, since $\langle v_8, v_4, v_9\rangle$ is the only path (in $G_1$) that starts at $v_8$ and ends at $v_9$. To mitigate this change, HoD inserts into $G_1$ a shortcut $\langle v_8, v_9\rangle$ that has the same length with $\langle v_8, v_4, v_9\rangle$, as illustrated in Figure~\ref{fig:over-steps}c. As such, the distance between any two nodes in $G_1$ remains unchanged after $v_4$ is removed. Similarly, when HoD eliminates $v_6$, it constructs a shortcut $\langle v_9, v_7\rangle$ with a length $2$ to reconnect the two neighbors of $v_6$. Meanwhile, $v_5$ is removed without creating any shortcut, since deleting $v_5$ does not change the distance between its two neighbors. Figure~\ref{fig:over-steps}c illustrates the resulting reduced graph $G_2$.

To explain why HoD chooses to remove $v_4$, $v_5$, and $v_6$ from $G_1$, observe that each of those nodes has only two neighbors. For any of such nodes, even if the removal of the node changes the distance between its neighbors, HoD only needs to construct one shortcut to reconnect its neighbors. In other words, the number of shortcuts required is minimum, which helps reduce the space consumption of HoD. In contrast, if HoD chooses to remove $v_9$ from $G_1$ (which has a larger number of neighbors than $v_4$, $v_5$, and $v_6$), then much more shortcuts would need to be constructed.

%One may question why $v_7$ or $v_8$ is not removed by HoD, given that each of them also has only two neighbors. The answer is that $v_7$ and $v_8$ cannot be removed with $v_4$, $v_5$, and $v_6$ in the same iteration, as will be clarified in Section~\r
%
%On the other hand, although each
%
% since the elimination of those nodes has a relatively small effect on the other nodes in $G_1$. In particular, each of $v_4$, $v_5$, and $v_6$ has only two neighbors, and hence, even if removing the node changes the distance between its neighbors, we can easily mitigate the effect by creating one shortcut that reconnects the two neighbors of the node\footnote{Each of $v_7$ and $v_8$ also has two neighbors only, but they cannot be selected along with $v_4$, $v_5$, and $v_6$, as will be explained in Section~\ref{sec:pre}.}. For instance, if we remove $v_4$ from $G_1$, the distance from $v_8$ to $v_9$ would change from $2$ to $+\infty$, since $\langle v_8, v_4, v_9\rangle$ is the only path (in $G_1$) that starts at $v_8$ and ends at $v_9$. However, this change can be mitigated by adding into $G_1$ a shortcut $\langle v_8, v_9\rangle$ with a length $2$, as illustrated in Figure~\ref{fig:over-steps}c. As such, the distance between any two nodes in $G_1$ remains unchanged after the removal of $v_4$. Similarly, when HoD eliminates $v_6$, it constructs a shortcut $\langle v_9, v_7\rangle$ with length $2$ to reconnect the two neighbors of $v_6$. Meanwhile, $v_5$ is removed without creating any shortcut, since deleting $v_5$ does not change the distance between its two neighbors. Figure~\ref{fig:over-steps}d illustrates the resulting reduced graph $G_2$.

%, as they can be eliminated without significant effects on the remaining nodes

\vgap

Finally, in the third iteration, HoD removes $v_7$ and $v_8$ from $G_2$ as they are considered unimportant. The removal of $v_7$ leads to a new shortcut $\langle v_9, v_{10}\rangle$ with a length $3$, since $\langle v_9, v_7, v_{10}\rangle$ is the only path connecting $v_9$ to $v_{10}$, and the length of the path equals $3$. On the other hand, $v_8$ is directly eliminated as it is not on the shortest path between its only two neighbors $v_9$ and $v_{10}$. Figure~\ref{fig:over-steps}d shows the reduced graph $G_3$ after the removal of $v_7$ and $v_8$.

Assume that the reduced graph $G_3$ is considered sufficiently small by HoD. Then, the iterative procedure of HoD would terminate. The three shortcuts created during the procedure (i.e., $\langle v_8, v_9\rangle$, $\langle v_9, v_7 \rangle$, and $\langle v_9, v_{10}\rangle$) are added into the original graph $G$, which leads to the augmented graph $\G$ in Figure~\ref{fig:over-steps}d. \done
\end{example}

\vgap

The above discussion leaves several issues open, i.e., (i) the specific criterion for identifying less important nodes in the reduced graph, (ii) the detailed algorithm for generating shortcuts after node removal, and (iii) the exact termination condition of the reduction procedure. We will clarify these issues in Section~\ref{sec:pre} by presenting the detailed preprocessing algorithm of HoD. For the discussions in the rest of this section, it suffices to know that when HoD terminates the reduction procedure, the reduced graph must fit in the main memory. We use $\GC$ to denote this memory-resident reduced graph, and we refer to it as the {\em core graph}. (Note that $\GC$ is a subgraph of the augmented graph $\G$.) In addition, we define the {\em rank} $r(v)$ of each node $v$ in $G$ as follows:
\begin{enumerate} [itemsep = 0.5mm, topsep = 2mm] %, partopsep = 0mm]
\item If $v$ is removed in the $i$-th iteration of the iterative procedure, then $r(v) = i$;

\item If $v$ is not removed in any iteration (i.e., $v$ is retained in the core graph $\GC$), then $r(v) = 1 + \max_{v \notin \GC} r(v)$, i.e., $r(v)$ is larger than the maximum rank of any node not in $\GC$.
\end{enumerate}
For instance, in Example~\ref{example:over-steps},
%the core graph $\GC$ is as illustrated in Figure~\ref{fig:over-steps}e, and
the ranks of $v_1$, $v_2$, and $v_3$ equal $1$, since they are removed from $G$ in the first iteration of the reduction procedure. Similarly, $r(v_4) = r(v_5) = r(v_6) = 2$, and $r(v_7) = r(v_8) = 3$. The ranks of $v_9$ and $v_{10}$ equal $4$, since they are in the core graph $\GC$. The ranks of the nodes are utilized in the query processing algorithms of HoD, as will be illustrated shortly. Unless otherwise specified, we use the term {\em edge} to refer to both a shortcut and an original edge in $\G$.

\subsection{Query Processing} \label{sec:overview-query}

Given an SSD query from a node $s$, HoD answers the query with two traversals of the augmented graph $\G$. The first traversal starts from $s$, and it follows only the {\em outgoing} edges of each node, ignoring any edge {\em whose starting point ranks higher than the ending point}. For instance, if HoD traverses from the node $v_9$ in Figure~\ref{fig:over-steps}e, it would ignore the outgoing edge $\langle v_9, v_7\rangle$, since $r(v_9) = 4 > r(v_7) = 3$. As such, the first traversal of HoD never moves from a high-rank node to a low-rank node, and it terminates only when no higher-rank nodes can be reached. For each node $v$ visited, HoD maintains the distance from $s$ to $v$ along the paths that have been seen during the traversal, denoted as $\d(s, v)$.

Let $V'$ be the set of nodes that are not in the core graph of $\G$. The second traversal of HoD is performed as a linear scan of the nodes in $V'$, {\em in descending order of their ranks}. For each node $v' \in V'$ scanned, HoD inspects each {\em incoming} edge $e$ of $v'$, and then checks the starting point $u$ of the edge. For any such $u$, HoD calculates $\d(s, u) + l(e)$ as an upperbound of the distance from $s$ to $v'$. (Our solution guarantees that $u$ should have been visited by HoD before $v'$.) Once all incoming edges of $v'$ are inspected, HoD derives the distance from $s$ to $v'$ based on the upperbounds, and then it moves on to the next node in $V'$. This process terminates when all nodes in $V'$ are examined.

We illustrate the above query algorithm of HoD with an example.

%\parskip = \PAR

\vgap

\begin{example} \label{example:over-query}

\em Consider an SSD query from node $v_1$ in Figure~\ref{fig:over-steps}a. Given the augmented graph $\G$ in Figure~\ref{fig:over-steps}e, HoD first traverses $\G$ starting from $v_1$, following the outgoing edges whose ending points rank higher than the starting points. Since $v_1$ has only one outgoing edge $\langle v_1, v_9\rangle$, and since $v_9$ ranks higher than $v_1$, HoD moves from $v_1$ to $v_9$. $v_9$ has three outgoing edges: $\langle v_9, v_6\rangle$, $\langle v_9, v_7\rangle$, and $\langle v_9, v_{10}\rangle$. Among them, only $\langle v_9, v_{10}\rangle$ has an ending point that ranks higher than the starting point. Therefore, HoD moves from $v_9$ to $v_{10}$. $v_{10}$ has outgoing edges to three unvisited nodes, i.e., $v_3$, $v_5$, and $v_8$. Nevertheless, all of those nodes rank lower than $v_{10}$, and hence, they are ignored. As none of the remaining nodes can be reached without violating the constraints on node ranks, the first traversal of HoD ends. Based on the edges visited, HoD calculates $\d(v_1, v_9) = 1$ and $\d(v_1, v_{10}) = 4$.

\vgap

The second traversal of HoD examines the nodes {\em not} in the core graph in descending order of their ranks, i.e., it first examines $v_7$ and $v_8$ (whose ranks equal $3$), followed by $v_4$, $v_5$, and $v_6$ (whose ranks equal $2$), and finally $v_2$ and $v_3$ (whose ranks equal $1$), ignoring $v_1$ (as it is the source node of the query). $v_7$ has two incoming edges, $\langle v_6, v_7\rangle$ and $\langle v_9, v_7\rangle$. Among $v_6$ and $v_9$, only $v_9$ has been visited by HoD. Therefore, HoD calculates $\d(v_1, v_7) = \d(v_1, v_9) + l(\langle v_9, v_7\rangle) = 3$. Similarly, after inspecting $v_8$'s only incoming edge $\langle v_{10}, v_8\rangle$, HoD computes $\d(v_1, v_8) = \d(v_1, v_{10}) + l(\langle v_{10}, v_8\rangle) = 5$. The remaining nodes are processed in the same manner, resulting in
\vspace{-0mm}
\begin{align}
dist(v_1, v_4) \; &= \; dist(v_1, v_8) + l(\langle v_8, v_4\rangle)  \;\;\;\; = \, 6 \nonumber \\
dist(v_1, v_5) \; &= \; dist(v_1, v_{10}) + l(\langle v_{10}, v_5\rangle) \; = \, 5 \nonumber \\
dist(v_1, v_6) \; &= \; dist(v_1, v_{9}) + l(\langle v_{9}, v_6\rangle) \;\;\;\, = \; 2 \nonumber \\
dist(v_1, v_2) \; &= \; dist(v_1, v_{4}) + l(\langle v_{4}, v_1\rangle) \;\;\;\, = \; 7. \nonumber
\end{align}
\vspace{-0mm}
Observe that all the above distances computed from $\G$ are identical with those from the original graph in Figure~\ref{fig:over-steps}a. \done

%Meanwhile, the traversal from $v_4$ first inspects the only incoming edge of $v_4$, i.e., $\langle v_8, v_4\rangle$. Since $r(v_8) > r(v_4)$, HoD travels along the edge to visit $v_8$. $v_8$ has only one incoming edge $\langle v_8, v_{10}\rangle$, and $r(v_{10}) > r(v_8)$. Therefore, HoD traverses from $v_8$ to $v_{10}$. $v_{10}$ has two incoming edges, but only one of them (i.e., $\langle v_9, v_{10}\rangle$) has a starting point that ranks no lower than the ending point. As a consequence, HoD moves from $v_{10}$ to $v_9$. After that, the traversal terminates since, under the requirements on node ranks, no remaining nodes can be visited. The subgraph of $\G$ covered by the traversal is as illustrated in Figure~\ref{fig:over-traverse}c. Based on this subgraph, HoD derives that $dist(v_9, v_4) = 5$, $dist(v_{10}, v_4) = 2$, and $dist(v_8, v_4) = 1$.

%\begin{figure}[t]
%%\centering
%\centering
%\begin{tabular}{cc}
%%\hspace{-6mm} \multirow{3}{*}[6mm]{\includegraphics[width = 55mm]{./figures/ranks.eps}}
%%& \hspace{-2mm} \includegraphics[width = 30mm]{./figures/ppd1.eps} \\
%%& \hspace{-2mm} (b) Traversal from $v_1$. \\
%%& \hspace{-2mm} \gape[t]{\includegraphics[width = 26.5mm]{./figures/ppd2.eps}} \\
%%\hspace{-6mm} (a) Ranks of nodes. & \hspace{-2mm} (c) Traversal from $v_4$.
%\hspace{-3mm} \multirow{3}{*}[6mm]{\includegraphics[width = 55mm]{./figures/ranks.eps}}
%& \hspace{-2mm} \includegraphics[width = 30mm]{./figures/ppd1.eps} \\
%& \hspace{-2mm} (b) Traversal from $v_1$. \\
%& \hspace{-2mm} \gape[t]{\includegraphics[width = 26.5mm]{./figures/ppd2.eps}} \\
%\hspace{-3mm} (a) Ranks of nodes. & \hspace{-2mm} (c) Traversal from $v_4$.
%\end{tabular}
%\figcapup \caption{Traversals based on node ranks.} \figcapdown
%\label{fig:over-traverse}
%\end{figure}

\vgap

%Combining the distance information obtained from each traversal, HoD identifies that $v_9$ and $v_{10}$ are the only two nodes visited by both traversals. In addition, since $dist(v_1, v_9) + dist(v_9, v_4) = 6$, we can infer that the shortest path from $v_1$ to $v_4$ via $v_9$ has a length $6$. Similarly, $dist(v_1, v_{10}) + dist(v_{10}, v_4) = 6$, which indicates that the shortest path from $v_1$ to $v_4$ via $v_{10}$ also has a length $6$. Based on this, HoD reports that the distance from $v_1$ to $v_4$ equals $6$. This answer is correct, as $dist(v_1, v_4) = 6$ in the original graph (see Figure~\ref{fig:intro-hod}a). \done
\end{example}

\vgap

The query algorithm of HoD has an interesting property: the first traversal of the algorithm always visits nodes in ascending order of their ranks (as it never follows an edge that connects a high-rank node to low-rank node), while the second phase always visits nodes in descending rank order. Intuitively, if we maintain two copies of the augmented graph, such that the first (resp.\ second) copy stores nodes in ascending (resp.\ descending) order of their ranks, then HoD can answer any SSD query with a linear scan of the two copies. This leads to high query efficiency as it avoids random disk accesses. In Section~\ref{sec:pre}, we will elaborate how such two copies of the augmented graph can be constructed.


%Intuitively, the above query algorithm is efficient since HoD only visits a small portion of the augmented graph $\G$, due to the constraints on node ranks. For instance, in Example~\ref{example:over-query}, the traversal from $v_1$ (resp.\ $v_4$) only visits $3$ (resp.\ $4$) nodes, which is considerably smaller than the total number of nodes in $\G$. Furthermore, as HoD travels from low-rank nodes to high-rank nodes, we can arrange the nodes on the disk in ascending order of their ranks, so that the traversals of HoD can be performed with a linear scan of the nodes, without incurring a significant number of random disk accesses. One complication here is that, when two nodes $u$ and $v$ have the same rank, it might be difficult to decide (in advance) whether $u$ or $v$ would be traversed first by HoD. Fortunately, this complication can be easily avoided if $u$ and $v$ are not in the core graph $\GC$, as will be shown in Sections \ref{sec:sssp} and \ref{sec:p2p}. On the other hand, if $u$ and $v$ are in $\GC$, then we can address the issue by loading $\GC$ into the main memory, in which case HoD can traverse $u$ and $v$ in any order. (Recall that the preprocessing algorithm of HoD ensures that the core graph $\GC$ fits in the main memory.)
%
%\vgap
%
%One natural question is, why does HoD return correct results for PPD queries, given that it bypasses a large number of low-rank nodes during query processing? An intuitive answer is as follows. Recall that, for any node $v$ whose rank is not the highest (i.e., $v$ is not in the core graph $\GC$), $r(v) = i$ if and only if $v$ is removed by HoD in the $i$-th iteration of the reduction procedure (which constructs shortcuts in $\G$). In addition, when $v$ is removed in the $i$-th iteration, HoD creates shortcuts to ensure that the distance between any two remaining nodes would not be affected by the removal of $v$. In other words, if a node $u$ has a higher rank than $v$ (i.e., $u$ is not eliminated in the $i$-th iteration), then the distance from $u$ to other high-rank nodes should be retained by the shortcuts constructed by HoD. Therefore, even if HoD ignores $v$ during query processing, it can still correctly compute the distance from $u$ to other high-rank nodes. In Section~\ref{sec:p2p}, we will formalize this intuition and prove that HoD correctly answers PPD queries.
%
%\vgap
%
%%So far we have only discussed PPD queries.
%
%For SSD queries, HoD adopts an algorithm that also exploits the ranks of nodes to improve efficiency. In particular, given a source node $s$, HoD first traverses the augmented graph $\G$ from $s$ in exactly the same way as the traversal from $s$ for a PPD query, i.e., HoD follows the outgoing edges of each node, ignoring any edge with an ending point that ranks lower than its starting point. Upon termination of the traversal, HoD linearly scans the unvisited nodes in {\em descending} order of their ranks; for each node $v$, HoD derives the distance from $s$ to $v$ by inspecting the {\em incoming} edges of $v$. In Section~\ref{sec:sssp}, we will present the details of this algorithm and prove its correctness.




%% [inlined 3-Preprocessing/preprocessing.tex]
\section{Index Construction}\label{sec:pre}

As discussed in Section~\ref{sec:overview-pre}, the preprocessing algorithm of HoD takes as input a copy $G_0$ of the graph $G$, and it iteratively reduces $G_0$ into smaller graphs $G_1, G_2, \ldots$, during which it creates shortcuts to augment $G$. More specifically, the ($i{+}1$)-th ($i \ge 0$) iteration of the algorithm has four steps:
\begin{enumerate} [itemsep = 0.5mm, topsep = 2mm] %, partopsep = 0mm]
\item Select a set $R_i$ of nodes to be removed from $G_{i}$.

\item For each node $v \in R_i$, construct shortcuts in $G_{i}$ to ensure that the removal of $v$ does not change the distance between any two remaining nodes.

\item Remove the nodes in $R_i$ from $G_{i}$ to obtain a further reduced graph $G_{i+1}$. Store information about the removed nodes in the index file of HoD.

\item Pass the $G_{i+1}$ to the ($i{+}2$)-th iteration as input.
\end{enumerate}
In the following, we first elaborate Steps $2$ and $3$, and then clarify Step $1$. After that, we will discuss the termination condition of the preprocessing algorithm, as well as its space and time complexities.

For ease of exposition, we represent each edge $e = \langle u, v \rangle$ as a triplet $\left\langle u, v, l(e)\right\rangle$ or $\left\langle v, u, -l(e)\right\rangle$, where $l(e)$ is the length of $e$. For example, the edge $\langle v_9, v_7 \rangle$ in Figure~\ref{fig:over-steps}a can be represented as either $\left\langle v_9, v_7, 2\right\rangle$ or $\left\langle v_7, v_9, -2\right\rangle$. That is, a negative length in the triplet indicates that the second node in the triplet is the starting point of the edge. In addition, we assume that the input graph $G$ is stored on the disk as adjacency lists, such that (i) for any two nodes $v_i$ and $v_j$, the adjacency list of $v_i$ precedes that of $v_j$ if $i < j$, and (ii) each edge $\langle v_i, v_j\rangle$ with length $l$ is stored twice: once in the adjacency list of $v_i$ (as a triplet $\langle v_i, v_j, l\rangle$), and another in the adjacency list of $v_j$ (as a triplet $\langle v_j, v_i, -l\rangle$).

\subsection{Node Removal and Shortcut Generation} \label{sec:pre-removal}

Let $\v$ be a node to be removed from $G_{i}$. We define an {\em outgoing neighbor} of $\v$ as a node $u$ to which $\v$ has an outgoing edge. Similarly, an {\em incoming neighbor} of $\v$ is a node $w$ from which $\v$ has an incoming edge. We have the following observation:

%\parskip = \PAR

\vspace{-1mm}
\begin{observation} \label{observ:pre-neighbor}
\em
For any two nodes $v_j$ and $v_k$ in $G_i$, the distance from $v_j$ to $v_k$ changes after $\v$ is removed, if and only if the shortest path from $v_j$ to $v_k$ contains a sub-path $\langle u, \v, w\rangle$, such that $u$ (resp.\ $w$) is an incoming (resp.\ outgoing) neighbor of $\v$. \done
\end{observation}
\vspace{-1mm}


By Observation~\ref{observ:pre-neighbor}, we can preserve the distance between any two nodes in $G_i$ after removing $\v$, as long as we ensure that the distance between any incoming neighbor and any outgoing neighbor of $\v$ remains unchanged. This can be achieved by connecting the incoming and outgoing neighbors of $\v$ with shortcuts, as demonstrated in Section~\ref{sec:overview-pre}. Towards this end, a straightforward approach is to generate a shortcut $\langle u, w\rangle$ for any incoming neighbor $u$ and any outgoing neighbor $w$. The shortcuts thus generated, however, are often {\em redundant}. For example, consider the graph $G_i$ in Figure~\ref{fig:pre-contract}a. Suppose that we are to remove $v_2$, which has an incoming neighbor $v_1$ and an outgoing neighbor $v_3$. If we construct a shortcut from $v_1$ to $v_3$, it is useless since (i) $v_1$ already has an outgoing edge to $v_3$, and (ii) the edge $\langle v_1, v_3\rangle$ is even shorter than the path from $v_1$ to $v_3$ via $v_2$. As another example, assume that $v_4$ in Figure~\ref{fig:pre-contract}a is also to be removed. $v_4$ has an incoming neighbor $v_1$ and an outgoing neighbor $v_5$, but the path $\langle v_1, v_4, v_5\rangle$ is no shorter than another path from $v_1$ to $v_5$, i.e., $\langle v_1, v_3, v_5\rangle$, which does not go through $v_4$. As a consequence, even if we remove $v_4$ from $G_i$, the distance from $v_1$ to $v_5$ is still retained, and hence, it is unnecessary to insert a shortcut from $v_1$ to $v_5$.

%\multicolumn{2}{c}{\hspace{0mm} \includegraphics[height=3.3mm]{./figures/legend1.eps}}

In general, for any incoming neighbor $u$ and outgoing neighbor $w$ of $\v$, a shortcut from $u$ to $w$ is unnecessary if there is a path $P$ from $u$ to $v$, such that (i) $P$ does not go through $\v$, and (ii) $P$ is no longer than $\langle u, \v, w\rangle$. To check whether such a path $P$ exists, one may apply Dijkstra's algorithm to traverse $G_i$ from $u$ (or $w$), ignoring $\v$ during the traversal. However, when $G_i$ does not fit in main memory (as is often the case in the pre-computation process of HoD), this approach incurs significant overhead, due to the inefficiency of Dijkstra's algorithm for disk-resident graphs (as discussed in Section~\ref{sec:intro}). To address this issue, we adopt a heuristic approach that is not as effective (in avoiding redundant shortcuts) but much more efficient. Specifically, for each $\v$ in the node set $R_i$ to be removed from $G_i$, we generate a {\em candidate edge} $e_c = \langle u, w\rangle$ from each incoming neighbor $u$ of $\v$ to each outgoing neighbor $w$ of $\v$, setting the length of the shortcut to $l(\langle u, \v, w\rangle)$. For any such candidate edge $e_c$, we insert it into a temporary file $T$ as two triplets: $\langle u, w, l(e_c)\rangle$ and $\langle w, u, -l(e_c)\rangle$.

In addition to the candidate edges, we also insert two additional groups of edges (referred to as {\em baseline edges}) into the temporary file $T$ as triplets. The first group consists of any edge $e$ in $G_i$ connecting two nodes not in $R_i$, i.e., the two endpoints of $e$ are not to be removed. The second group is generated as follows: for each node $v$ not in $R_i$, we select $v$'s certain incoming neighbor $u'$ and outgoing neighbor $w'$, and we construct a baseline edge $\langle u', w'\rangle$, setting its length to $l(\langle u', v, w'\rangle)$.

The purpose of inserting a baseline edge $e$ into the temporary file $T$ is to help eliminate any redundant candidate edge that (i) shares the same endpoints with $e$ but (ii) is not shorter than $e$. Towards this end, once all baseline edges are added into $T$, we sort the triplets in $T$ using a standard algorithm for external sort, such that a triplet $t_1 = \langle v_a, v_b, l_1\rangle$ precedes another triplet $t_2 = \langle v_\alpha, v_\beta, l_2\rangle$, if any of the following conditions hold:
\begin{enumerate} [itemsep = 0.5mm, topsep = 2mm]%, partopsep = 0mm]
\item $a < \alpha$, or $a = \alpha$ but $b < \beta$.

\item $a = \alpha$, $b = \beta$, and $l_1 > 0 > l_2$. That is, any outgoing edge of a node precedes its incoming edges.

\item $a = \alpha$, $b = \beta$, $l_1 \cdot l_2 > 0$ (i.e., $t_1$ and $t_2$ are both incoming edges or both outgoing edges), and $|l_1| < |l_2|$. That is, $t_1$ and $t_2$ share the same starting and ending points, but $t_1$ is shorter than $t_2$.

\item $a = \alpha$, $b = \beta$, $l_1 \cdot l_2 > 0$, $|l_1| = |l_2|$, and $t_1$ is a baseline edge while $t_2$ is a candidate edge.
\end{enumerate}
Once $T$ is sorted, the outgoing (resp.\ incoming) edges with the same endpoints are grouped together, and the first edge in each group should have the smallest length within the group. If the first edge $e$ in a group
is a candidate edge, then we retain $e$ as it is shorter than any other baseline or candidate edges that we have generated. On the other hand, if $e$ is a baseline edge, then the distance between the endpoints of $e$ must not be affected by the removal of any nodes in $R_i$. In that case, all candidate edges in the group can be omitted. With one linear scan of the sorted $T$ and the adjacency lists of $G_i$, we can remove the information of any node in $R_i$, and merge the retained candidate edges into the adjacency lists of the remaining nodes. We illustrate the above algorithm with an example.

\begin{example} \label{example:pre-removal}
\em
Suppose that, given the graph $G_i$ in Figure~\ref{fig:pre-contract}a, we are to remove a node set $R_i = \{v_2, v_4\}$ from $G_i$. $v_2$ has only one incoming neighbor $v_1$ and one outgoing neighbor $v_3$, and $l(\langle v_1, v_2, v_3\rangle) = 2$. Accordingly, HoD generates a candidate edge $\langle v_1, v_3\rangle$ by inserting into the temporary file $T$ two triplets, $\langle v_1, v_3, 2\rangle$ and $\langle v_3, v_1, -2\rangle$. Similarly, for $v_4$, HoD creates a candidate edge $\langle v_1, v_5\rangle$, which is represented as two triplets in $T$: $\langle v_1, v_5, 2\rangle$ and $\langle v_5, v_1, -2\rangle$.

Meanwhile, the edge $\langle v_1, v_3\rangle$ in $G_i$ is selected as a baseline edge and is inserted into $T$, since neither $v_1$ nor $v_3$ is in $R_i$. In addition, HoD also generates a baseline edge $\langle v_1, v_5\rangle$ from the neighbors of $v_3$. This is because that (i) $v_1, v_3, v_5$ are not in $R_i$, (ii) $v_1$ is an incoming neighbor of $v_3$, and (iii) $v_5$ is an outgoing neighbor of $v_3$. Figure~\ref{fig:pre-contract}c illustrates the temporary file $T$ after all candidate and baseline edges are inserted, with some triplets omitted for simplicity. Figure~\ref{fig:pre-contract}d shows the file $T$ after it is sorted. The baseline edge $\langle v_1, v_3, 1\rangle$ precedes the candidate edge $\langle v_1, v_3, 2\rangle$, which indicates that we do not need to add a shortcut from $v_1$ to $v_3$. Similarly, the baseline edge $\langle v_3, v_1, -1\rangle$ precedes the candidate edge $\langle v_3, v_1, -2\rangle$, in which case the latter is omitted. Overall, each of the candidate edges in $T$ is preceded by a baseline edge, and hence, no shortcut will be created. Consequently, HoD removes from $G_i$ the adjacency lists of $v_2$ and $v_4$, as well as all edges to and from $v_2$ and $v_4$ in any other adjacency lists. This results in the reduced graph illustrated in Figure~\ref{fig:pre-contract}b. \done
\end{example}

In summary, HoD decides whether a candidate edge $e$ should be retained, by comparing it with all edges in $G_i$ as well as some two-hop paths in $G_i$. This heuristic approach may retain unnecessary candidate edges, but it does not affect the correctness of SSD queries. To understand this, recall that each candidate edge $e = \langle u, w\rangle$ has the same length with a certain path $\langle u, \v, w\rangle$ that {\em exists} in $G_i$, where $\v$ is the node whose removal leads to the creation of $e$. In other words, the length of $e$ is at least larger than or equal to the distance from $u$ to $w$. Adding such an edge into $G_i$ would not decrease the distance between any two nodes in $G_i$, and hence, retaining $e$ does not change the results of any SSD queries.

The above discussions assume that HoD has selected a set $R_i$ of nodes to be removed from $G_i$, and has decided which baseline edges are to be generated from the neighbors of the nodes not in $R_i$. We will clarify these two issues in Section \ref{sec:pre-select} and \ref{sec:pre-baseline}.

%\parskip = \PAR


\subsection{Selecting Nodes for Removal} \label{sec:pre-select}

Consider any node $v$ in $G_i$. Intuitively, if the removal of $v$ requires us to insert a large number of shortcuts into $G_i$, then $v$ may lie on the shortest paths between many pair of nodes, in which case $v$ should be considered important. Let $B_{in}$ and $B_{out}$ be the set of incoming and outgoing neighbors of $v$, respectively. The maximum number of shortcuts induced by $v$'s removal is:
\begin{equation} \label{eqn:pre-score}
s(v) = {\big |}B_{in}{\big |} \cdot {\big |}B_{out} \setminus B_{in}{\big |} + {\big |}B_{out}{\big |} \cdot {\big |}B_{in} \setminus B_{out}{\big |}.
\end{equation}
We refer to $s(v)$ as the {\em score} of $v$ in $G_i$, and we consider $v$ {\em unimportant} if $s(v)$ is no more than the median score in $G_i$. (For practical efficiency, we use an approximated value of the median score computed from a sample set of the nodes.)

%\parskip = \PAR

Ideally, we would like to remove all unimportant nodes from $G_i$, but this is not always feasible. To explain, consider that we are given the reduced graph $G_1$ in Figure~\ref{fig:over-steps}b, and we aim to eliminate both $v_6$ and $v_7$. $v_6$ has only one incoming neighbor $v_9$ and one outgoing neighbor $v_7$, and hence, HoD creates one candidate edge $\langle v_9, v_7\rangle$, setting its length to $2$ (i.e., the length of the path $\langle v_9, v_6, v_7\rangle$). Similarly, for $v_7$, HoD generates a candidate edge $\langle v_6, v_{10}\rangle$. These two candidate edges are intended to preserve the distance between any two nodes in $G_i$ after $v_6$ and $v_7$ are removed. However, none of the two candidate edges is valid if both $v_6$ and $v_7$ are eliminated. In particular, $\langle v_9, v_7\rangle$ points from $v_9$ to $v_7$, i.e., it connects $v_9$ to a node that no longer exists. To avoid the above error, whenever HoD chooses to delete a node $v$ from $G_i$, it will retain all neighbors of $v$ in $G_i$, even if some neighbor might be unimportant. For example, in Figure~\ref{fig:over-steps}b, if HoD decides to remove $v_6$, then it will prevent $v_7$ from being removed at the same time, and vice versa.

%This does imply that the nodes removed from $G_i$ depend on the order in which HoD examines the nodes. Nonetheless, we find that the order of node examination does not have a significant effect on the performance of HoD.

\begin{figure}[t]
\begin{small}
\centering
\begin{tabular}{ccc}
\multicolumn{3}{c}{\includegraphics[height=3mm]{./figures/legend2.eps}} \vspace{1mm}\\
\includegraphics[width = 15mm]{./figures/contract1.eps} & \multirow{3}{*}[16.5mm]{\includegraphics[width = 14mm]{./figures/contract3.eps}} & \multirow{3}{*}[16.5mm]{\includegraphics[width = 14mm]{./figures/contract4.eps}} \vspace{1mm}\\
(a) Graph $G_i$. & & \vspace{4mm}\\
\includegraphics[width = 14mm]{./figures/contract2.eps} & & \\
(b) $v_2, v_4$ removed. & \hspace{3mm} (c) Before sorting. & \hspace{3mm} (d) After sorting. \\
\end{tabular}
\figcapup \caption{Node removal and shortcut generation.} \figcapdown \vspace{-1mm}
\label{fig:pre-contract}
\end{small}
\end{figure}

\subsection{Generation of Baseline Edges} \label{sec:pre-baseline}

As mentioned in Section~\ref{sec:pre-removal}, a baseline edge generated by the preprocessing algorithm of HoD is either (i) an edge in $G_i$ whose endpoints are not to be removed, or (ii) an artificial edge $\langle u, w\rangle$ that corresponds to certain two-hop path $\langle u, v, w\rangle$ in $G_i$, such that none of $u, v, w$ is to be removed. The construction of baseline edges from two-hop paths is worth discussing. First, given that there exists an enormous number of two-hop paths in $G_i$, it is prohibitive to convert each two-hop path into a baseline edge. Therefore, we only select a subset of the two-hop paths in $G_i$ for baseline edge generation. In particular, the total number of two-hop paths selected is set to $c \cdot \sum_{v \in R_i} s(v)$, where $c$ is a small constant, $s(v)$ is as defined in Equation~\ref{eqn:pre-score}, and $\sum_{v \in R_i} s(v)$ is the total number of candidate edges induced by the removal of nodes in $R_i$. In other words, we require that the number of baseline edges generated from two-hop paths is at most $c$ times the number of candidate edges. In our implementation of HoD, we set $c = 5$.

%To explain, observe that $\sum_{v \in R_i} s(v)$ is the total number of candidate edges induced by the removal of nodes in $R_i$. In other words, we require that the number of baseline edges generated from two-hop paths should be $c$ times the number of candidate edges. In our implementation of HoD, we set $c = 5$.

Those $c \cdot \sum_{v \in R_i} s(v)$ baseline edges are generated as follows. First, we randomly choose an edge in $G_i$, and arbitrarily select an endpoint $v$ of the edge that is not in $R_i$. (Note that such an endpoint always exists.) After that, from the incoming (resp.\ outgoing) neighbors of $v$, we randomly select a node $u$ (resp.\ $w$), and we generate a baseline edge $\langle u, w\rangle$, setting its length to $l(\langle u, v, w\rangle)$. As such, if a node $v$ is adjacent to a large number of edges, then it has a high chance of being selected to produce baseline edges. This is intuitive since such a node $v$ tends to lie on the shortest paths between many pairs of nodes, and hence, the baseline edges generated from $v$ may be more effective in eliminating redundant shortcuts.

\subsection{Termination Condition} \label{sec:pre-terminate}

As mentioned in Section~\ref{sec:overview}, HoD requires that the core graph $\GC$ fits in the main memory, where $\GC$ is the reduced graph obtained in the last iteration of HoD's preprocessing algorithm. Accordingly, we do not allow the pre-computation procedure of HoD to terminate before the reduced graph $G_i$ has a size no more than $M$. In addition, even after $G_i$ fits in the main memory, we will still continue the preprocessing procedure, until the size of $G_i$ is reduced by less than $5\%$ in an iteration of the processing algorithm. This is intended to reduce the size of the core graph $\GC$ to improve query efficiency, as will be explained in Section~\ref{sec:sssp}.

\begin{figure}[t]
\centering
\begin{small}
\begin{tabular}{ccc}
\hspace{-3mm}
\includegraphics[width = 30mm]{./figures/forward.eps} &
\hspace{-3mm}
\includegraphics[width = 18mm]{./figures/legend3.eps} &
%\hspace{-2mm}
\includegraphics[width = 30mm]{./figures/backward.eps} \\
\hspace{-3mm} (a) $\GF$. & \multicolumn{2}{c}{\hspace{15mm} (b) $\GB$.}
\end{tabular}
\figcapup \vspace{-2mm} \caption{Forward graph $\GF$ and backward graph $\GB$.} \figcapdown \vspace{-1mm}
\label{fig:pre-graphs}
\end{small}
\end{figure}

\subsection{Index File Organization} \label{sec:pre-file}
%In general, a smaller $\GC$ leads to higher query efficiency, as will be clarified in Section~\ref{sec:sssp}. Therefore, even after the reduced graph $G_i$ starts to fit in memory, we still allow the preprocessing algorithm to proceed, so as to remove more nodes from $G_i$. However, if the size of $G_i$ is reduced by less than $5\%$ after an iteration of the algorithm, we consider $G_i$ sufficiently small and we terminate the algorithm.

%\parskip = \PAR

Once the preprocessing procedure completes, the core graph $\GC$ is written to the disk in the form of adjacency lists. Meanwhile, the adjacency list of each node not in $\GC$ is separated into two parts that are stored in two different files, $\FF$ and $\FB$. These two files are created at the beginning of HoD's preprocessing algorithm, and they are initially empty. Whenever a node $v$ is removed from the reduced graph $G_i$, we inspect the adjacency list of $v$ in $G_i$, and we append all of $v$'s outgoing (resp.\ incoming) edges to $\FF$ (resp.\ $\FB$). Upon termination of the preprocessing procedure, we reverse the order of nodes in $\FB$, but retain the order of nodes in $\FF$. We refer to the graph represented by $\FF$ as the {\em forward graph}, denoted as $\GF$. Meanwhile, we refer to the graph represented by $\FB$ as the {\em backward graph}, denoted as $\GB$. When combined, $\GC$, $\GF$, and $\GB$ form the augmented graph that is used by HoD for query processing. For example, for the augmented graph in Figure~\ref{fig:over-steps}a, its core graph is as illustrated in Figure~\ref{fig:over-steps}e, while its forward and backward graphs are as shown in Figure~\ref{fig:pre-graphs}. For ease of exposition, we will abuse notation and use $\GF$ (resp.\ $\GB$) to refer to both $\GF$ (resp.\ $\GB$) and its underlying file structure $\FF$ (resp.\ $\FB$).

% As such, when the preprocessing procedure ends, $\FF$ (resp.\ $\FB$) would contain the outgoing (resp.\ incoming) edges of all nodes not in $\GC$.





%\begin{figure}[t]
%%\centering
%\centering
%\begin{tabular}{ccc}
%\hspace{-6mm} \includegraphics[width = 30mm]{./figures/forward.eps} & \hspace{-6mm} \includegraphics[width = 10mm]{./figures/legend3.eps} & \hspace{-2mm} \includegraphics[width = 30mm]{./figures/backward.eps} \\
%%\hspace{-6mm} \multicolumn{3}{c}{(a) $\GF$. $\qquad$ (b) $\GB$ (equivalently, $\GR$).}
%%\hspace{-6mm} (a) $\GF$. & \hspace{-2mm} \multicolumn{2}{c}{\hspace{0mm} \includegraphics[height=3.3mm]{./figures/legend1.eps}} %\multicolumn{2}{c}{(b) $\GB$ (equivalently, $\GR$).}
%\end{tabular}
%\figcapup \caption{Forward graph $\GF$, backward graph $\GB$, and reversed graph $\GR$.} \figcapdown
%\label{fig:pre-graphs}
%\end{figure}


$\GF$ and $\GB$ have two interesting properties. First, all nodes in $\GF$ (resp.\ $\GB$) are stored in ascending (resp.\ descending) order of their ranks. To explain this, recall that any node $v$ removed in the $i$-th iteration of the preprocessing algorithm has a rank $r(v)=i$. Consequently, if a node $u$ is stored in $\GF$ before another node $w$, then $r(u) \le r(w)$. As for $\GB$, since we reverse the order of all edges in $\GB$ upon termination of the preprocessing produce, we have $r(u) \ge r(w)$ for any node $u$ that precedes another node $w$ in $\GB$. Second, for any node $v$, its edges in $\GF$ and $\GB$ only connect it to the nodes whose ranks are strictly higher than $v$. This is because, by the time $v$ is removed from the reduced graph, all nodes that rank lower than $v$ must have been eliminated from the reduced graph, and hence, any edge in $v$'s adjacency list only links $v$ to the nodes whose rank is at least $r(v)$. Meanwhile, any neighbor $u$ of $v$ in the reduced graph should have a rank higher than $r(v)$. Otherwise, we have $r(u) = r(v)$, which, by the definition of node ranks, indicates that $u$ and $v$ are removed in the same iteration of the preprocessing algorithm. This is impossible as the pre-computation procedure of HoD never eliminates two adjacent nodes in the same iteration, as explained in Section~\ref{sec:pre-removal}. In Section~\ref{sec:sssp}, we will show how HoD exploits the above two properties of $\GF$ and $\GB$ to efficiently process SSD queries.

%For SSD queries, HoD requires a {\em reversed} version of $\GB$ that stores the adjacency lists of high-rank nodes {\em before} those of low-rank nodes. For this purpose, we add an extra step in the preprocessing algorithm of HoD: we scan $\GB$ and write its adjacency lists into a file $\FR$ in reverse order. The resulting file represents a graph that is identical to $\GB$. To avoid confusion, we refer to this graph as the {\em reversed} graph, and denote it as $\GR$.



%\begin{figure}[t]
%\centering
%%\begin{small}
%\begin{tabular}{ccc}
%\hspace{-3mm}
%\includegraphics[width = 30mm]{./figures/forward.eps} &
%\hspace{-3mm}
%\includegraphics[width = 18mm]{./figures/legend3.eps} &
%%\hspace{-2mm}
%\includegraphics[width = 30mm]{./figures/backward.eps} \\
%\hspace{-3mm} (a) $\GF$. & \multicolumn{2}{c}{\hspace{15mm} (b) $\GB$ (equivalently, $\GR$).}
%\end{tabular}
%\figcapup \caption{Forward graph $\GF$, backward graph $\GB$, and reversed graph $\GR$.} \figcapdown
%\label{fig:pre-graphs}
%\end{figure}

\subsection{Cost Analysis} \label{sec:pre-complexity}

The preprocessing algorithm of HoD requires $O(n)$ main memory, where $n$ is the number of nodes in $G$. This is because (i) when removing a node $v$ from the reduced graph, HoD needs to record the neighbors of $v$ and exclude them from the node removal process, and (ii) $v$ may have $O(n)$ neighbors. Other parts of the preprocessing algorithm do not have a significant memory requirement.

The major I/O and CPU costs of the preprocessing algorithm are incurred by sorting the edge triplets in each iteration. In the worst case when the input graph $G$ is a complete graph, there are $O(n^2)$ triplets generated in each iteration, leading to a prohibitive I/O and CPU overhead. Fortunately, real-world graphs are seldom complete graphs, and they tend to contain a large number of nodes with small degrees. In that case, each iteration of HoD's preprocessing procedure would only generates a moderate number of edge triplets, leading to a relatively small overhead.
%In Section~\ref{sec:exp}, we will demonstrate that HoD's pre-computation procedure is efficient on large real data.

Lastly, the space consumption of HoD's index is $O(n^2)$ since, in the worst case, HoD may construct a shortcut from each node $v$ to every node that ranks higher than $v$. This space complexity is unfavorable, but it is comparable to the space complexity of VC-Index \cite{ckcc12}. In addition, as shown in our experiments, the space requirement of HoD in practice is significantly smaller than the worst-case bound.



%\begin{figure}[t]
%\centering
%\begin{tabular}{cc}
%\multicolumn{2}{c}{\hspace{-4mm} \includegraphics[height=3.3mm]{./figures/legend1.eps}} \vspace{1mm} \\
%\hspace{-2mm} \includegraphics[width=26.6mm]{./figures/pre1.eps}
%&
%\hspace{3mm} \includegraphics[width=26.6mm]{./figures/pre2.eps} \\
%\hspace{-2mm} (a) After removing $v_1, v_2, v_3$. & \hspace{3mm} (b) After removing $v_4$.  \vspace{1mm}\\
%\hspace{-2mm} \includegraphics[width=26mm]{./figures/pre3.eps}
%&
%\hspace{3mm} \includegraphics[width=25mm]{./figures/pre4.eps} \\
%\hspace{-2mm} (c) After removing $v_5,v_6$. & \hspace{3mm} (d) After removing $v_7$.
%\end{tabular}
%\figcapup \caption{Graph reduction and shortcut construction.} \figcapdown
%\label{fig:over-steps}
%\end{figure}




%This section presents {\em FC} (\underline{f}irst-\underline{c}ut), an index structure designed for road networks with small arterial dimensions. FC is worst-case efficient for distance queries, and its space consumption is modest; nevertheless, FC is unsuitable for large road networks as it incurs significant pre-processing cost. The reasons that we introduce FC are (i) it is a conceptually simple method that demonstrates the key idea of our proposal, and (ii) with a few modifications and optimizations, FC can be turned into a scalable method that handles both distance and shortest path queries (see Section~\ref{sec:ah}).
%
%%In what follows, we will first present an overview of FC (Section~\ref{sec:fc-pre}), and then present the query processing algorithms (Section~\ref{sec:fc-query}), followed by a complexity analysis (Section~\ref{sec:fc-complexity}) and a correctness proof (Section~\ref{sec:fc-correctness}).
%
%\subsection{Index Construction} \label{sec:fc-pre}
%
%Given a road network $G$, FC first assigns a {\em level} to each node in $G$, such that nodes with higher levels tend to be more {\em important}. After that, FC organizes the nodes into a hierarchy based on their levels, and it adds auxiliary edges between various nodes to facilitate query processing. In the following, we will elaborate how the node levels are decided and how the auxiliary edges are created.
%
%\header {\bf Deciding Node Levels.} First, FC imposes on $G$ a ($4 {\times} 4$)-cell square grid that tightly covers all nodes in $G$. After that, FC recursively splits each grid cell into $2 \times 2$ smaller cells, until each cell contains at most one node in $G$. This results in a sequence of square grids with increasing resolutions. Let $h$ be the number of grids thus constructed. We use $R_i$ to denote the grid with $2^{h+2 - i} \times 2^{h+2 - i}$ cells, i.e., $R_h$ is the ($4 {\times} 4$)-cell grid that FC first constructed, and $R_1$ is the grid with the highest resolution.
%
%Let $d_{max}$ (resp.\ $d_{min}$) be the largest (resp.\ smallest) $L_\infty$ distance between any two nodes in $G$. It can be verified that $h \le \log_2 (d_{max}/d_{min}) - 1$. We note that $h$ is always a small number for practical road networks: Even if $d_{max}$ is as large as the length of the Equator ($\approx 4 \times 10^7$ meters) and $d_{min}$ is as small as $1$ meter, the value of $h$ is no more than $26$.
%
%Given each $R_i$ ($i \in [1, h]$), FC computes the arterial edges in any ($4 {\times} 4$)-cell region in $R_i$. Let $A_i$ be the set of arterial edges obtained from $R_i$. For any edge in $A_i$, we define it as a {\em level-$i$ edge} if it does not appear in $A_{i+1}, \ldots, A_h$. If an edge does not appear in any $A_i$, then we refer to it as a {\em level-$0$ edge}. In other words, an edge has a higher level if it is an arterial edge for a larger region. Similarly, we also define the level of each node $v$ in $G$: we say that $v$ is a {\em level-$i$ node} if it is adjacent to some edge at level $i$ but not any edge at level $i+1, \ldots, h$. Intuitively, a higher-level node tend to be more important for shortest path and distance queries.
%
%\header {\bf Creation of Shortcuts.} Once the node levels are decided, FC organizes the nodes in $G$ into a hierarchy $\G$ of $h+1$ levels $L_0, L_1, \ldots, L_h$, such that all level-$i$ ($i \in [0, h]$) nodes are contained in $L_i$. For example, Figure~\ref{fig:intro-hierarchy} illustrates a $3$-level hierarchy of the nodes in Figure~\ref{fig:intro-original}.
%
%The hierarchy $\G$ retains all edges in $G$. In addition, FC inserts into $\G$ some auxiliary edges, referred to as {\em shortcuts}. For any two nodes $v_s$ and $v_t$, FC creates a shortcut $c$ from $v_s$ to $v_t$, if the shortest path from $v_s$ to $v_t$ only passes through nodes whose levels are lower than both $v_s$'s and $v_t$'s. Furthermore, the length of $c$ equals to the distance from $v_s$ to $v_t$, i.e., $l(c) = dist(v_s, v_t)$. For instance, consider the nodes $v_6, v_8, v_9, v_{10}$ in Figure~\ref{fig:intro-hierarchy}, whose levels are $0$, $1$, $1$, and $2$, respectively. There is a shortcut from $v_9$ to $v_{10}$, since the shortest path from $v_9$ to $v_{10}$ only goes through $v_6$, and the level of $v_6$ is lower than those of $v_9$ and $v_{10}$. On the other hand, there is no shortcut from $v_8$ to $v_9$, since the shortest path from $v_8$ to $v_9$ passes through $v_{10}$, whose level is higher than both $v_8$'s and $v_9$'s.
%
%The shortcuts inserted into $\G$ enable us to avoid visiting unimportant nodes when processing distance queries. For example, given the shortcut $c$ from $v_9$ to $v_{10}$ in Figure~\ref{fig:intro-hierarchy}, we can determine that $dist(v_9, v_{10}) = l(c)$, without having to compute the actual shortest path from $v_9$ to $v_{10}$. In general, for any two nodes $v_s$ and $v_t$ in $\G$, there exists a path from $v_s$ to $v_t$ that bypasses unimportant nodes with shortcuts, as will be explained in Section~\ref{sec:fc-query}. For simplicity, we will use the term ``edge'' to refer to either an original edge or a shortcut in $\G$, unless otherwise specified.
%
%\subsection{Query Processing} \label{sec:fc-query}
%
%Consider a query $\QC$ that asks for the distance from a node $s$ in $G$ to another node $t$. Given the node hierarchy $\G$, FC answers $\QC$ with two concurrent traversals of $\G$ that start from $s$ and $t$, respectively. Each traversal is performed using a {\em constrained} version of Dijkstra's algorithm \cite{d59}, as explained in the following.
%
%\header {\bf Traversal Algorithm.}  The traversal from $s$ maintains a hash table $T_s$ and a priority queue $Q_s$. The hash table $T_s$ maps each node $v$ in $G$ to a value $\kappa_s(v)$, which equals the length of the shortest path from $s$ to $t$ that has been found so far. Initially, we have $\kappa_s(s) = 0$ and $\kappa(v) = +\infty$ for any other node $v$.
%
%Meanwhile, each entry in the priority queue $Q_s$ corresponds to a certain node $v'$ in $G$, and the key of the entry equals $\kappa_s(v')$. In the beginning of the traversal, $Q_s$ contains only one entry, which corresponds to $s$. Subsequently, FC iteratively extracts (from $Q_s$) the node $u$ with the smallest key. For each $u$ extracted, FC inspects every edge $\langle u, v\rangle$ in $\G$ that starts from $u$, and it checks whether $v$ satisfies certain constraints. (We will clarify these constraints shortly). If $v$ violates any of the constraints, it would be ignored; otherwise, FC would further check whether $\kappa_s(u) + l(\langle u, v\rangle) < \kappa_s(v)$, i.e., whether the path from $s$ to $v$ via $u$ is shorter than all known paths from $s$ to $v$. If the inequality holds, then FC sets $\kappa_s(v) = \kappa_s(u) + l(\langle u, v\rangle)$ and inserts $v$ into $Q_s$ (if $v$ has not been inserted before).
%
%The traversal from $t$ also maintains a hash table $T_t$ and a priority queue $Q_t$. It is performed in a manner similar to the traversal from $s$, with one notable difference: Whenever FC extracts a node $u$ from $T_t$, it only inspects the edges $\langle v, u\rangle$ that {\em points to} $u$. In other words, the traversal from $t$ focuses on paths that end at $t$.
%
%FC conducts the above two traversals in an round-robin fashion, i.e., it extracts nodes from the two priority queues $Q_s$ and $Q_t$ in turns. To determine when the traversals can be terminated, FC maintains a variable $\theta$ that records the length of the shortest path from $s$ to $t$ that is seen so far. Initially, $\theta = +\infty$. After that, for each node $u$ extracted from either priority queue, FC retrieves its key $\kappa_s(u)$ in the hash table $T_s$, as well as its key $\kappa_t(u)$ in $T_t$. Recall that $\kappa_s(u)$ (resp.\ $\kappa_t(u)$) records the length of the shortest path from $s$ to $u$ (resp.\ from $u$ to $t$) found so far. Therefore, the shortest path from $s$ to $t$ should be no longer than $\kappa_s(u) + \kappa_t(u)$. Based on this, if $\kappa_s(u) + \kappa_t(u) < \theta$, then FC would update $\theta$ and set it to $\kappa_s(u) + \kappa_t(u)$.
%
%Whenever $\theta$ is no more than the smallest key value in $Q_s$, we know that $dist(s, u) \ge \theta$ for any node $u$ remaining in $Q_s$, which indicates that $u$ cannot be on the shortest path from $s$ to $t$. In that case, FC would terminate the traversal from $s$. Similarly, the traversal from $t$ is stopped if $\theta$ is no more than any key values in $Q_t$. When both traversals are terminated, FC returns $\theta$ as the answer to the distance query.
%
%\header {\bf Constraints on Node Traversals.} As mentioned above, whenever FC extracts a node $u$ from a priority queue (either $Q_s$ or $Q_t$), it inspects the neighbors of $u$, and it processes only those neighbors $v$ that satisfy certain constraints. Specifically, there are two constraints on $v$:
%\begin{enumerate} %[itemsep = 1mm]
%\item {\em Level Constraint:} $v$ should not be at a level lower than $u$'s.
%
%\item {\em Proximity Constraint:} Let $i$ be the level of $v$ ($i \in [0, h]$). If $u$ is extracted from $Q_s$ (resp.\ $Q_t$), then $v$ and $s$ (resp.\ $t$) should be covered in the same ($3 {\times} 3$)-cell region in $R_{i+1}$. (Recall that $R_i$ is a square grid with $2^{h+2 - i} \times 2^{h+2 - i}$ cells.)
%\end{enumerate}
%
%Both of the above constraints are intended to improve the efficiency of FC. In particular, the level constraint helps FC bypass unimportant nodes during query processing. For example, consider that we use FC to compute the distance from $v_8$ to $v_{11}$. As explained previously, FC would invoke two traversals starting from $v_8$ and $v_{11}$, respectively. Since $v_{11}$ is at level $3$, the traversal from $v_{11}$ would only visit the neighbors of $v_{11}$ that are at levels no lower than $3$. As a consequence, $v_{10}$ is visited (since its level equals $3$), while $v_1$ and $v_9$ are bypassed. After that, the traversal would visit any neighbor of $v_{10}$ whose level is no lower than that of $v_{10}$. Since none of the neighbors of $v_{10}$ fulfills this requirement, the traversal terminates. Similarly, the traversal from $v_8$ would first visit two of $v_8$'s neighbors, $v_7$ and $v_{10}$, ignoring the remaining neighbor $v_3$, since $v_3$'s level is lower than that of $v_8$. After that, the traversal visits only $v_{11}$ and terminates, as all other remaining nodes violate the level constraint. In summary, the two traversals by FC visit only four nodes: $v_7$, $v_8$, $v_{10}$, and $v_{11}$; all other nodes in $\G$ are omitted.
%
%\begin{figure}[t]
%%\begin{small}
%\centering
%\includegraphics[width = 2.8in]{figures/proximity.eps}
%\vspace{-1mm} \figcapup \caption{Illustration of the proximity constraint.} \label{fig:fc-proximity} \figcapdown \vspace{-1mm}
%\end{figure}
%
%
%%Meanwhile, the proximity constraint ensures that FC only searches a small number of grid cells in each level of the node hierarchy $\G$. For example, suppose that we are given the node hierarchy in Figure~\ref{fig:fc-proximity}, and we use FC to compute the distance from $v_1$ to $v_6$. Among the two traversals invoked by FC, the one starting from $v_1$ would first visit $v_2$, and then $v_3$ and $v_7$. The node $v_7$ has two unvisited neighbors, $v_8$ and $v_{10}$, both of which are at the same level as $v_7$, i.e., they satisfy the level constraint. However, they would be ignored by FC, as they violate the proximity constraint. In particular, both $v_8$ and $v_{10}$ are level-$1$ nodes, but there does not exist any $(2 \times 2)$-cell region in $R_2$ that can cover $v_1$ and $v_8$ (resp.\ $v_{10}$) simultaneously. Note that ignoring $v_8$ and $v_{10}$ does not affect the correctness of the query result, since neither $v_8$ nor $v_{10}$ connects $v_1$ to $v_6$. In contrast, the node $v_4$, which is a neighbor of $v_3$, would be visited by FC as it satisfies both the level and proximity constraints. Specifically, the level of $v_4$ equals $2$, which is no less than that of $v_3$; furthermore, $v_4$ and $v_1$ is contained in the same ($4 {\times} 4$)-cell region in $R_3$.
%
%Meanwhile, the proximity constraint ensures that FC only searches a small number of grid cells in each level of the node hierarchy $\G$. For example, suppose that we are given the node hierarchy in Figure~\ref{fig:fc-proximity}, and we use FC to compute the distance from $v_1$ to $v_6$. Among the two traversals invoked by FC, the one starting from $v_1$ would first visit $v_2$, and then $v_3$ and $v_7$. The node $v_7$ has a neighbors $v_8$, which is at the same level as $v_7$, i.e., $v_8$ satisfies the level constraint. However, $v_8$ would still be ignored by FC, as it violates the proximity constraint. In particular, $v_8$ is a level-$1$ node, but there does not exist any ($3 {\times} 3$)-cell region in $R_2$ that can cover $v_1$ and $v_8$ simultaneously. In contrast, the node $v_4$, which is a neighbor of $v_3$, would be visited by FC as it satisfies both the level and proximity constraints. Specifically, the level of $v_4$ equals $2$, which is no less than that of $v_3$; furthermore, $v_4$ and $v_1$ is contained in the same ($3 {\times} 3$)-cell region in $R_3$. Note that, although FC ignores $v_8$, the correctness of the query result is not affected, since $v_8$ is not on the shortest path from $v_1$ to $v_6$.
%
%In general, the proximity constraint guarantees that in each level $i$ of the node hierarchy, FC only traverses the nodes contained in two ($5 {\times} 5$)-cell regions, which are centered at the source $s$ and destination $t$ of the query, respectively. In particular, the region centered at $s$ (resp.\ $t$) is the union of all ($3 {\times} 3$)-cell regions that cover $s$ (resp.\ $t$). This, when combined with the level constraint, ensures that FC is worst-case efficient in terms of query time, as will be shown in Section~\ref{sec:fc-complexity}.
%
%
%%the proximity constraint prevents FC from visiting a node $v$ too far away from $s$ and $t$, unless $v$ is of high importance. Intuitively, this is consistent with the way that we drive between two distant locations $s$ and $t$ in practice: In the beginning, we may start from some small streets connected to $s$; but as we move further away from $s$, we turn into the main roads and no longer traverse the small streets, i.e., we focus on more important road network nodes as we travel away from $s$. After a while, we may further move on to the highways; but upon approaching $t$, we will exit the highways to get into some main roads, which then lead us to some small streets connected with $t$. In other words, we would traverse less important nodes only when they are close to $t$.
%
%\subsection{Complexity Analysis} \label{sec:fc-complexity}
%
%In this section, we will prove that FC takes $O(h n)$ space, and it answers any distance query in $O(h^2)$ time, where $h$ is the maximum level in the node hierarchy $\G$, and $n$ is the number of nodes in the road network $G$. In addition, we will discuss the pre-computation time of FC.
%
%\header {\bf Query Time.} As explained in Section~\ref{sec:fc-query}, FC answers any distance query by two traversals on the node hierarchy $\G$, starting from the source $s$ and destination $t$ of the query, respectively. Due to the level and proximity constraints, each traversal of FC visits any level of $\G$ at most once; in addition, for the $i$-th level ($i \in [1, h]$), each traversal only examines the nodes in a ($5 {\times} 5$)-cell region in the grid $R_{i+1}$. A natural question is: How many level-$i$ nodes are there in the ($5 {\times} 5$)-cell region? The following lemma provides an answer.
%%\begin{lemma} \label{lmm:fc-node-num}
%%Any ($\alpha {\times} \alpha$)-cell region in $R_i$ contains $O(\alpha^2 \a)$ level-$i$ nodes in $\G$, where $\a$ is the arterial dimension of $G$.
%%\end{lemma}
%\begin{lemma} \label{lmm:fc-node-num}
%Any ($\alpha {\times} \alpha$)-cell region in $R_i$ contains $O(\alpha^2 \a)$ level-$i$ nodes in $\G$, where $\a$ is the arterial dimension of $G$.
%\end{lemma}
%\begin{proof}
%The proofs for all lemmas and theorems in this paper can be found in our anonymous technical report online \cite{AhReport}.
%\end{proof}
%To explain the rationale behind Lemma~\ref{lmm:fc-node-num}, recall that each level-$i$ node in $\G$ is adjacent to an arterial edge in a ($4 {\times} 4$)-cell region in $R_i$. Furthermore, each ($4 {\times} 4$)-cell region in $R_i$ has at most $\a$ arterial edges. For any ($\alpha \times \alpha$)-cell region in $R_i$, it can overlap with $O(\alpha^2)$ regions of $4 \times 4$ cells, and hence, it contains $O(\alpha^2 \a)$ level-$i$ nodes.
%
%Observe that any ($5 {\times} 5$)-cell region in $R_{i+1}$ corresponds to a ($10 \times 10$)-cell region in $R_i$. By Lemma~\ref{lmm:fc-node-num}, this region contains $O(\a)$ level-$i$ node in $\G$. In other words, the number of level-$i$ nodes visited by each traversal of FC is $O(\a)$. Given that $\G$ has $h+1$ levels, the total number of nodes traversed by FC is $O(h\a)$.
%
%Next, we will show that each node in $\G$ has $O(\a)$ edges that satisfy the level constraint. (The edges that violate the constraint can be removed from $\G$ beforehand, as they would never be traversed by FC for any query). Consider any level-$i$ node $u$, and any node $v$ whose level is at least $i$. By the way that $\G$ is constructed, there is a shortcut connecting $u$ to $v$, if and only if the shortest path between $u$ and $v$ only goes through nodes at levels lower than $i$. Intuitively, this indicates that $u$ and $v$ should not be too far apart from either other; otherwise, the shortest path between $u$ and $v$ in $G$ would be a path that connects two distant locations, in which case the path might contain some highly important node at a level higher than $i$, due to which there would not be any shortcut between $u$ and $v$. More formally, we have the following lemma:
%
%\begin{lemma} \label{lmm:fc-sliding}
%Let $P$ be a shortest path in $G$, such that no ($3 {\times} 3$)-cell region in $R_i$ ($i \in [1, h]$) can cover all nodes in $P$ simultaneously. Then, $P$ must contain an arterial edge of some ($4 {\times} 4$)-cell region in $R_i$.
%\end{lemma}
%
%By Lemma~\ref{lmm:fc-sliding}, $u$ and $v$ must be covered in the same ($3 {\times} 3$)-cell region in $R_{i+1}$; otherwise, the shortest path between $u$ and $v$ must pass through a level-($i{+}1$) node, for which there cannot exist any shortcut between $u$ and $v$. This implies that $v$ must be in the ($5 {\times} 5$)-cell region in $R_{i+1}$ that is centered at $u$. By Lemma~\ref{lmm:fc-node-num}, this region contains $O(\a)$ level-$i$ nodes in $\G$. With a similar analysis, it can be shown that the region also covers $O(\a)$ nodes at any level higher than $i$. Therefore, the total number of edges adjacent to $u$ is $O(h \a)$.
%
%In summary, FC answers any distance query with two constrained Dijkstra search, each of which traverses $O(h\a)$ nodes and $O(h^2 \a^2)$ edges. As such, the time complexity of each traversal equals $O(h\a \log(h\a) + h^2 \a^2)$. Given that the arterial dimension $\a$ of the road network $G$ is a constant, the overall time complexity of FC is $O(h^2)$.
%
%\header {\bf Space Complexity.} Recall that the node hierarchy contains $h+1$ levels, each of which contains $O(n)$ nodes. In addition, each node in $\G$ has $O(\a)$ edges. Therefore, the space consumption of FC is $O(h n)$ when $\a$ is constant.
%
%\header {\bf Preprocessing Cost.} The pre-computation of FC consists of two steps: First, we identify the arterial edges in any ($4 {\times} 4$)-cell region in any grid $R_i$ ($i \in [1, h]$); After that, we decide the level of each node and we connect pairs of nodes with shortcuts. The identification of arterial edges requires computing the shortest paths in all ($4 {\times} 4$)-cell regions in all $R_i$, which incurs considerable overhead, especially when the granularity of the grid is low. Similarly, the construction of shortcuts is time consuming as it requires deriving a largest number of shortest paths (between nodes that are potentially far apart). Such significant precessing cost renders FC only applicable for small road networks. In Section~\ref{sec:ah}, we will address this issue and present a modified and scalable version of FC.
%
%\subsection{Correctness Proof} \label{sec:fc-correctness}
%
%Let $P = \langle v_1, v_2, \ldots, v_k\rangle$ be a shortest path in $G$. Let $P'$ be a path from $v_1$ to $v_k$ on the node hierarchy $\G$, such that FC reports $l(P')$ as the distance from $v_1$ to $v_k$. We will prove the correctness of FC's result by showing that $l(P') = l(P)$. In particular, we will show that both $l(P') \ge l(P)$ and $l(P') \le l(P)$ hold.
%
%\noindent {\bf Proving $\boldsymbol{l(P') \ge l(P)}$}. Recall that every shortcut on $\G$ corresponds to a path in $G$. Therefore, if we replace each shortcut in $P'$ with the corresponding path, we can transform $P'$ into a path $P''$, such that (i) $P''$ does not contain any shortcut, (ii) $P''$ connects $v_1$ to $v_k$, and (iii) $l(P') = l(P'')$. On the other , we have $l(P'') \ge l(P)$, since $P$ is the shortest path from $v_1$ to $v_k$ in $G$. Therefore, $l(P') = l(P'') \ge l(P)$.
%
%\header {\bf Proving $\boldsymbol{l(P') \le l(P)}$}. Assume for simplicity that $P$ contains a node $v_j$ ($j \in [1, k]$) whose level is higher than that of any other node on $P$. (Our analysis can be easily extended to the case when the highest-level node on $P$ is not unique.) Let $P_1 = \langle v_1, v_2, \ldots, v_j\rangle$ and $P_2 = \langle v_j, v_{j+1}, \ldots, v_k\rangle$. In the following, we will show that the node hierarchy $\G$ contains a path $P'_1$ from $v_1$ to $v_j$ that has the same length with $P_1$. Furthermore, we will prove that the sequence of nodes on $P'_1$ satisfies both the level and proximity constraints, i.e., $P'_1$ can be identified by FC with a traversal starting from $v_1$. In a similar manner, it can be shown that $\G$ contains a path $P'_2$ from $v_j$ to $v_k$, such that $l(P'_2) = l(P_2)$, and that $P'_2$ can be found by FC with a constrained Dijkstra search starting from $v_k$. This would lead to
%\begin{displaymath}
%l(P') \,\; \le \,\; l(P'_1) + l(P'_2) \,\; = \,\; l(P_1) + l(P_2) \,\; = \,\; l(P).
%\end{displaymath}
%
%Consider the path $P_1 = \langle v_1, v_2, \ldots, v_j\rangle$. Suppose that we remove from $P_1$ any node $v_i$ ($i \in [1, j]$) that has a smaller level than some node $v_a$ ($a < i$) preceding it. Let $S = \langle v'_1, v'_2, \ldots v'_b \rangle$ be the sequence of nodes remaining on $P_1$. We have $v'_1 = v_1$ (since no node precedes $v_1$), and $v'_b = v_j$ (since $v_j$ is the highest-level node in $P$). For instance, if $P_1$ contains six nodes $v_1$, $v_2$, $v_3$, $v_4$, $v_5$, $v_6$ at levels $1$, $0$, $2$, $2$, $1$, $3$, respectively, then $S = \langle v_1, v_3, v_4, v_6 \rangle$.
%
%By the way that $S$ is constructed, for any $v'_i$ ($i \in [1, b{-}1]$), the shortest path from $v'_i$ to $v'_{i+1}$ contains only nodes whose levels are smaller than those of $v'_i$ and $v'_{i+1}$. As such, the node hierarchy $\G$ would contain a shortcut cut from $v'_i$ to $v'_{i+1}$, and the length of the shortcut equals $dist(v'_i, v'_{i+1})$. This indicates that $\G$ contains a path $P'_1 = \langle v'_1, v'_2, \ldots v'_b \rangle$ that connects $v_1$ to $v_k$, such that $l(P'_1) = l(P)$. Furthermore, $P'_1$ satisfies the level constraint, since the level of $v'_i$ ($i \in [1, b{-}1]$) is no larger than that of $v'_{i+1}$.
%
%%Assume to the contrary that $P'_1$ does not satisfy the proximity constraint. Then, there should exist a node $v'_a$ ($a \in [1, b-1]$) on $P'_1$, such that (i) $v'_a$ is at level $i$ ($i \in [0, h]$), but (ii) no ($3\times3$)-cell region in $R_{i+1}$ covers $v'_1$ and $v'_{a+1}$ simultaneously. Then, $v'_1$ and $v'_a$ must locate in two distant cells in $R_{i+1}$. Intuitively, this indicates that the shortest path from $v'_1$ to $v'_a$ must run through a certain ($4 {\times} 4$)-cell region $B$ in $R_{i+1}$, in which case the path should contain an arterial edge of $B$. More formally, we have the following lemma:
%
%%\begin{lemma} \label{lmm:fc-sliding}
%%Let $P$ be a shortest path in $G$, such that no ($3 {\times} 3$)-cell region in $R_i$ ($i \in [1, h]$) can cover all nodes in $P$ simultaneously. Then, $P$ must contain an arterial edge of some ($4 {\times} 4$)-cell region in $R_i$.
%%\end{lemma}
%%\begin{proof}
%%The proofs for all lemmas and theorems in this paper can be found in our online and anonymous technical report \cite{???}.
%%\end{proof}
%
%Assume to the contrary that $P'_1$ does not satisfy the proximity constraint. Then, there should exist a node $v'_a$ ($a \in [1, b{-}1]$) on $P'_1$, such that (i) $v'_a$ is at level $i$ ($i \in [0, h]$), but (ii) no ($3\times3$)-cell region in $R_{i+1}$ covers $v'_1$ and $v'_{a+1}$ simultaneously. Then, by Lemma~\ref{lmm:fc-sliding}, the shortest path from $v'_1$ to $v'_a$ must contain an arterial edge $e$ of a ($4 {\times} 4$)-cell region in $R_{i+1}$, since none of the ($3 {\times} 3$)-cell regions in $R_{i+1}$ covers both $v'_1$ and $v'_a$. In that case, each endpoint of $e$ has a level at least $i+1$. In other words, on the shortest path from $v'_1$ to $v'_a$, there exists some node whose level is higher than that of $v'_a$ (recall that $v'_a$ is at level $i$). This contradicts the assumption that $v'_a$ has a level no lower than any node preceding it on $P_1$.
%
%In summary, the node hierarchy $\G$ contains a path $P'_1$ from $v_1$ to $v_j$, such that $P'_1$ has the same length with $P_1$ and satisfy both the level and proximity constraints. Therefore, FC can correctly identify the distance from $v_1$ to $v_j$ with a traversal starting from $v_1$. Similarly, we can show that FC can correctly compute the distance from $v_j$ to $v_k$ with a traversal starting from $v_k$. This proves the correctness of the query processing algorithm of FC.
%%
%%
%%
%%Among the neighbors of $v_{11}$, only $v_{10}$ is at a level is no lower than that of $v_{11}$. Therefore, the traversal from $v_{11}$ would first visit $v_{10}$, ignoring all other neighbors of $v_{11}$. After that, the traversal would not visit any other nodes in $\G$, since all unvisited neighbors of $v_{10}$ are at levels lower than $v_{10}$. Similarly, the traversal from $v_8$ would visit only thee nodes: $v_7$ and $v_{10}$ (since they are neighbors of $v_8$ with levels no lower than $v_8$), as well as $v_{11}$ (i.e., the only neighbor of $v_7$ or $v_{10}$ that satisfies the level constraint).
%%
%%
%%
%%In general, for any shortest path $P$ in the road network, there always exists a path $P'$ in $\G$, such that (i) $w(P) = w(P')$ and (ii) $P'$ bypasses unimportant nodes with shortcuts. To formalize this claim, we introduce that concepts of {\em level sequence} and {\em unimodal path}:
%%
%%\begin{property}[Level Sequence \& Unimodal Paths] \label{proper:fc-unimodal}
%%Let $P'$ be a path in $\G$, and $\langle v_1, v_2, \ldots, v_k\rangle$ be the sequence of nodes appearing on $P'$. The {\bf level sequence} of $P'$ is a sequence $\langle l_1, l_2, \ldots, l_k \rangle$, such that $l_i$ denotes the level of $v_i$ in $\G$. $P'$ is a {\bf unimodal path}, if its node sequence can be divided into two subsequence $\langle l_1, l_2, \ldots, l_j\rangle$ and $\langle l_{j}, l_{j+1}, \ldots, l_k\rangle$, such that $l_1 \le l_2 \le \ldots \le l_j$ and $l_j \ge l_{j+1} \ge \ldots \ge l_k$.
%%\end{property}
%
%
%%\begin{property}[Equivalent Paths in $\G$] \label{proper:fc-hrchy}
%%For any two nodes $s, t$ in $G$, if the distance from $s$ to $t$ equals $d$, then there exists a path $P$ in $\G$ that connects $s$ to $t$, such that (i) $P$'s weight $w(P)$ equals $d$, and (ii) the node sequence in
%%\end{property}
%
%
%
%%Given a road network $G$, AH preprocesses $G$ by transforming it into a hierarchy $\G$ with a constant number of levels (see Figure~\ref{fig:intro-hierarchy} for an example). Each node in $V$ is assigned to a unique level in $\G$ based on its {\em importance} for shortest path and distance queries, such that more important node would be at a higher level. To explain, let us consider a node $v_h$ on an inter-state highway and a node $v_l$ on a suburb street. Intuitively, $v_h$ would be more frequently accessed by shortest path or distance queries than $v_l$, since $v_h$ is on an important route that connects distant regions, while $v_l$ may only be used for local transportation in a small area. Therefore, $v_h$ is considered more important than $v_l$. Accordingly, AH would assign $v_h$ to a higher level in the hierarchy than $v_l$. Besides the nodes in $V$, the hierarchy $\G$ also contains all edges in $E$, as well as a set of {\em shortcuts} (i.e., auxiliary edges) that facilitate query processing (details to be clarified shortly).
%%
%%To answer a distance query from a node $s$ to a node $t$, AH performs two alternating traversals starting from $s$ and $t$, respectively. The traversal from s (t) only inspects the outgoing (incoming) edges of each node, i.e., it searches forward from s to t (backwards from t to s). Furthermore, both traversals mimic Dijkstra's algorithm in the sense that they both visit nodes in ascending order of their distances to the starting point. However, the two traversals differ from Dijkstra's algorithm in one crucial aspect: they never travel from a node at a higher level to a node at a lower level. For example, suppose that AH performs a traversal of the hierarchy in Figure~\ref{fig:intro-hierarchy}, starting from node $v_8$. Among the four edges adjacent to $v_8$, AH would follow $\langle v_8, v_{10}\rangle$ and $\langle v_8, v_7\rangle$, but it would omit $\langle v_8, v_3\rangle$ and $\langle v_8, v_5\rangle$, as both $v_3$ and $v_5$ are at a level lower than that of $v_8$. This traversal strategy significantly improves query efficiency, since it reduces the number of nodes that need to be inspected during query processing.
%%
%%Intuitively, the way that AH traverses the road network is similar to how we drive between two distant locations $s$ and $t$ in practice: We start from some small streets (connected to $s$), which take us to the main roads; following the main roads, we turn on to the highways, after which we keep traveling on the the highways for a period of time; upon approaching the destination, we exit the highways to get to some main roads, which then lead us to some small streets connected with $t$. In other words, in the first half of our journey, we always travel from less important locations in the road network (i.e., nodes on the small streets) to more important ones (i.e., nodes on the highways), while the second half of our journey visits a sequence of locations with decreasing importance. Similarly, when AH answers a distance query from $s$ to $t$, the traversal starting from $s$ keeps moving from less important nodes to more important ones, while the backward search from $t$ omits any edge that points from a node $u$ to a node $v$ with higher importance.
%%
%%The aforementioned query algorithm of AH, however, requires adding shortcuts to the node hierarchy. To illustrate, consider the hierarchy $\G$ in Figure~\ref{fig:intro-hierarchy} that contains two bi-directional shortcuts, $\langle v_9, v_{10} \rangle$ and $\langle v_{10}, v_{11} \rangle$. Assume that we remove $\langle v_{10}, v_{11} \rangle$ from $\G$, and we use the resulting hierarchy to answer a distance query from $v_{10}$ to $v_{11}$. In that case, neither the traversal from $v_{10}$ nor the one from $v_{11}$ would be able to proceed, since (i) both $v_{10}$ and $v_{11}$ are at level $2$ of the hierarchy, and (ii) all nodes connected with $v_{10}$ or $v_{11}$ are at lower levels. On the other hand, if we insert into $\G$ a shortcut $\langle v_{10}, v_{11} \rangle$ whose weight equals the travel time from $v_{10}$ to $v_{11}$, then the distance query from $v_{10}$ to $v_{11}$ can be easily answered by AH. In general, we will construct a set of shortcuts in $\G$ to ensure that AH can return exact answers for every shortest path and distance query. The weight of a shortcut $c$ usually equals the weight of shortest path $P$ in the road network $G$ that connects the two endpoints of $c$. Nevertheless, $\G$ may also contain some {\em redundant} shortcuts whose weight is larger than the weight of the corresponding path in $G$. Fortunately, such shortcuts do not affect the correctness of query results, as will be shown in Section~\ref{sec:ah-query}.
%%
%%In Sections \ref{sec:ah-query} and \ref{sec:ah-pre}, we will elaborate the query and preprocessing algorithms of AH. Before that, we first summarize the crucial properties of the hierarchy $\G$ constructed by AH, so as to form a basis for detailed discussions.
%%
%%\subsection{Properties of Hierarchy} \label{sec:ah-properties}
%%
%%Let $d_{max}$ and $d_{min}$ be the maximum and minimum $L_\infty$ distances between any two nodes in $G$, and $h = \left\lceil \log_2(d_{max}/d_{min}) \right\rceil - 1$. By Assumption~\ref{assu:def-log}, $h$ can be regarded as a constant. The hierarchy $\G$ constructed by AH has the following properties.
%%
%%\begin{property}[Number of Levels] \label{ppt:ae-height}
%%The number of levels in $\G$ equals $h$.
%%\end{property}
%%
%%Property~\ref{ppt:ae-height} is essential for the $O(1)$ time complexity of AH. This is because, when answering a distance query, AH may need to visit every level of $\G$, i.e., the time complexity of AH is at least at the order of the number of levels in $\G$.
%%
%%\begin{property}[Number of Shortcuts] \label{ppt:ae-scnum}
%%Each node $v$ in $\G$ is adjacent to at most $2 h \lambda$ shortcuts and edges, where $\lambda$ is the arterial dimension of $G$.
%%\end{property}
%%
%%Property~\ref{ppt:ae-scnum} indicates that the space complexity of AH is $O(\lambda n)$, which is linear to $n$ given that $\lambda$ is a constant. It also shows that, when searching in $\G$, we only need $O(1)$ time to inspect the edges and shortcuts connected to each node.
%%
%%The third property of $\G$ concerns the characteristics of the nodes in each level of $\G$. The formulation of the property is based on a few concepts as follows. Let $L_i$ ($i \in [0, h-1]$) be the set of nodes at the $i$-th level of $\G$ (assuming that $L_0$ corresponds to the lowest level). Let $V_i = L_i \cup L_{i+1} \cup \ldots \cup L_{h-1}$ be the set of all nodes at level $i$ and above. Let $R_i$ ($i > 0$) be a $2^{h - i + 1} \times 2^{h - i + 1}$ grid imposed on $G$, such that the sidelength of $R_i$ equals the maximum $L_\infty$ distance between any two nodes in $G$. In particular, $R_{h-1}$ contains $4 \times 4$ cells, and $R_{h-2}$ has a resolution of $8 \times 8$, etc.
%%
%%\begin{property}[Density of $V_i$] \label{ppt:ae-density}
%%Any $b \times b$ cell region in $R_i$ contains at most $(b+4)^2 \cdot 2\lambda^2$ nodes in $V_i$.
%%\end{property}
%%
%%Intuitively, Property~\ref{ppt:ae-density} can be understood as follows. Let $b = 2^{h-i+1}$. By Property~\ref{ppt:ae-density} and the fact that $R_i$ contains $2^{h-i+1} \times 2^{h-i+1}$ cells, we know that the total number of nodes in $V_i$ is upper-bounded by $(2^{h-i+1}+4)^2 \cdot 2\lambda^2$, which is a number that decreases exponentially with the increase of $i$. In other words, Property~\ref{ppt:ae-density} indicates that the density of the nodes in the higher levels of $\G$ are much smaller than that in the lower levels.
%%
%%The fourth property of $\G$ also concerns about $V_i$. For convenience, we say that a set $S$ of nodes {\em covers} a path $P$, if $P$ passes through at least one node in $S$. In addition, we say that two nodes are {\em far apart} with respect to the grid $R_i$, if they are contained in two different cells in $R_i$, such that each cell is outside of the other cell's outer shell. For example, in Figure~\ref{fig:intro-hierarchy}, the cell containing $v_8$ is beyond the outer shell of the cell that encloses $v_9$, and hence, $v_8$ and $v_9$ are far apart with respect to the grid.
%%
%%\begin{property}[Coverage of $V_i$] \label{ppt:ae-cover}
%%For any two nodes $s$ and $t$ that are far apart with respect to $R_{i}$, the shortest path from $s$ to $t$ is covered by $V_{i+1}$.
%%\end{property}
%%
%%By Property~\ref{ppt:ae-cover}, given $V_{h-1}$ (i.e., the nodes in the highest level of $\G$), we can cover all shortest paths between any pair of of nodes that are far apart with respect to the $4 \times 4$ grid $R_{h-1}$. On top of that, if we also want to cover the shortest path between any two nodes that are far apart with respect to the $8 \times 8$ grid $R_{h-2}$ (but not with respect to $R_{h-1}$), we can use $V_{h-2}$ instead, and so on. This implies that the nodes at a higher (lower) level of $\G$ are more (less) important for query processing, as they cover the shortest paths between locations that are faraway from (relatively close to) each other.
%%
%%The last property of AH characterizes the shortest paths in $\G$. For any path $P$ in $\G$, we define its {\em level sequence} as a sequence of numbers, such that the $i$-th number equals the level of the $i$-th node in $P$. For instance, in Figure~\ref{fig:intro-hierarchy}, there is a path that goes from $v_1$ to $v_{10}$ via $v_{11}$, and the levels of $v_1$, $v_{10}$, and $v_{11}$ are $0$, $2$, and $2$, respectively. The level sequence of the path is $(0, 2, 2)$.
%%
%%\begin{property}[Shortest Paths with Shortcuts] \label{ppt:ae-path}
%%For any two nodes $s$ and $t$, there exists a path $P$ from $s$ to $t$ consisting of edges and shortcuts in $\G$, such that (i) the weight of $P$ equals the weight of the shortest path from $s$ to $t$ in $G$, and (ii) the level sequences of the path is weakly unimodal.
%%\end{property}
%%
%%By Property~\ref{ppt:ae-path}, for any two nodes $s, t$ in $G$, we can find a path $P$ in $\G$ that connects $s$ to $t$, such that the weight of $P$ is identical to the travel time from $s$ to $t$ in $G$. Furthermore, the level sequence of $P$ is non-decreasing from the first number to a certain number, and is non-increasing from that number to the last number. For example, in Figure~\ref{fig:intro-hierarchy}, the path from $v_1$ to $v_{10}$ via $v_{11}$ has a level sequence $(0, 2, 2)$, which is unimodal. This property is crucial to the correctness of AH's query algorithm. To understand this, recall that when AH answers a distance query from $s$ to $t$, it employs two traversals starting from $s$ and $t$, respectively, such that each traversal only moves from a node $v$ to its neighbors whose levels are not lower than $v$. Since the shortest path $P$ from $s$ to $t$ in $\G$ has a unimodal level sequence, the traversal of AH from $s$ ($t$) can discover the part of $P$ with non-decreasing (non-increasing) levels. Once the two parts of $P$ are identified, AH can return an answer for the distance query by returning the weight of $P$, which, by Property~\ref{ppt:ae-path}, is equal to the weight of the real shortest path from $s$ to $t$ in $G$.
%%
%%\subsection{Query Algorithm for Distance Queries} \label{sec:ah-query}
%%
%%[To Do: Transition Paragraph]
%%
%%As mentioned in Section~\ref{sec:ah-overview}, given a distance query from a node $s$ to a node $t$, AH answers the query with a {\em forward traversal} from $s$ and a {\em backward traversal from $t$}, both of which adopt a modified version of Dijkstra's algorithm. In particular, the forward traversal maintains two values for each node $v$ that it encounters: (i) a variable $d_f(v)$ that records the minimum travel time from $s$ to $v$ via the nodes that have been visited in the traversal, and (ii) a boolean value that indicates whether $v$ is {\em settled}, i.e., whether $d_f(v)$ has been confirmed as the minimum travel time from $s$ to $v$ in the road network. In addition, the algorithm also maintains a priority queue $Q_f$ of nodes, such that $Q_f$ uses $d_f(v)$ as the sort key of a node $v$.
%%
%%The forward traversal starts by inserting $s$ into $Q_f$, with $d_f(s)$ set to $O$. After that, it iteratively extracts from $Q_f$ the node $u$ with the smallest key, and it inspects each {\em outgoing} edge (or shortcut) $e$ adjacent to $u$. If $e$ points to a node $v$ whose level is lower than $u$'s, $e$ is omitted. Otherwise, the algorithm checks whether $v$ has been inserted into $Q_f$ before. If $v$ is in $Q_f$, then we set $d(v) = \min\{d_f(v), d_f(u) + w(e)\}$; otherwise, we set $d_f(v) = d_f(u) + w(e)$ and push $v$ into $Q_f$. In other words, we check whether the path from $s$ to $v$ via $u$ has a weight smaller than $d(v)$, and we update $d(v)$ accordingly. After all outgoing edge from $u$ are examined, $u$ is marked as settled (since the $d_f(u)$ would remain unchanged in the subsequent part of the traversal).
%%
%%The backward traversal from $t$ works in the same way as the forward traversal, with a few exceptions as follows. First, for each node $v$, the traversal maintains a value $d_b(v)$ that equals the minimum travel time from $v$ to $t$ (instead of from $s$ to $v$) via the nodes that we been seen during the traversal. Second, the algorithm maintains a priority queue $Q_b$ of nodes that uses $d_b(v)$ as the sort key. Third, whenever the algorithm retrieves from $Q_b$ the node $u$ with the smallest key, it only examines the {\em incoming} edges (or shortcuts) of $u$ that start from nodes at a level above or equal to $u$'s.
%%
%%The above two traversals are performed in an alternating manner, i.e., the nodes in the two priority queues are extracted one at a time in a round-robin fashion. To determine when the traversals can be terminated, AH maintains a value $d_{min}$ that records the weight of the shortest path from $s$ to $t$ via the nodes encountered in either traversal. In addition, AH keeps track of the node $v_f$ ($v_b$) with the smallest key in $Q_f$ ($Q_b$). At the beginning of the algorithm, $d_{min} = +\infty$. After that, whenever AH retrieves a node $u$ from a priority queue, it examines whether $u$ has been marked as settled by both the forward and backward traversals. If $u$ is settled by both traversals, AH sets $d_{min} = \min\{d_{min}, d_f(u) + d_b(u)\}$, i.e., $d_{min}$ is compared to the minimum travel time from $s$ to $t$ via $u$, and it is updated if necessary. Given $d_{min}$, the forward traversal is terminated whenever (i) $Q_f$ becomes empty, or (ii) $d_{min} < d_f(v_f)$ (i.e, when the travel time from $s$ to any node in $Q_f$ is larger than $d_{min}$). Similarly, the backward traversal stops when $Q_b$ is emptied or $d_{min} < d_b(v_b)$. After both traversals are finished, AH returns $d_{min}$ as the answer for the distance query.
%%
%%\begin{theorem}\label{thrm:ae-dist}
%%AH answers any distance query in $G$ in $O\left(h^2 \lambda^3\right)$ time.
%%\end{theorem}
%%
%%%\begin{figure}[t]
%%%\begin{small}
%%%\begin{tabbing}
%%%\line(1,0){240}\\
%%%\hspace{4mm} \= \hspace{4mm} \= \hspace{4mm} \= \hspace{4mm} \=
%%%\hspace{4mm} \= \hspace{4mm} \= \hspace{4mm} \= \hspace{4mm} \=
%%%\hspace{4mm} \kill
%%%
%%%{\bf Algorithm {\em Distance}} ($\G$, $s$, $t$) \\[1mm]
%%%1.\> create a hash table $T_1$ that maps a node $v$ in $\G$ to a value $key_1(v)$ \\
%%%2.\> create a hash table $T_2$ that maps a node $v$ in $\G$ to a value $key_2(v)$ \\
%%%3.\> set $key_1(s) = 0$; set $key_1(v) = +\infty$ for any $v \ne s$ \\
%%%4.\> set $key_2(t) = 0$; set $key_2(v) = +\infty$ for any $v \ne t$ \\
%%%5.\> initialize three variables $lb_1 = 0$, $lb_2 = 0$, $ub = +\infty$ \\
%%%6.\> while
%%%4.\> insert entries $(s, 0)$ and $(t, 0)$ into $M_1$ and $M_2$, respectively \\
%%%5.\> while (($lb_1 + lb_2 < ub$) and ($M_1$ and $M_2$ are not both empty)) \\
%%%6.\>\> de-heap the top entry $(v, key_1(v))$ from $M_1$ \\
%%%7.\>\>\> for each outgoing edge (or shortcut) $e$ adjacent to $v$ \\
%%%8.\>\>\>\> if $e$ points to a node $u$ not at a lower level than $v$ \\
%%%9.\>\>\>\>\>
%%%
%%%2.\> for each client $c$ \\
%%%3.\>\> employ Dijkstra's algorithm to traverse the vertices in $G$ in \\
%%%\>\> ascending order of their distances to $c$; for each vertex $v$ traversed \\
%%%4.\>\>\> if $d(c, v) \le a(c)$ \\
%%%5.\>\>\>\> add an entry $\big\langle c, d(c, v)\big\rangle$ to the attraction set of $v$ \\
%%%6.\>\>\> else goto Line 2 \\
%%%7.\> return \\
%%%\line(1,0){240}
%%%\end{tabbing}
%%%\end{small}
%%%\codecapup
%%%\caption{The {\em Blossom} Algorithm} \label{code:ah-query}
%%%\codecapdown
%%%\end{figure}
%%
%%\subsection{Preprocessing Algorithm} \label{sec:ah-pre}
%%
%%
%%%Consider $V_0$ and the square grid $R_0$. Recall that $R_i$ contains $2^{h-i+1} \times 2^{h-i+1}$ cells, and $h = \left\lceil \log_2(d_{max}/d_{min}) \right\rceil - 1$, where $d_{max}$ and $d_{min}$ are the maximum and minimum $L_\infty$ distances between any two nodes in $G$. As a consequence, each cell $C$ in $R_0$ has a sidelength at most $l_min$, which implies that $C$ contains at most one node in $V_0$.
%%
%%Consider $V_i$ and the $2^{h-i+1} \times 2^{h-i+1}$ square grid $R_i$. We compute $V_{i+1}$ from $V_i$ in three steps.
%%
%%\header
%%
%%
%%%\subsection{Hierarchical road network}
%%%Our intuition of the proposed method is from the real-life travelling experience. Suppose a tourist intends to drive to a far away city, we then make the following two observations.
%%%
%%%The first observation is that when travelling far away, we usually take a number of highways, the shortcuts connecting major cities, by which we can directly and quickly reach a far-away place without bypassing the secondary city roads. If a certain highway is closed, then we have to choose a substitute path which usually costs longer travelling time than the highway. This real-life experience indicates that highways, and accordingly the entrances of highways, are deemed important in the road network because they frequently appear in the shortest paths, in particular in the long distance paths. Contrastingly, road intersections at the countryside are less important because they are commonly accessed only when the source is in the countryside. We use the term "importance" to measure the ability of a node to cover shortest paths. i.e. the more important a node, the more likely it is, compared to the other nodes, to appear in shortest paths. Using this observation we differentiate the nodes by their importances. We construct a hierarchical road network, in which the nodes are partitioned into \emph{levels}. At the bottom of the hierarchy is a set containing all the nodes in the road network. Then we pick out some important nodes to construct a higher level. After repeating this step, a series of levels, each of which contains a set of nodes, are constructed. Intuitively, the topmost level in our hierarchical road network contains the entrances of major highways.
%%%
%%%To select the important nodes, we introduce the second observation: no matter where the starting place or destination is, as long as they are far away from each other, there will be a few number of roads (including highways) through which we have to pass before we eventually leaving the starting place to the destination. We say those roads, and accordingly the road intersections, are important to the starting place because they cover all the shortest paths from the starting place to another place faraway. Also, in some sense, these roads are important in the road network because they usually cover a number of shortest paths. Consequently, we simulate this observation to pick out the important nodes of each level.
%%%
%%%%\begin{figure}
%%%%\centering
%%%%\includegraphics[width = 3in]{figures/original.eps}
%%%%\caption{Road Network} \label{fig:original}
%%%%\end{figure}
%%%%
%%%%\begin{figure}
%%%%\centering
%%%%\includegraphics[width = 3in]{figures/hierarchy.eps}
%%%%\caption{Hierarchical Road Network} \label{fig:hierarchical}
%%%%\end{figure}
%%%
%%%
%%%To exemplify the hierarchical road network and further the query method, let us consider a toy road network in Figure ~\ref{fig:original} that contains eleven nodes $v_1,v_2\ldots,v_{11}$ and twelve edges. The edges are bidirectional and we omit the arrows for simplicity. The weights of the thick edges equal 2, while the weights of the thin edges are 1. Figure \ref{fig:hierarchical} is its corresponding constructed hierarchy. We divide the nodes of the road network into three levels $\mathcal{L}=\{L_0,L_1,L_2\}$. $L_0=\{v_1,v_2,v_3,v_4,v_5,v_6\}$, $L_1=\{v_7,v_8,v_9\}$ and $L_2=\{v_{10},v_{11}\}$. For a more clear description of the construction process, in the following we use two different terms \emph{layer} and \emph{level}. At the very beginning all the nodes are at layer 0. In the first round we find the nodes $v_7,v_8,v_9,v_{10},v_{11}$ more important compared to the reset, so we upgrade $v_7,v_8,v_9,v_{10},v_{11}$ to layer 1. At the next round we find that $v_{10}$ and $v_{11}$ are the most important two nodes and upgrade them to layer 2. Hence finally the three levels are constructed. The \emph{level} of a node is determined by the topmost \emph{layer} it can be upgraded to. Note that although $v_{10}$ and $v_{11}$ are at level 2, they also appear in layer 1, as depicted in Figure ~\ref{fig:hierarchical}, which appropriately explains the construction process: important nodes are repeatedly selected from lower layers and upgraded to higher layers.
%%%
%%%After partitioning the nodes into levels, we add \emph{shortcuts} (additional edges) among them to accelerate the query answering. We add two types of shortcuts into the hierarchy. The first type is \emph{jumping shortcuts}. From a node $v$, the jumping shortcuts connect $v$ to its important nodes. Consistent with the second observation, the important nodes of $v$ is the ones that would appear on the shortest path from $v$ to another far-away node. As shown in Figure ~\ref{fig:hierarchical}, the dotted curve ($v_4$, $v_{11}$) is a jumping shortcut because the shortest path from $v_4$ to another far-away node, like $v_1$ or $v_2$, has to pass $v_{11}$. The other type is the \emph{query shortcuts}. The query shortcuts connect the important nodes of the road network, like the dashed line $(v_{10},v_{11})$\footnote{To keep figure ~\ref{fig:hierarchical} clear, not all the shortcuts are drawn out.}. The details of these two types of shortcuts will be fully discussed in the rest of the paper. Here we give a rough overview of how the query works. Suppose we are to find the shortest path between $v_2$ and $v_4$. We perform a bidirectional Dijkstra search. Starting from $v_2$ (resp. $v_4$) we can be quickly directed to $v_{10}$ and $v_{11}$ by jumping shortcuts, then we continue the bidirectional Dijkstra search from $v_{10}$ and $v_{11}$ with the aid of query shortcuts for acceleration. This process is quite similar to our real-life experience. Recall that driving example. The jumping shortcuts quickly direct us from the source (resp. target) to the important nodes, i.e. from a starting place (resp. destination) to an entrance of a highway without traversing other secondary roads to find the exact route. Meanwhile, the query shortcuts help us to quickly find the highways that connects the two entrances.
%%%
%%%
%%%\subsection{major steps of our method}\label{subsec:major_steps}
%%%In this subsection we describe the execution flow of the hierarchy construction.
%%%
%%%As outlined above, we construct the hierarchical road network in two steps. At first we compute the sets of nodes for each layer, from which we can partition the nodes into levels, and then we add shortcuts to the graph. The sets of nodes for each layer are computed in an iterative way. Let $V_i$ be the set of selected nodes at layer $i$.
%%%
%%%We start with $V_0$ = $V$. At the $i$th iteration, suppose we have already computed $V_0, V_1, \ldots V_{i-1}$, we pick out a set of nodes that constitute $V_i$. Corresponding to the running example, at the beginning we have $V_0=\{v_0,v_1,\ldots,v_{11}\}$, subsequently we compute $V_1=\{v_7,v_8,v_9,v_{10},v_{11}\}$, and finally we have $V_2=\{v_{10},v_{11}\}$.
%%%
%%%%\begin{figure}
%%%%\centering
%%%%\includegraphics[width = 2in]{figures/grid2.eps}
%%%%\caption{Transit Nodes} \label{fig:grid}
%%%%\end{figure}
%%%
%%%When computing $V_i$ we select the important nodes that cover the shortest paths between two far-away nodes. We define what is the so called "far-away" and formally describe the set $V_i$. We impose a grid on the graph. The grid is composed of a number of equally sized $g_i\times g_i$ squares, which are called \emph{cells}. For each cell $C$ in the grid, let us define the \emph{inner shell} (resp. \emph{outer shell}) as the boundary of a 3 $\times$ 3 (resp. 5 $\times$ 5) sized square centred at $C$. As illustrated in Figure \ref{fig:grid}, the dashed-line (resp. dotted-line) square illustrates the inner shell (resp. outer shell) of the cell that contains the node $v_9$. $s$ and $t$ are far away if $t$ is beyond the outer shell of the cell that contains $s$. In the example, $v_9$ and $v_7$ are far way.
%%%
%%%Let $C$ be a cell that contains a node $s$. We observe from the grid that, when $s$ and $t$ are far away, $SP(s,t)$ (resp. $SP(t,s)$) would usually pass through certain edges or nodes located inside the outer shell of $C$. We define a set of edges (resp. a set of nodes) to represent the important edges (resp. nodes) of $C$. Given a cell $C$, we say an edge $e\in E$ \emph{overlaps} its outer shell, if at least one of the two endpoints of $e$ locates inside the outer shell. For example, $(v_{11},v_{7})$ and $(v_8,v_{10})$ \footnote{In this paper we use $<u,v>$ to represent a directional edge from $u$ to $v$. For simplicity we use an undirected graph as an example, and accordingly we use $(u,v)$ to represent an edge with both directions.} overlap the outer shell while $(v_7,v_8)$ does not, although $(v_7,v_8)$ geometrically intersects the outer shell. We consider a sub-graph composed of the edges that overlap the outer shell of $C$ together with the endpoints of those edges, and we call that a sub-graph centred at $C$. A set $I_C$ of nodes is called \emph{inner border nodes} of $C$, if $v\in I_C$ is an endpoint of the edge that intersects the boundary of $C$, and $v$ geometrically locates inside $C$. A set $O_C$ of nodes called \emph{outer border nodes} of $C$ refers to those nodes in the sub-graph which geometrically locates beyond the outer shell of $C$. For example, the nodes $V-\{v_3,v_4\}$ and the edges $E-\{(v_7,v_8),(v_7,v_4),(v_8,v_3)\}$ constitute the sub-graph centred at $C$ in Figure ~\ref{fig:grid}. The inner border nodes of $C$ are $I_C=\{v_9\}$ and the outer border nodes $O_C=\{v_7,v_8\}$.
%%%
%%%A set $TE_C$ of edges in $E$ is called \emph{transit edges} of a cell $C$ if and only if it satisfies the following conditions. Firstly, each $e\in TE_C$ intersects the inner shell of $C$. Secondly, for any inner border node of $C$, say $s$, and any outer border node, say $t$, the shortest path from $s$ to $t$ (or reversely from $t$ to $s$) in the sub-graph centred at $C$, passes at least one edge in $TE_C$. A set $TN_C$ of nodes in $V$ is called \emph{transit nodes} of $C$ if $TN_C$ covers $TE_C$, i.e. for each $e\in TE_C$, $TN_C$ contains at least one of its two endpoints. Consider the example in Figure ~\ref{fig:grid}. $TE_C=\{(v_{11},v_7),(v_6,v_{10})\}$ is a set of transit edges respect to $C$, while the set $TN_C=\{v_{11},v_{10}\}$ a set of transit nodes.
%%%
%%%Although we only consider the shortest paths between the inner border nodes and outer border nodes in a locally sub-graph centred at $C$, we can prove that the transit edges (resp. transit nodes) computed in this way \emph{cover} all the shortest paths between the interior of $C$ and the exterior of its outer shell. Due to space limits, here we give an informal prove. Firstly, the shortest path originated from the node inside $C$ to the exterior of the outer shell must pass through an inner border node of $C$. Hence it is sufficient to consider only the inner border nodes instead of all the nodes inside $C$. The same reason applies to the outer border nodes. Secondly, although the local shortest path $p$ from an inner border node $s$ to an outer border node $t$ found in the sub-graph is not a real shortest path, say $SP(s,t)$ in $G$, we prove that $TE_C$ (resp. $TN_C$) still covers $SP(s,t)$. Let $u$ on $SP(s,t)$ be an outer border node that is closet to $s$. Such $u$ exists, because, as stated above, the shortest path to the exterior of outer shell must pass an outer border node. Since $u$ is the closet outer border node to $s$, which means that the parent of $u$ on $SP(s,t)$ is inside the outer shell, then $SP(s,u)$ could be found in the sub-graph. $TE_C$ (resp. $TN_C$) covers $SP(s,u)$, which in turn covers $SP(s,t)$.
%%%
%%%In the following we simply call the union of $TE_C$ (resp. $TN_C$) for all the cells $C$ in the grid transit edges (resp. transit nodes) without referring to a specific cell $C$.
%%%
%%%Hence our hierarchy construction works in the following steps.  At the very beginning we impose a grid on the graph. The side length of a cell $g_1$ equals the minimal horizontal or vertical distance between two nodes, i.e. the difference of the x-coordinates (resp. y-coordinates) of the nodes, so that each cell contains at most one node. %In the following we use $g_i$ to denote the side length of a cell imposed on the graph.
%%%Subsequently we compute the transit nodes for each cell, and the union of all the transit nodes constitute $V_1$. At the second round we enlarge the size of a cell to four times than before. i.e. merge four previous small cells into a larger one. As shown in Figure \ref{fig:grid}, the shadowed cell on the upper left corner is a larger sized cell with side length $g_2=g_1\times 2$, and is used to be imposed on $G$ at this round. After that we compute $V_2$ in a similar way. To accelerate the computation we reuse the outcomes when computing $V_1$. The details will be discussed in section \ref{subsec:levelconstruct}. Repeat the process above until the imposed grid is so large that none of the cells has outer shell. Suppose when the construction stops we have computed $V_0, V_1, \ldots, V_h$. We call $h$ the height of the road network hierarchy. In our experiments $h$ is usually less than twenty. After computing $\mathcal{V}=\{V_0, V_1, \ldots, V_h\}$ we partition the nodes into levels and add shortcuts among them.
%%%
%%%\subsection{highway dimension assumption}\label{subsec:assumption}
%%%Before elaborating on the details, here we discuss an interesting, and quite important feature of transit nodes which is a fundamental support that can guarantee the complexity bounds. On a road network, no matter how large a cell is, or no matter where the cell locates, the number of its transit nodes (as long as the cell has outer shell so that its transit nodes can be computed as described above) is always limited. We have verified this observation by experiments on the US road network datasets: arbitrarily draw a horizontal (resp. vertical) base line that cross the road network. The number of edges which intersect that base line is always less than $\sqrt{n}$ with $n$ the number of nodes in the road network. This result means that the number of edges that intersect the inner shell of a cell (the two vertical and two horizontal segments that constitute the boundary of a square) is linear to $\sqrt{n}$, which, by definition of transit nodes, further implies the number of transit nodes of a cell is linear to $\sqrt{n}$.  Practically that number is far smaller than $\sqrt{n}$ because the perimeter of the inner shell of a cell is far smaller than the length of the base line that cross the entire road network.\textbf{ [mention the experimental result here]} Hence it is sufficient and practical to assume that for any cell the number of its transit nodes is bounded to a parameter $b$, by which we can bound the space and time complexity of our method.


%
%% [inlined 5-P2P/p2p.tex]
\section{Algorithm for PPD Queries} \label{sec:p2p}

Suppose that we are given a PPD query from a source node $s$ to a destination node $t$. If both $s$ and $t$ are in the core graph $\GC$, then HoD first reads $\GC$ into the main memory, and then answers the query using an in-memory algorithm. On the other hand, if $s$ or $t$ is not in $\GC$, then the algorithm for processing the query is more sophisticated. In particular, in the most challenging case when neither $s$ nor $t$ is in $\GC$, HoD processes the query in three steps:
\begin{enumerate}%[itemsep = 1mm, topsep = -1mm, partopsep = 0mm]
\item {\bf Forward Search:} HoD traverses the forward graph $\GF$ starting from $s$, and for each node $v$ visited, computes the distance from $s$ to $v$ in $\GF$.

\item {\bf Backward Search:} This step is similar to the forward search, except that the graph being traversed is $\GB$, and the starting point of the traversal is $t$.

\item {\bf Core Search:} Given the distance information obtained in the first two steps, HoD loads the core graph $\GC$ into the main memory, and then derives the distance from $s$ to $t$.
\end{enumerate}
In what follows, we first present the details of each step for the case when $s$ and $t$ are not in $\GC$, and then we explain how the other cases can be handled. For convenience, we define an {\em index} value $\theta(v)$ for each node {\em not} in the core graph $\GC$, such that $\theta(v) = i$ only if the $i$-th adjacency list in $\GF$ belongs to $v$. By the way $\GF$ is constructed (see Section~\ref{sec:pre-terminate}), for any two nodes $u$ and $v$ with $\theta(u) < \theta(v)$, the rank of $u$ is no larger than the rank of $v$.

%use $v_i$ ($i \ge 1$) to denote the owner of the $i$-th adjacency list in $\GF$ (i.e., the adjacency list stores the outgoing edges of $v_i$). Note that for any node $v_j$ ($j > i$), its rank $r(v_j)$ must be no smaller than the rank $r(v_i)$ of $v_i$, as explained in Section~\ref{sec:pre-terminate}. In addition, in $\GR$, the adjacency list of $v_j$ is placed before that of $v_i$, since $\GR$ store adjacency list in the opposite order.

\subsection{Forward Search} \label{sec:p2p-forward}

The forward search of HoD maintains a hash table $\HF$ and a min-heap $\QF$. In particular, $\HF$ maps each node $v$ to a key $\kf(v)$, which equals the length of the shortest path from $s$ to $v$ that is seen so far. Initially, $\kf(s) = 0$, and $\kf(v) = +\infty$ for any node $v \ne s$. On the other hand, each entry in $\QF$ corresponds to a node $v$, and the key of the entry equals $\theta(v)$, i.e., the index of $v$. As will become evident shortly, $\QF$ ensures that the forward search visits nodes in ascending order of their indices, and hence, it scans the file structure of $\GF$ only once, without the need to backtrack to any disk block that it has visited before.

HoD starts the forward search by inspecting each edge $e = \langle s, v\rangle$ adjacent to $s$ in $\GF$, and then inserting $v$ into $\HF$ with a key $\kappa(v) = l(e)$. (Note that $\GF$ contains only outgoing edges.) In addition, HoD also inserts $v$ into $\QF$. After that, HoD iteratively removes the top entry $u$ in $\QF$, and processes $u$ as follows: for each edge $e = \langle u, v\rangle$ adjacent to $u$, if $\kf(v) = +\infty$ in the hash table $\HF$, HoD sets $\kf(v) = \kf(u) + l(e)$ and inserts $v$ into $\QF$; otherwise, if HoD sets $\kf(v) = \min\{\kf(v), \kf(u) + l(e)\}$. When $\QF$ becomes empty, HoD terminates the forward search, and retains the hash table $\HF$ for the third step of the algorithm, i.e., the core search.

\subsection{Backward Search} \label{sec:p2p-backward}

The backward search of HoD works in a manner very similar to the forward search, and it also maintains a hash table $\HB$ and a min-heap $\QB$. $\HB$ maps each node $v$ to a key $\kb(v)$, while $\QB$ sorts nodes according to their index values.

When the backward search starts, HoD sets $\kb(t) = 0$ and $\kb(v) = +\infty$ for any $v \ne t$, after which it inserts $t$ into $\QB$. Then, HoD iteratively extracts the top entry $u$ in $\QB$, and processes $u$ as follows. First, HoD examines the neighbors of $u$ in $\GB$ (recall that $\GB$ only records the incoming edges of each node). For each neighbor $v$, if $\kb(v) = +\infty$ in $\HB$, then HoD adds $v$ into $\QB$, setting $\kb(v)$ to the length of the edge $\langle u, v \rangle$; otherwise, if $\kb(v) > \kb(u) + l(\langle u, v \rangle)$, then HoD updates $\kb(v)$ to $\kb(u) + l(\langle u, v \rangle)$. The backward search completes when $\QB$ is empty, upon which HoD retrieves the hash table $\HB$ and moves on to the core search.

\subsection{Core Search} \label{sec:p2p-core}

The core search of HoD performs two traversals of the core graph $\GC$, and the first (resp.\ second) traversal is essentially a continuation of the forward (resp.\ backward) search. In particular, the core search uses the hash tables $\HF$ and $\HB$ created during the forward and backward searches, along with two new min-heaps $\QFF$ and $\QBB$. Each entry in $\QFF$ (resp.\ $\QBB$) corresponds to a node $v$, and the key of the entry equals the value $\kf(v)$ in $\HF$ (resp.\ $\kb(v)$ in $\HB$). This contrasts the forward and backward searches, which uses min-heaps where the key of each entry is the index value of a node. In addition, HoD also maintains a variable $\dm$ that records the smallest distance from $s$ to $t$ in the part of the augmented graph that has been traversed. Initially, $\dm = \min_v \big\{\kf(v) + \kb(v)\big\}$. (Observe that $\kf(v)$ is an upper-bound of the distance from $s$ to $v$ and $\kb(v)$ is an upper-bound of the distance from $v$ to $t$.)

As the first step of the core search, HoD inserts into $\QFF$ any node $u \in \GC$ with $\kf(u) \ne +\infty$ (i.e., $u$ is reached by the forward search previously). Similarly, it also adds into $\QBB$ any node $w \in \GC$ with $\kb(u) \ne +\infty$. After that, HoD removes the top entries from $\QFF$ and $\QBB$ in a round-robin fashion. Whenever it extracts an entry $u$ from $\QFF$, it inspects both $\kf(u)$ and $\kb(u)$. If $\kf(u) + \kb(u) < \dm$, then HoD sets $\dm = \kf(u)+\kb(u)$. In addition, for each outgoing edge $e = \langle u, v\rangle$ from $u$, HoD sets $\kf(v) = \min\{\kf(v), \kf(u) + l(e)\}$, and inserts $v$ into $\QFF$ (if $v$ has not been inserted before).

The core search terminates when $\dm$ is smaller than or equal to $\kf(u) + \kb(w)$, where $u$ and $w$ are the top entries in $\QFF$ and $\QBB$, respectively. The rationale is that, for any unseen entry $v$, the distance from $s$ to $v$ is at least $\kf(u)$, and the distance from $v$ to $t$ is at least $\kb(w)$. As a consequence, if $\dm \le \kf(u) + \kb(w)$, then the path from $s$ to $t$ via $v$ cannot be shorter than $\dm$. Upon termination of the core search, HoD returns $\dm$ as the result of the PPD query from $s$ to $t$.

\subsection{Degenerated Cases and Correctness} \label{sec:p2p-proof}

The above discussion assumes that neither $s$ or $t$ is in the core graph $\GC$. Nevertheless, HoD can easily handle the case when $s$ or $t$ is in $\GC$. In particular, if $s$ is in $\GC$, then HoD simply omits the forward search. Accordingly, the core search will use a hash table $\HF$ where $\kf(s) = 0$ and $\kf(v) = +\infty$ for every $v \ne s$. The case when $t$ is not in $\GC$ is handled similarly.

The following lemma shows the correctness of HoD's query algorithm for PPD queries:
\begin{lemma} \label{lmm:p2p-correct}
Given any source node $s$ and any destination node $t$, HoD correctly computes the distance from $s$ to $t$.
\end{lemma}
%\vspace{0mm}
\begin{proof}[Sketch]
Let $v_i$ ($i \ge 1$) be the node in $G_f$ with the $i$-th largest index value. We define $S_0$ as the set of nodes in $\GC$, and $S_i$ ($i \ge 1$) as a set that contains all nodes in $\GC$ as well as $v_1, v_2, \ldots, v_i$. We will use an induction on $i$ to show that, for any $S_i$, HoD can correctly compute the distance between any two nodes in $S_i$.

Consider the base case when $i = 0$. When both $s$ and $t$ are in $S_0$, HoD computes the distance from $s$ to $t$ based only on the nodes and edges in $\GC$. Recall that, in the preprocessing procedure of HoD, whenever we remove a node from the reduced graph, we always construct shortcuts between the remaining nodes, to ensure that the distance between any two remaining nodes is not affected. Therefore, the distance from $s$ to $t$ in $\GC$ should be equal to $\d(s, t)$, i.e., HoD should return correct result for the PPD query.

Assume that HoD can correctly compute the distance between any two nodes in $S_i$ ($i \ge 0$). Let $u$ be the node in $S_{i+1}$ but not $S_i$. We will first prove that HoD can correctly derive the distance from $u$ to any other node $v$ in $S_{i+1}$. First, since $v \in S_{i+1}$ and $v \ne u$, $v$ must be in $S_i$. Let $P$ be the shortest path from $u$ to $v$ in the original graph. Among the nodes in $S_i$ that also appear on $P$, let $w$ be the one that minimizes $\d(u, w)$. By the preprocessing algorithm of HoD, the distance from $u$ to $w$ should be preserved in the reduced graph, when all nodes not in $S_{i+1}$ are removed. This indicates that $u$ should have a shortcut to $w$, and the shortcut should have a length $\d(u, w)$. Therefore, when HoD performs the forward search from $u$, it can correctly identify the distance from $u$ to $w$. Meanwhile, by the induction hypothesis, HoD can also derive the correct distance from $w$ to $v$. Therefore, the distance from $u$ to $v$ can be correctly calculated by HoD. In a similar manner, we can also show that HoD can return the correct distance from any other node $v \in S_{i+1}$ to $u$, which completes the proof.
%Since $\kappa(v) \ne +\infty$, HoD must have traversed a path $P$ from $s$ to $v$, such that $l(P) = \kappa(v)$. This indicates that $\d(s, v) \le \kappa(v)$. To prove the lemma, it suffices to show that $\d(s, v) \ge \kappa(v)$ also holds.
%
%Assume to the contrary that there exists a path $P'$ in the original graph $G$, such that $l(P') < \kappa(v)$. Without loss of generality, assume that $P' = \langle v_1, v_2, \ldots, v_k\angle$, where $v_1 = s$ and $v_k = v$. %Suppose that we remove from $P'$ any node $v_i$ ($i \in [1, j]$) that has a smaller rank than some node $v_a$ ($a < i$) preceding it. Let $S = \langle v'_1, v'_2, \ldots v'_b \rangle$ be the sequence of nodes remaining on $P_1$. We have $v'_1 = v_1$, since no node precedes $v_1$. In addition, $v'_b = v_j$, since (i)
\end{proof}

\subsection{Space and Time Complexities} \label{sec:p2p-complex}

Overall, the PPD query algorithm of HoD requires $O(n + m_c)$ main memory, where $n$ is the number of nodes in $G$ and $m_c$ is the number of edges in the core graph $\GC$. This is because, during the core search, the algorithm reads $\GC$ into the memory, and it needs to maintain two hash tables and two min-heaps, each of which takes $O(n)$ space.

When processing a query, HoD inserts each node into each min-heap at most once, and it visits each edge in $\GC$, $\GF$, and $\GB$ no more than once. Therefore, the time complexity of the query algorithm is $O(n \log n + m')$, where $m'$ is the total number of edges in $\GC$, $\GF$, and $\GB$. In practice, the query time of HoD tends to be much smaller than this worst-case bound, since the forward and backward searches of HoD usually visit only a small portion of the augmented graph, as illustrated in Section~\ref{sec:overview-query}. The I/O costs of the query algorithm is $O((n + m')/B)$, because HoD scans $\GC$, $\GF$, and $\GB$ exactly once.

%\subsection{Reversed Search} \label{sec:sssp-reversed}
%
%Given the hash table $H$ obtained from the core search, the reversed search of HoD is performed as a sequential scan of the reversed graph $\GR$, where the nodes' adjacency list is stored in descending order their ranks. For each node $v$ visited, HoD checks each edge $e = \langle u, v\rangle$ adjacent to $v$ (recall that $\GR$ contains only incoming edges). If $\kappa(u) \ne +\infty$ and $\kappa(u) + l(e) < \kappa(v)$, then HoD sets $\kappa(v) = \kappa(u) + l(e)$. Once all nodes in $\GR$ are scanned, HoD terminates the reversed search and, for each node $v$, returns $\kappa(v)$ as the distance from $s$ to $v$.


%\section{Shortest Path Queries} \label{sec:path}
%
%Talk about how to store path information for reconstruction.

%\section{Extensions}\label{sec:exten}
%In this section we present two major optimizations for the basic algorithm presented in section ~\ref{sec:details}. The first optimization is called \emph{Node Extension} which more accurately selects transit nodes from the set of local shortest paths in $SelectTransitNode$. The second one is an alternative way of adding shortcuts which greatly improve the query efficiency.
%
%\subsection{Node extension}
%In $SelectTransitNode$, we select the transit nodes to $V_i$ by $V_{i-1}$ and $E_{i-1}$. We examine each local shortest path $p\in\mathcal{P}$ to find an edge $e$ that intersects the inner shell. If $e$ is a shortcut, we need to expand $e$ to find the real transit edge contracted by $e$. However, this expansion would probably upgrade a node which is not previously in $V_{i-1}$. As shown in Figure ~\ref{fig:realedge}, $(u',v')$ is the real-edge and is put into $TE_i$. In this case either $u'$ or $v'$ would be selected into $V_i$, but $u'$ or $v'$ is not necessarily in $V_{i-1}$. As a matter of fact, $(u,v)$ could also be considered a transit edge if all the local shortest paths in $\mathcal{P}$ which passes $(u',v')$ also passes $(u,v)$. When we say a path passes a shortcut we mean that the path passes the shortest path contracted by the shortcut. For example, as shown in Figure ~\ref{fig:level2}, when computing $V_2$, we already have four local shortest paths in $\mathcal{P}$: $p_1$ from $v_1$ to $v_3$, $p_2$ from $v_1$ to $v_4$, $p_3$ from $v_2$ to $v_3$ and $p_4$ from $v_2$ to $v_4$. As for $p_3$, the shortcut $(v_9,v_{10})$ contains a real-edge $(v_6,v_{10})$. Since in $\mathcal{P}$ all the shortest path which passes $(v_6,v_{10})$ also passes $(v_9,v_{10})$ (the sub-path $(v_9,v_6,v_{10})$), in this case it is reasonable to take $(v_9,v_{10})$ as an transit edge instead of $(v_6,v_{10})$.
%
%In addition, here we perform a node extension. We consider not only the real-edge $(u',v')$, but a sub-path extended from $u'$ and $v'$. Let $p=(s,\ldots,u,v,\ldots,t)\in\mathcal{P}$ be a local shortest path found by $FindLocalShortestPath$, then $p$ is a path composed of the nodes in $V_{i-1}$, except that neither of its source $s$ nor target $t$ is necessarily in $V_{i-1}$. Let $(u,v)$ be the edge on $p$ that intersects the inner shell of a cell that contains $s$. We further use $(u',v')$ to denote the real-edge contracted by $(u,v)$. If $(u,v)$ is an original edge in $G$, then let $(u',v')=(u,v)$. We extend the sub-path $(u',v')$ on $p$ by $\alpha$ hops ($\alpha$ is a pre-defined parameter), and we call it \emph{$\alpha$-hops extended path from $(u',v')$ of $p$}. For example, if we set $\alpha=2$, we extend the sub-path $(u',v')$ to $(u_1,u,u',v',v,v_1)$, where $u_1$ is the parent of $u$ (on $p$) and $v_1$ likewise. Note that $(u_1,u,v,v_1)$ is a sub-path on $p$, and $(u',v')$ is contracted  by $(u,v)$.  Let $\mathcal{P}_{(u',v')}\subseteq\mathcal{P}$ be the set of local shortest paths that has $(u',v')$ as its real-edge, and the sub-path $p'=(x,\ldots,y)$ be the longest common sub-path among all the $\alpha$-hops extended path from $(u',v')$ on each $p\in\mathcal{P}_{(u',v')}$. Then we put $(x,y)$ as an imaginary transit edge into $TE_i$ although $x$ and $y$ are not really connected by an original edge or shortcut. Further, either $x$ or $y$ would be selected from $TE_i$ as an transit node into $V_i$. In our experiments we select $\alpha=3$, namely, starting from the real-edge $(u',v')$, we trace three hops backward ($u,u_1,u_2$) and tree forward($v,v_1,v_2$). Because we found that the local shortest paths that share the same real-edge would usually diverge at the third hop.
%
%The key point of why Node-extension optimizes the method is that, we chose a longer common sub-path $(x,y)$ shared by a number of local shortest paths instead of the real-edge $(u',v')$ as the candidates of transit edges, which in turns makes the estimation of the node importance more accurate and plausible. This is consistent with our real-life experience. For example, if a path $p'$ is the common sub-path shared by a number of shortest-paths, when planing a route, we would usually pay more attention at the endpoints of $p'$ instead of the middle points. Observe that in this optimization we only substitute the real-edge with a longest common sub-path as a candidate transit edge for selecting transit node. This substitution would not vary the number of transit nodes since the real-edge and the common sub-path cover the same number of local shortest paths.
%
%\subsection{Query indexing}
%The second optimization is that, we use an alternative way of adding shortcuts instead of the one described in subsection ~\ref{label{sub:shortcuts}}. After we obtain a partition of nodes into levels $\mathcal{L}=\{L_0,L_1,\ldots,L_h\}$, we impose a total order on the nodes: The nodes in $L_{i+1}$ has a higher rank than those in $L_i$, and within the same level, the node covers more local shortest paths has a higher rank. Now we remove all the shortcuts added before and reconstruct the index in this way: we add two types of shortcuts: \emph{Jumping shortcuts} and \emph{Query shortcuts}. To accelerate the construction of jumping shortcuts, we also add a third type of shortcut \emph{level shortcuts}. Actually not all the level shortcuts are used in the query, hence, after constructing jumping shortcuts, we only preserve a subset of level shortcuts which are useful in the queries, and we call them query shortcuts.
%
%\subsubsection{Jumping shortcuts and query shortcuts}
%Before constructing the shortcuts, we slightly adjust the nodes layers. We define $\mathcal{R}=\{R_0,\ldots,R_h\}$, where $R_i=\bigcup\limits_{j=i}^h{R_j}$ for $0\leq i\leq h$, i.e. $R_i$ is the set of nodes that are at least in level $i$. We start from $i=1$ and $E_0=E$. At the $i$th iteration, firstly we impose a grid on $G$. The grid scale is the same as the one imposed before when computing the node levels, i.e. $g_i\times g_i$. For each $v\in R_i$, we use a similar algorithm to compute the layer-i neighbours of $v$ as described in \ref{sub:shortcuts}. For $v\in R_i$, the layer-i neighbours of $v$ refer to those nodes: if for any node $w$ such that $SP(v,w)$ passes at least one node in $R_i$, and $u$ on $SP(v,w)$ is a node in $R_i$ that is closet to $v$, then $u$ is the layer-i neighbour of $v$. The set of edges (or shortcuts) that connect each $v\in R_i$ to its layer-i neighbours is represented as $E_i$, and we call the edges in $E_i$ level shortcuts. Then at the $i+1$th iteration, the computation of the layer-(i+1) neighbours of a node $v\in R_{i+1}$ can reuse the edges in $E_i$ as shown in Figure ~\ref{fig:levelshortcuts}. The level shortcuts in $E_i$ guarantee the connectivity because $V_{i+1}\subseteq V_i$. %For a node $v\in V_{i+1}$, we perform a Dijkstra traversal from $v$, meanwhile we maintain a counter $ufb$ to count the unfinished branches as described in algorithm ~\ref{alg:addshortcuts}. Different from algorithm ~\ref{alg:addshortcuts}, during the traversal, we only visit the nodes in $V_i$. The level shortcuts in $E_i$ guarantee the connectivity because $V_{i+1}\subseteq V_i$.
%
%%\begin{figure}
%%\centering
%%\includegraphics[width = 3in]{figures/levelshortcuts.eps}
%%\caption{Level shortcuts} \label{fig:levelshortcuts}
%%\end{figure}
%
%\begin{figure}[t]
%%\begin{small}
%\begin{minipage}[t]{1.63 in}
%\centering
%\includegraphics[width = 1.67in]{figures/levelshortcuts.eps}
%\vspace{-3.49mm}
%\figcapup \caption{Level shortcuts $\quad$} \label{fig:levelshortcuts} \figcapdown
%\end{minipage}
%\begin{minipage}[t]{1.85 in}
%\centering
%%\hspace{3mm}
%\includegraphics[width = 1.67in]{figures/jumpingshortcut.eps}
%\figcapup \caption{Jumping shortcuts $\quad$} \label{fig:jumpingshortcut} \figcapdown
%\end{minipage}
%%\end{small}
%\end{figure}
%
%With the level shortcuts we can construct the jumping shortcuts in this way. Starting with $i=1$, at the $i$th iteration, we impose a grid on $G$ sized of $g_i\times g_i$. For each $s\in V-R_i$, (ATTENTION: or $s\in V-R_{i+1}$??) we add a shortcut to its transit node in $R_i$. Let $C$ be the cell that contains $s$. We consider a sub-graph that overlaps with the outer shell of $C$. Then we build a Dijkstra spanning tree rooted from $s$. During the Dijkstra traversal we only visit the nodes in $R_{i-1}$ connected by the level shortcuts. Then from a leaf node $t$, backtrack the path to its root $s$. Let $u$ on the path be the node in $R_i$ that is closet to $v$, then we add a shortcut from $s$ to $u$, which is called jumping shortcut. Such $u$ exists because the leaf node $t$ is outside the outer shell of $C$, and by definition the local shortest path from $s$ to $t$ would pass a transit node, which would be at least in $L_i$ or even higher level. In a word, a jumping shortcut connects a node to the transit node of its cell as illustrated in Figure ~\ref{fig:jumpshortcuts}.
%
%%\begin{figure}
%%\centering
%%\includegraphics[width = 3in]{figures/jumpingshortcut.eps}
%%\caption{Jumping shortcuts} \label{fig:jumpingshortcut}
%%\end{figure}
%
%The query shortcuts preserve a subset of the level shortcuts. For each $v\in R_i$, previously we add a level shortcut from $v$ to its layer-i neighbour $u$.  Now we preserve the shortcut only if $Rank_v<Rank_u$.
%
%Although we present the construction of level shortcuts and jumping shortcuts in two separate processes, these two constructions could be performed at the same time, i.e. at the $i$th iteration, we add level shortcuts that connect the nodes in $R_i$, and jumping shortcuts to the nodes in $R_i$ by the level shortcuts added in the $i-1$th iteration.
%
%\subsubsection{Query method}
%Now the query method works in the following way: for a given source $s$ and target $t$, we instantiate two Dijkstra traversal respectively from $s$ and $t$. During the traversal we use two priority queues $Q_f$ and $Q_b$ which respectively maintain the forward search from $s$ and the backward search from $t$. Initially $Q_f$ (resp. $Q_b$) contains $s$ (resp. $t$) only. Then we compute how far it is between $s$ and $t$. Based on the geometrical location of $s$ and $t$, we compute the highest level $l$, such that in the $g_l\times g_l$ scale cell, $t$ is beyond the outer shell of the cell that contains $s$, and in the $g_{l+1}\times g_{l+1}$ level scale, $t$ is inside the outer shell. If $l$ is greater than $l_s$, where $l_s$ is the level of $s$, then we dequeue $s$ and enqueue the transit nodes of $s$ into $Q_f$ by jumping shortcuts at level $l$. As illustrated in Figure ~\ref{fig:jumpingshortcut}, we first reach the transit nodes of $s$ (at level $l$) by jumping shortcuts, because the jumping shortcuts direct the search to the nodes that are highly probable on the right shortest path. After that we can perform a similar bidirectional Dijkstra search as described in section ~\ref{sub:query}. If the level of $s$ is greater than $l$, such that from $s$ there is no jumping shortcut at level $l$, then we can start the Dijkstra traversal from $s$. Since the search would only visit a node with higher rank, and $s$ itself is a node of higher rank, so the search would quickly converge to the node with highest rank in the path.
%
%\subsection{Directed graphs}\label{subsec:directedgraphs}



%% [inlined 4-SingleSource/singlesource_proof.tex]

\section{Algorithm for SSD Queries} \label{sec:sssp}

Given an SSD query from a node $s$ that is not in the core graph $\GC$, HoD processes the query in three steps:
\begin{enumerate}[itemsep = 0.5mm, topsep = 2mm]
\item {\bf Forward Search:} HoD traverses the forward graph $\GF$ starting from $s$, and for each node $v$ visited, computes the distance from $s$ to $v$ in $\GF$.

\item {\bf Core Search:} HoD reads the core graph $\GC$ into the main memory, and continues the forward search by following the outgoing edges of each node in $\GC$.

\item {\bf Backward Search:} HoD linearly scans the backward graph $\GB$ to derive the exact distance from $s$ to any node not in $\GC$.
\end{enumerate}
On the other hand, if $s$ is in $\GC$, then HoD would answer the query with a core search followed by a backward search, skipping the forward search. In what follows, we will present the details of three searches performed by HoD. For convenience, we define an {\em index} value $\theta(v)$ for each node {\em not} in the core graph $\GC$, such that $\theta(v) = i$ only if the $i$-th adjacency list in $\GF$ belongs to $v$. By the way $\GF$ is constructed (see Section~\ref{sec:pre-file}), for any two nodes $u$ and $v$ with $\theta(u) < \theta(v)$, the rank of $u$ is no larger than the rank of $v$.

\subsection{Forward Search} \label{sec:sssp-forward}

The forward search of HoD maintains a hash table $\HF$ and a min-heap $\QF$. In particular, $\HF$ maps each node $v$ to a key $\kf(v)$, which equals the length of the shortest path from $s$ to $v$ that is seen so far. Initially, $\kf(s) = 0$, and $\kf(v) = +\infty$ for any node $v \ne s$. On the other hand, each entry in $\QF$ corresponds to a node $v$, and the key of the entry equals $\theta(v)$, i.e., the index of $v$. As will become evident shortly, $\QF$ ensures that the forward search visits nodes in ascending order of their indices, and hence, it scans the file structure of $\GF$ only once, without the need to backtrack to any disk block that it has visited before.

HoD starts the forward search by inspecting each edge $e = \langle s, v\rangle$ adjacent to $s$ in $\GF$, and then inserting $v$ into $\HF$ with a key $\kappa(v) = l(e)$. (Note that $\GF$ contains only outgoing edges.) In addition, HoD also inserts $v$ into $\QF$. After that, HoD iteratively removes the top entry $u$ in $\QF$, and processes $u$ as follows: for each edge $e = \langle u, v\rangle$ adjacent to $u$, if $\kf(v) = +\infty$ in the hash table $\HF$, HoD sets $\kf(v) = \kf(u) + l(e)$ and inserts $v$ into $\QF$; otherwise, HoD sets $\kf(v) = \min\{\kf(v), \kf(u) + l(e)\}$. When $\QF$ becomes empty, HoD terminates the forward search, and retains the hash table $\HF$ for the second step of the algorithm, i.e., the core search.

%Given an SSD query from a source node $s$, HoD processes the query in three steps:
%\begin{enumerate}%[itemsep = 1mm, topsep = -1mm, partopsep = 0mm]
%\item {\bf Forward Search:} HoD traverses the forward graph $\GF$ and, for each node $v$ visited, computes the distance from $s$ to $v$ in $\GF$.
%
%\item {\bf Core Search:} HoD loads the core graph $\GC$ into the main memory, and computes the exact distance from $s$ to every node in $\GC$.
%
%\item {\bf Reversed Search:} HoD linearly scans the reverse graph $\GR$ to derive the exact distance from $s$ to any node not in $\GC$.
%\end{enumerate}
%In the following, we will present the details of each step.
%
%%For convenience, we use $v_i$ ($i \ge 1$) to denote the owner of the $i$-th adjacency list in $\GF$ (i.e., the adjacency list stores the outgoing edges of $v_i$). Note that for any node $v_j$ ($j > i$), its rank $r(v_j)$ must be no smaller than the rank $r(v_i)$ of $v_i$, as explained in Section~\ref{sec:pre-terminate}. In addition, in $\GR$, the adjacency list of $v_j$ is placed before that of $v_i$, since $\GR$ store adjacency list in the opposite order.
%
%\subsection{Forward Search} \label{sec:sssp-forward}
%
%The forward search of HoD maintains a hash table $H$ and a min-heap $Q^f$. In particular, $H$ maps each node $v$ to a key $\kappa(v)$, which equals the length of the shortest path from $s$ to $v$ that is seen so far. Initially, $\kappa(s) = 0$, and $\kappa(v) = +\infty$ for any node $v \ne s$. On the other hand, each entry in $Q^f$ corresponds to a node $v$, and the key of the entry equals {\em the position of $v$'s adjacency list in the forward graph $\GF$}. As will become evident shortly, $Q^f$ ensures that HoD can perform the forward search by scanning the file structure of $\GF$ once, without the need to backtrack to any disk block that it has visited before.
%
%In the beginning of the forward search, HoD inspects each edge $e = \langle s, v\rangle$ adjacent to $s$ in $\GF$, and then inserts $v$ into $H$ with a key $\kappa(v) = l(e)$. (Note that $\GF$ contains only outgoing edges.) In addition, HoD also inserts $v$ into $Q^f$.
%
%After that, HoD iteratively removes the top entry $u$ in $Q^f$, and processes $u$ as follows: for each edge $e = \langle u, v\rangle$ adjacent to $u$, if $\kappa(v) = +\infty$ in the hash table $H$, HoD sets $\kappa(v) = \kappa(u) + l(e)$ and inserts $v$ into $Q^f$; otherwise, if $\kappa(v) < +\infty$ and $\kappa(v) > \kappa(u) + l(e)$, HoD sets $\kappa(v) = \kappa(u) + l(e)$. When $Q^f$ becomes empty, the forward search terminates, and the hash table $H$ is handed over to the next step of the query algorithm, i.e., the core search.

\subsection{Core Search} \label{sec:sssp-core}

The core search of HoD is a continuation of the forward search, and it inherits the hash table $\HF$ created during the forward search. In addition to $\HF$, HoD creates a min-heap $\QC$, such that $\QC$ stores entries of the form $\langle v, \kf(v)\rangle$, where $v$ is a node and $\kf(v)$ is the key of $v$ in $\HF$. For any node $u$ with $\kf(u) \ne +\infty$ (i.e., $u$ is visited by the forward search), HoD inserts $u$ into $\QC$.

Given $\HF$ and $\QC$, HoD performs the core search in iterations. In each iteration, HoD extracts the top entry $v$ from $\QC$, and examines each outgoing edge $e$ of $v$. For every such edge, HoD inspects its ending point $w$, and sets $\kf(w) = \min\{\kf(w), \kf(v) + l(e)\}$. Then, HoD adds $w$ into $\QC$ if $w$ is currently not in $\QC$. This iterative procedure is repeated until $\QC$ becomes empty. After that, the hash table $\HF$ is sent to the last step (i.e., the backward search) for further processing.

\subsection{Backward Search} \label{sec:sssp-backward}

Given the hash table $\HF$ obtained from the core search, the reversed search of HoD is performed as a sequential scan of the backward graph $\GB$, which stores nodes in descending order of their index values. For each node $v$ visited during the sequential scan, HoD checks each edge $e = \langle u, v\rangle$ adjacent to $v$. (Note that $\GB$ contains only incoming edges). If $\kf(u) \ne +\infty$ and $\kf(u) + l(e) < \kf(v)$, then HoD sets $\kf(v) = \kf(u) + l(e)$. Once all nodes in $\GB$ are scanned, HoD terminates the backward search and, for each node $v$, returns $\kf(v)$ as the distance from $s$ to $v$.

One interesting fact about the backward search is that it does not require a heap to decide the order in which the nodes are visited. This leads to a much smaller CPU cost compared with Dijkstra's algorithm, as it avoids all expensive heap operations.

\subsection{Correctness and Complexities} \label{sec:sssp-proof}

Compared with Dijkstra's algorithm, the main difference of HoD's query algorithm is that it visits nodes in a pre-defined order based on their ranks. The correctness of this approach is ensured by the shortcuts constructed by the preprocessing algorithm of HoD. In particular, for any two nodes $s$ and $t$ in $G$, it can be proved that the augmented graph $\G$ always contains a path $P$ from $s$ to $t$, such that (i) $P$'s length equals the distance from $s$ to $t$ in $G$, and (ii) $P$ can be identified by HoD with a forward search from $s$, followed by a core search and a backward search. More formally, we have the following theorem.
\begin{theorem} \label{thrm:sssp-correct}
Given a source node $s$, the SSD query algorithm of HoD returns $\d(s, v)$ for each node $v \in G$.
\end{theorem}
The proof of Theorem~\ref{thrm:sssp-correct} in included in Appendix~\ref{sec:proofs}.
%\begin{proof}
%Let $w$ be any node in $G$. By Lemma~\ref{lmm:p2p-correct}, HoD can correctly answer a PPD query from $s$ to $w$, which indicates that (i) there exists a path $P$ from $s$ to $w$ in the augmented graph, and (ii) $P$ can be identified by the combination of a forward search from $s$, a backward search from $w$, and a core search in $\GC$. Let $P_f$, $P_b$, and $P_c$ be the sub-paths of $P$ that is traversed by the forward, backward, and core searches of HoD, respectively.
%
%When processing the SSD query from $s$, HoD should be able to identify $P_f$, since the SSD query algorithm of HoD performs a forward search from $s$ that is exactly the same as in the PPD query algorithm. Meanwhile, since the semi-core search of HoD continues the forward search in $\GC$ using Dijkstra's algorithm, the path $P_c$ should also be found by HoD. Therefore, we can prove the lemma by showing that the reversed search of HoD can identify $P_b$.
%
%Without loss of generality, assume that $P_b = \langle v_1, v_2, \ldots, v_k, w\rangle$. Since $P_b$ is a path traversed by the backward search of HoD, we have $r(v_1) > r(v_2) > \ldots > r(v_k) > r(w)$. This indicates that the index values of $v_1, v_2, \ldots, v_k, w$ are in descending order, i.e., $\theta(v_1) > \theta(v_2) > \ldots > \theta(v_k) > \theta(w)$. Now recall that the reversed search of HoD (for any SSD query) is performed by a sequential scan of $\GB$, where nodes are stored in descending order of their index values. Therefore, the reversed search visits $v_i$ before $v_{i+1}$ ($i \in [1, k-1]$), and traverses $v_k$ before $w$. In other words, the path $P_b = \langle v_1, v_2, \ldots, v_k, w\rangle$ can be identified by the reversed search, which completes the proof.
%\end{proof}

%As with the case for PPD queries, HoD's algorithm for SSD queries requires $O(n + m_c)$ main memory. This is due to the fact that (i) the forward, semi-core, and reversed searches of HoD all maintain a hash table that takes $O(n)$ space, and (ii) the semi-core search requires reading the core graph $\GC$ into the memory. The time complexity of the algorithm is $O(n \log n + m')$, where $m'$ is the total number of edges in $\GC$, $\GF$, and $\GB$. The reason is that, when processing an SSD query, HoD may need to scan $\GF$, $\GC$, and $\GB$ once, and it may need to put $O(n)$ nodes into a min-heap. On the other hand, the I/O costs of the query algorithm is $O((n + m')/B)$, since it requires at most one scan of $\GF$, $\GC$, and $\GB$.


The query algorithm of HoD requires $O(n + m_c)$ main memory, where $m_c$ is the size of the core graph. This is due to the fact that (i) the forward, core, and back searches of HoD all maintain a hash table that takes $O(n)$ space, and (ii) the core search requires reading the core graph $\GC$ into the memory. The time complexity of the algorithm is $O(n \log n + m')$, where $m'$ is the total number of edges in $\GC$, $\GF$, and $\GB$. The reason is that, when processing an SSD query, HoD may need to scan $\GF$, $\GC$, and $\GB$ once, and it may need to put $O(n)$ nodes into a min-heap. Finally, the I/O costs of the algorithm is $O((n + m')/B)$, since it requires at most one scan of $\GF$, $\GC$, and $\GB$.



%\subsection{Optimization}

%This section presents {\em Arterial Hierarchy (AH)}, a scalable indexing method built upon the FC approach introduced in Section~\ref{sec:fc}. Compared with FC, AH has the same space complexity, a similar time complexity for distance queries, but significantly smaller pre-computation cost. In addition, AH also supports shortest path queries in a worst-case efficient manner.
%
%%In the following, we will XXXX
%
%\subsection{Overview} \label{sec:ah-overview}
%
%The main structure of AH is a node hierarchy $\H$ that resembles FC's node hierarchy $H$. In particular, both $\H$ and $H$ have $h+1$ levels, and both of their $i$-th levels ($i \in [0, h]$) is associated with a square grid $R_i$ of $2^{h+2 - i} \times 2^{h+2 - i}$ cells. However, AH and FC differ substantially in the ways that they decide node levels, construct shortcuts, and process queries.
%
%\header
%{\bf Differences in Node Levels.} To compute the level of each node, FC first imposes each $R_i$ on the road network $G$, and then computes the arterial edges in each ($4 {\times} 4$)-cell region in $R_i$, after which FC decides the node levels based on the arterial edges. As discussed in Section~\ref{sec:fc-complexity}, the derivation of arterial edges could incur significant overheads, since each ($4 {\times} 4$)-cell region in a coarse grid may cover a large number of nodes and edges in $G$.
%
%In contrast, AH computes node levels with an incremental algorithm that substantially improves efficiency. Given $G$, it first imposes the grid $R_1$ on $G$. Based on $R_1$, it identifies a set of unimportant nodes in $G$, and it assigns them to level $0$ of the node hierarchy $\H$. Then, it removes a subset of the unimportant nodes from $G$, and it constructs shortcuts between the remaining nodes. This results in a {\em reduced} graph $G_1$ that is considerably smaller than $G$. After that, AH recursively reduces $G_1$ into smaller graphs $G_2, G_3, \ldots G_h$, during which it assigns nodes to higher levels of $\H$. For the reduction from $G_{i}$ to $G_{i+1}$, AH needs to impose the grid $R_{i+1}$ on $G_{i}$ and compute the shortest paths in each ($4 {\times} 4$)-cell region. However, this computation is inexpensive since (i) $G_i$ has a much smaller size than $G$, and hence, (ii) each ($4 {\times} 4$)-cell region in $R_{i+1}$ contain only a small number of nodes or edges in $G_i$.
%
%\header
%{\bf Differences in Shortcuts.} FC creates only necessary shortcuts to ensure correct results for distance queries under the level and proximity constraints. In contrast, the shortcuts constructed by AH are not only for processing distance queries under the level and proximity constraints, but also for computing the actual shortest path between any two given nodes. Specifically, every shortcut $\langle v_a, v_c\rangle$ in AH's node hierarchy $\H$ is associated with a node $v_b$, such that (i) both $\langle v_a, v_b \rangle$ and $\langle v_b, v_c \rangle$ are edges in $\H$, and (ii) the length of $\langle v_a, v_c\rangle$ equals the lengths of $\langle v_a, v_b \rangle$ and $\langle v_b, v_c \rangle$ combined. In other words, $\langle v_a, v_c\rangle$ can be transformed into a two-hop shortest path $\langle v_a, v_b, v_c \rangle$. As such, given any path $P'$ in $\H$, we can transform $P'$ into a path in $G$, by recursively replacing each shortcut in $P'$ with its corresponding two-hop path.
%
%For example, Figure~\ref{fig:ah-recon} illustrates a shortest path $\langle v_1, v_2, \ldots, v_6\rangle$ in $G$, as well as three shortcuts $\langle v_1, v_4\rangle$, $\langle v_2, v_4 \rangle$, and $\langle v_4, v_6\rangle$. The shortcut $\langle v_1, v_4\rangle$ is associated with the node $v_2$, since $v_1$ is directly connected with $v_2$ and $v_2$ is directly connected with $v_4$. Similarly, $\langle v_2, v_4 \rangle$ and $\langle v_4, v_6\rangle$ are associated with $v_3$ and $v_5$, respectively. Now suppose that, given a distance query from $v_1$ to $v_6$, AH identifies $P' = \langle v_1, v_4, v_6\rangle$ as the shortest path from $v_1$ to $v_6$ in $\H$. To derive the actual shortest path from $v_1$ to $v_6$ in $G$, AH first replaces the shortcut $\langle v_1, v_4 \rangle$ in $P'$ with a two-hop path $\langle v_1, v_2, v_4 \rangle$, since $\langle v_1, v_4 \rangle$ is associated with $v_2$. This transforms $P'$ into another path $\langle v_1, v_2, v_4, v_6\rangle$. After that, we can replace $\langle v_2, v_4 \rangle$ with $\langle v_2, v_3, v_4 \rangle$, and substitute $\langle v_4, v_6\rangle$ with $\langle v_4, v_5, v_6\rangle$. As such, we obtain the shortest path $\langle v_1, v_2, \ldots, v_6\rangle$ from $v_1$ to $v_6$ in $G$.
%
%In general, given any shortest path query from a node $s$ to another node $t$, AH first computes the shortest path $P'$ from $s$ to $t$ in $\H$, and then it converts $P'$ into the corresponding path $P$ in the original road network. The conversion from $P'$ to $P$ takes only $O(k)$ time, where $k$ is the number of edges in $P$. This is because (i) for any shortcut in $\H$, we can identify its corresponding two-hop path in $O(1)$ time, and (ii) converting $P'$ to $P$ requires only $O(k)$ replacements of shortcuts.
%
%\header
%{\bf Differences in Query Processing.} Besides the aforementioned shortcuts (for reconstructing shortest paths), the node hierarchy $\H$ of AH also contains some extra shortcuts that can be leveraged for higher query efficiency. As a consequence, AH's query processing algorithm is slightly more sophisticated than FC's, as will be elaborated in Section~\ref{sec:ah-query}.
%
%\subsection{Index Construction} \label{sec:ah-pre}
%
%Similar to the case of FC, AH constructs its node hierarchy $\H$ in two steps: it first assigns each node in $G$ to a level in $\H$, and then it constructs shortcuts in $\H$ for query processing.
%
%%for all nodes in the same level, it orders the nodes based on a measurement of their relative importance. Finally, it creates shortcuts between nodes based on both the node levels and orderings.
%%
%\header
%{\bf Deciding Node Levels.}
%Given the road network $G$, AH first imposes on $G$ the grid $R_1$, where each cell contains at most one node. After that, AH identifies all ($4 {\times} 4$)-cell regions in $R_1$ that cover at least one node in $G$. For each of the ($4 {\times} 4$)-cell region identified, AH computes the arterial edges of the region in $O(1)$ time, and it marks each endpoint of an arterial edge as a {\em level-$1$ core}. At the same time, AH assigns all unmarked nodes to level $0$ of the node hierarchy $\H$ since, intuitively, those nodes are less important than the level-$1$ cores. After that, if any ($4 {\times} 4$)-cell region $B$ contains a local shortest path $P$ from a level-$1$ core $u$ to another level-$1$ core $v$, such that $P$ only goes through unmarked nodes, then AH inserts into $G$ a shortcut $\langle u, v\rangle$ with the same length as $P$. We say that $\langle u, v\rangle$ is a shortcut generated from $B$, and we use $G_1$ to denote the modified version of $G$ with all shortcuts added.
%%We also define the ID of $\langle u, v\rangle$ as the smallest ID of the edges on $P$.
%Overall, the computation of level-$1$ cores and the construction of $G_1$ takes only $O(n)$ time, since the number of non-empty ($4 {\times} 4$)-cell regions in $R_1$ is $O(n)$, and each of those regions contains $O(1)$ nodes and edges in $G$ (recall that $G$ is degree-bounded).
%
%For example, given the road network $G$ and the grid $R_1$ in Figure~\ref{fig:def-grid}, AH would identify $5$ level-$1$ cores: $v_7, v_8, v_9, v_{10}, v_{11}$. After adding shortcuts, $G$ is transformed into the graph $G_1$ in Figure~\ref{fig:ah-G0}. There exists a shortcut $\langle v_9$, $v_{10}\rangle$ in $G_0$ since (i) both $v_9$ and $v_{10}$ are level-$1$ cores, and (ii) in the ($4 {\times} 4$)-cell region $B$ illustrated in Figure~\ref{fig:def-grid}, the local shortest path between $v_9$ and $v_{10}$ goes only through $v_6$, which is unmarked. In general, the shortcuts in $G_1$ ensures that the level-$1$ cores form a connected graph even if we remove all unmarked nodes from $G_0$.
%
%%\header
%%{\bf Deciding Node Levels.}
%Given $G_1$, AH selects a subset of the level-$1$ cores in $G_1$ that are deemed more important than the others. The selected nodes are marked as the {\em level-$2$ cores}, while the remaining level-$1$ cores are assigned to level $1$ of $\H$. After that, AH converts $G_1$ into a smaller graph $G_2$ that retains all level-$2$ cores. This procedure is applied in a recursive manner: In the $i$-th recursion ($i \in [1, h-1]$), AH picks level-($i{+}1$) cores from the level-$i$ cores in $G_{i}$, and then it assigns the un-picked ones to level $i$ of $\H$, after which it transforms $G_{i}$ into a smaller graph $G_{i+1}$.
%
%A natural question is: Given $G_{i}$, how should AH select the level-($i{+}1$) cores from the level-$i$ cores? One straightforward solution is to construct a subgraph of $G_{i}$ that contains only the level-$i$ cores, and then compute the arterial edges in the subgraph to identify the more important nodes as level-($i{+}1$) cores. For example, given $G_1$ in Figure~\ref{fig:ah-G0}, we can first construct a subgraph of $G_1$ that contains only the five level-$1$ cores (i.e., $v_7, v_8, v_9, v_{10}, v_{11}$) and the edges connecting them (i.e., the five edges on the loop $\langle v_7, v_8, v_{10}, v_9, v_{11}, v_7\rangle$). After that, we impose the grid $R_2$ on the subgraph, compute the arterial edges, and then mark the endpoints of the arterial edges as level-$2$ cores. While this approach is intuitive, we find that (i) the resulting node hierarchy does not guarantee query correctness under the level and proximity constraints, and (ii) without the level and proximity constraints, it is difficult to achieve favorable asymptotic bounds on query time. To address this issue, we adopt a more careful approach to choose the level-($i{+}1$) cores without affecting the applicability of the level and proximity constraints in query processing. Specifically, our approach utilizes the concept of {\em border nodes}:
%\begin{definition}[Border Nodes] \label{def:ah-border}
%Let $B$ be a ($4 {\times} 4$)-cell region in $R_i$ ($i \in [1, h]$). A node $v$ in $G$ is a {\bf border node} of $B$, if (i) $v$ is not contained in the $2 \times 2$ cells centered at $B$, and (ii) $v$ is an endpoint of an edge in $G$ that intersects the boundary of the east, west, south, or north strip of $B$.
%\end{definition}
%
%\begin{figure}[t]
%%\begin{small}
%\begin{minipage}[t]{1.68 in}
%\centering
%%\hspace{3mm}
%\includegraphics[width = 1.64in]{./figures/PreAH6.eps}
%\figcapup \caption{Graph $G_1$.} \figcapdown \label{fig:ah-G0}
%\end{minipage}
%%\hspace{0.02in}
%\begin{minipage}[t]{1.68 in}
%\centering
%\includegraphics[width = 1.64in]{./figures/PreAH4.eps}
%%\vspace{-3.52mm}
%\figcapup \caption{Reduce graph $G'_1$.} \figcapdown \label{fig:ah-G0reduced}
%\end{minipage}
%%\end{small}
%\vspace{-2mm}
%\end{figure}
%
%For example, in Figure~\ref{fig:def-grid}, $v_1, v_2, v_9, v_{11}$ are all border nodes of the ($4 {\times} 4$)-cell region $B$, since each of them is an endpoint of an edge that intersects the boundary of $B$'s west strip, and none of them is contained in the $2\times2$ cells centered at $B$. Similarly, $v_3, v_4, v_7, v_8$ are also border nodes of $B$. On the other hand, $v_6$ and $v_{10}$ are not border nodes of $B$, since they are not adjacent to any edge that intersect the boundaries of $B$'s four strips.
%
%To select level-($i{+}1$) cores from $G_{i}$, we first reduce $G_{i}$ by removing any node in $G_{i}$ that is not a level-$i$ core and not a border node of any ($4 {\times} 4$)-cell region in $R_{i+1}$. We use $G'_{i}$ to denote the reduced graph thus obtained. For instance, given $G_1$ in Figure~\ref{fig:ah-G0}, we would remove $v_5$ and $v_6$, since none of them is a level-$1$ core or a border node in $R_2$. Figure~\ref{fig:ah-G0reduced} illustrates the reduced graph $G'_1$, with the border nodes in $R_2$ highlighted.
%
%Given $G'_{i}$, we impose $R_{i+1}$ on $G'_{i}$ and inspect each ($4 {\times} 4$)-cell region in $R_{i+1}$ that contains at least one node. For each such region $B$, we compute every spanning path of $B$ (see Definition~\ref{def:def-ae}) that satisfies two conditions:
%\begin{enumerate} [itemsep = 0.7mm, topsep = 1.5mm]
%\item {\em Border Condition:} The two endpoints of the path are border nodes of $B$, while the other nodes are all level-$i$ cores.
%
%\item {\em Coverage Condition:} Every shortcut on the path is generated from a region completely covered by $B$.
%\end{enumerate}
%For example, in the ($4 {\times} 4$)-cell region in Figure~\ref{fig:ah-G0reduced}, the spanning path $\langle v_2, v_9, v_{10}, v_8, v_3 \rangle$ satisfies both the border and coverage conditions, since (i) both $v_2$ and $v_3$ are border nodes, and (ii) the only shortcut on the path, $\langle v_9, v_{10} \rangle$, is generated from the region $B$ in Figure~\ref{fig:def-grid}, which is contained in the current ($4 {\times} 4$)-cell region.
%
%
%
%%We have the following lemma:
%%
%%\begin{lemma} \label{lmm:ah-pseudo}
%%Let $B$ be a ($4 {\times} 4$)-cell region in $R_i$ ($i \in [1, h]$), $P'$ be a path in $G'_{i-1}$, and $P$ be the path in $G$ that corresponds to $P'$. If $P'$ is a spanning path of $B$ that satisfies both the border and coverage conditions, then $P$ is a spanning path of $B$ in $G$.
%%\end{lemma}
%
%%By Lemma~\ref{lmm:ah-pseudo}, any spanning path $P$ in $G'_{i-1}$ that fulfills the border and coverage conditions can be mapped to a spanning path $P'$ in the original road network $G$. The correspondence between $P$ and $P'$ motivates us to identify important edges on $P$ in the same manner as on $P'$. In particular,
%
%For each spanning path $P$ that fulfills the border and coverage conditions, if it connects the west and east (resp.\ north and south) strips of $B$, we identify the edge\footnote{If multiple edges or shortcuts in $P$ intersect the bisector, we choose an arbitrary one among them as the pseudo-arterial edge.}
%%This does not affect the correctness or asymptotic bounds of AH.}
%in $P$ that intersects $B$'s vertical bisector (resp.\ horizontal bisector) as a {\em pseudo-arterial edge} of $B$. Observe that each pseudo-arterial edge of $B$ corresponds to a path in $G$ that contains an arterial edge of $B$. Intuitively, this indicates the importance of pseudo-arterial edges in the reduced graph $G'_{i}$. Accordingly, we mark the two endpoints every pseudo-arterial edge as level-($i{+}1$) cores, and we assign all unmarked level-$i$ cores to the $i$-th level of the node hierarchy $\H$. After that, for any local shortest path in a ($4 {\times} 4$)-cell region $B'$, if (i) the two endpoints of the path are either level-($i{+}1$) nodes or border nodes of $B'$, and (ii) other than its endpoints, the path does not go through any level-($i{+}1$) core, then we insert into $G'_{i}$ a shortcut between $u$ and $v$ with the same length as the local shortest path. Once all such shortcuts are added, we define the resulting graph as $G_{i+1}$, and we use it to recursively compute higher-level nodes in $\H$, as previously explained.
%
%It remains to show that we can efficiently derive the pseudo-arterial edges and construct shortcuts in $G'_{i}$. Let $B$ be a ($4 {\times} 4$)-cell region in $R_i$, and $u$ be a border node of $B$. Suppose that we invoke Dijkstra's algorithm to start a traversal of $G'_{i}$ from $u$; for each node visited, we follow the the outgoing edges of the node, ignoring any edge that violates the border condition or coverage condition. Once the traversal terminates, we can obtain the spanning paths of $B$ starting from $u$, as well as the pseudo-arterial edges on those edges. Similarly, with a traversal from $u$ that follows only the incoming edges of each node, we can compute the desired spanning paths of $B$ ending at $u$, along with the pseudo-arterial edges therein. By repeating this process on all border nodes of $B$, we can derive the set of all pseudo-arterial edges in $B$. With the same traversal algorithm, we can construct all shortcuts in $B$ using two traversals from each border node of $B$.
%
%%The aforementioned spanning paths can be easily computed with a slightly modified version of Dijkstra's algorithm. In particular, for each border node $u$ of $B$, we first apply Dijkstra's algorithm to perform a traversal starting from $u$. For each node $v$ that we visit, if $v$ is a level-($i{-}1$) core node, then we inspect any edge $\langle v, w\rangle$ going out form $v$, and we put $w$ into the priority queue if $w$ is contained in $B$. If $v$ is not a level-($i{-}1$) node, then $v$ must be a border node, if which case we omit all edges adjacent to $v$. Once the traversal terminates, we obtain all spanning paths starting from $u$ that satisfy both the border and coverage conditions. Similarly, we can also apply Dijkstra's algorithm to derive all spanning paths ending at $u$.
%%
%%converts $G_{i-1}$ into a smaller graph $G_i$ and assigns nodes to the $i$-th level of $\H$.
%%
%%would repeat a graph {\em reduction} procedure on $G_0$ for $h-1$ times, such that the $i$-th
%%
%%would focus on the core nodes in $G_0$, and then identify the
%%
%% such that the $i$-th step ($i \in [1, h-1]$) converts $G_{i-1}$ into a smaller graph $G_i$ and assigns nodes to the $i$-th level of $\H$. Specifically, to reduce $G_{i-1}$ to $G_i$, AH first removes from $G_{i-1}$ any edge connecting two non-core nodes. After that, it imposes the grid $R_i$ on $G_{i-1}$, and it eliminates any non-core node that is not a {\em border node} of any ($4 {\times} 4$)-cell region in $R_i$. The concept of border node is defined as follows:
%
%\header
%{\bf Creation of Shortcuts.} After the level of each node is decided, AH adds shortcuts in the node hierarchy $\H$ to facilitate query processing. The construction of shortcuts requires as input a strict total order on the nodes in the same level of $\H$. We will elaborate our ordering approach in Section~\ref{sec:ah-optimize}, but in general, any strict total order can be used without affecting the space and time complexities of AH. For our discussion that follows, it suffices to know that less important nodes tend to precede more important nodes in our strict total order. For convenience, we define a rank for each node in $G$, such that a node $u$ ranks lower than another node $v$, if (i) $v$ is at a higher level than $u$, or (ii) $u$ and $v$ have the same level, but $u$ precedes $v$ in the strict total order.
%
%AH constructs shortcuts in $\H$ in an incremental manner similar to the algorithm for deciding node levels. In particular, it first inspects $G$, and it inserts into $\H$ a set of shortcuts that are concern level-$0$ nodes. After that, it reduces $G$ to a smaller graph $\G_1$. Subsequently, it recursively reduces $\G_{i}$ into another graph $\G_{i+1}$ ($i \in [1, h-1]$), during which its constructs shortcuts that concern nodes at the $i$-th level of $\H$. In the following, we will elaborate the reduction from $\G_{i}$ to $\G_{i+1}$ ($i \in [0, h-1]$), assuming $\G_0 = G$. For convenience, we define the {\em level} of every edge in $\G_0$ as $-1$.
%
%Given $\G_{i}$, AH first imposes the grid $R_{i+1}$ on $\G_{i}$. For each node $u \in \G_i$, AH inspects the ($5 {\times} 5$)-cell region $C$ centered at $u$, as well as the subgraph of $\G_i$ that consists of any level-($i{-}1$) edge overlapping with $C$. Then, AH computes two {\em shortest path trees (SPT)} of the subgraph, as defined in the following:
%\begin{definition}[Shortest Path Trees (SPT)] \label{def:ah-spt}
%Let $G$ be a graph, and $T$ be a directed spanning tree of $G$ rooted at a node $u$. $T$ is a {\bf forward SPT} of $G$, if $T$ contains the shortest path from $u$ to any node in $G$. On the other hand, if $T$ contains the shortest path from any node in $G$ to $u$, then $T$ is a {\bf backward SPT} of $G$.
%\end{definition}
%
%Let $T_f$ (resp.\ $T_b$) be the forward (resp.\ backward) SPT of the aforementioned subgraph that is rooted at $u$. Observe that $T_f$ (resp.\ $T_b$) can be computed by one traversal of the subgraph using Dijkstra's algorithm. Let $v$ be any node in $T_f$, such that $u$ ranks lower than $v$ but higher than any ancestor of $v$ in $T_f$. For any such $v$, AH generates a shortcut $\langle u, v\rangle$ with a length equal to the distance from $u$ to $v$ in $T_f$. In addition, AH associates $\langle u, v\rangle$ with a node $w$ on the path from $u$ to $v$, such that $w$ ranks higher than any node on the path except $u$ and $v$. This is to indicate that, when answering shortest path queries, AH can replace $\langle u, v\rangle$ with a two-hop path $\langle u, w, v\rangle$. (Our algorithm guarantees that such a two-hop path always exists.) Similarly, for any node $v$ in $T_b$, AH creates a shortcut $\langle v, u\rangle$, if $u$'s rank is lower than $v$'s but higher than those of $v$'s ancestors in $T_b$. Furthermore, the shortcut is associated with the node $w'$ that ranks the highest among $v$'s ancestors except $u$.
%%For any node $v$ in the subgraph, let $P^+$ denote the path between $u$ and $v$ in $T^+$. If $v$ is the only node on $P^+$ that ranks higher than $u$, then AH adds a shortcut $\langle u, v\rangle$ into the node hierarchy $\H$, such that the length of $\langle u, v\rangle$ equals that of $P^+$.
%%
%%In addition, AH associates $\langle u, v\rangle$ with the node $w$ that immediately follows $u$ on $P^+$, i.e., $w$ is the neighbor of $u$ that lies on $P^+$.
%We refer to the shortcuts constructed above as {\em level-$i$ edges}\footnote{If multiple shortcuts are constructed from one node $u$ to another node $v$, AH retains only the shortest one.}. Intuitively, these shortcuts connect each level-$i$ node $u$ directly to its nearby higher-rank nodes. By following these shortcuts during query processing, AH can avoid visiting less important nodes, which helps improve efficiency.
%
%Besides the level-$i$ edges, AH creates a shortcut from $u$ and to a node $v$ in $T_f$ if (i) $u$ and $v$ are both at level $i$ or above, and (ii) all ancestors of $v$ except $u$ are below level $i$. Likewise, if $T_b$ contains a node $v$ with a level at least $i$, such that $u$ is the only ancestor of $v$ at level $i$ or higher, then AH generates a shortcut from $v$ to $u$. These shortcuts are to ensure that $\G_i$ would remain connected when we reduce $\G_i$ by removing some nodes below level $i$, as will be clarified shortly.
%
%When $i > 0$ (i.e., $\G_{i}$ is produced from a previous reduction step), AH also generates some extra shortcuts (referred to as {\em elevating edges}), in a manner slightly different from the construction of level-$i$ edges. First, AH inspects each node $u$ in $\G_{i}$ at a level lower than $i$, and it examines the ($5 {\times} 5$)-cell region $C$ in $R_{i+1}$ that is centered at $u$. Then, AH constructs a subgraph of $\G_{i}$ that comprises of all level-$i$ edges covered by $C$, {\em as well as all edges that connect $u$ with any node at level $i$ or above}. After that, AH computes the subgraph's forward and backward SPTs rooted at $u$. Let $P_f$ be any path in the forward SPT that connects $u$ to a node outside of $C$, and let $v$ the first node on $P_f$ at level $i$ or above (our algorithm ensures that such $v$ always exits). For every node $w$ preceding $v$ on $P_f$ (including $u$), we construct a shortcut $\langle w, v\rangle$, and we associate it with the node that immediately follows $w$ on $P_f$. The shortcut $\langle u, v\rangle$ is constructed to enable AH to efficiently traverse from $u$ to the $i$-th level of $\H$, while the other shortcuts $\langle w, v\rangle$ are generated to ensure that AH can map $\langle u, v\rangle$ to its corresponding path in $G$. Similarly, if the backward SPT contains a path $P_b$ that links $u$ with a node located beyond $C$, we identify the node $v$ that is closest to $u$ on $P_b$ among those nodes at level $i$ or above. For every node $w$ that precedes $v$ on $P_b$, we create a shortcut $\langle w, u\rangle$, and we associate it with the node that immediately precedes $w$ on $P_b$.
%
%
%
%%If the forward SPT contains a path $P_f$ from $u$ to a node $v$, such that $v$ is the only node on $P_f$ that ranks higher than $u$, then AH creates a shortcut $\langle u, v\rangle$ and associates it with the node that immediately succeeds $u$ on $P_f$. Similarly, for any path $P_b$ in the backward SPT that connects $u$ with a node $v$, if $u$ ranks higher than all other nodes on $P_b$ except $v$, then AH adds a shortcut $\langle v, u\rangle$ and associates it with the node the immediately follows $v$ on $P_b$. Intuitively, such elevating edges can help reduce query overhead, since they enable AH to traverse from lower-level nodes directly to higher-level nodes.
%
%%Let $P_f$ be a path in the forward SPT that connects $u$ to a node outside of $C$ (This occurs when the node is adjacent to an edge intersecting $C$'s boundary), and let $v$ be the first node on $P_f$ with a rank higher than $u$. For any such path $P_f$, we construct a shortcut $\langle w, v\rangle$ from every node $w$ preceding $v$ on $P_f$, and we associate the shortcut with the node that immediately follows $w$ on $P_f$.
%
%
%Once all level-$i$ edges and elevating edges are created, they are inserted into both $\G_i$ and $H$. After that, AH reduces $\G_i$ by retaining only (i) the border nodes in $R_{i+2}$ and (ii) nodes at level $i$ or above. The resulting graph is defined as $\G_{i+1}$ and is fed into the next reduction step.
%
%
%%It starts by imposing the grid $R_1$ on the road network $G$. For each node $u \in G$, AH inspects the ($5 {\times} 5$)-cell region $C$ centered at $u$, as well as the subgraph of $G$ that consists of any edge overlapping with $C$. Then, AH computes two {\em shortest path trees (SPT)} of the subgraph rooted at $u$, as defined in the following:
%%\begin{definition}[Shortest Path Trees (SPT)] \label{def:ah-spt}
%%Let $G$ be a graph, and $T$ be a directed spanning tree of $G$ rooted at a node $u$. $T$ is a {\bf forward SPT} of $G$, if $T$ contains the shortest path from $u$ to any node in $G$. On the other hand, if $T$ contains the shortest path from any node in $G$ to $u$, then $T$ is a {\bf backward SPT} of $G$.
%%\end{definition}
%%
%%Let $T^+$ (resp.\ $T^-$) be the forward (resp.\ backward) SPT of the aforementioned subgraph, such that both $T^+$ and $T^-$ are rooted at $u$. Observe that $T^+$ (resp.\ $T^-$) can be computed by one traversal of the subgraph starting from $u$, using the Dijkstra's algorithm. For any node $v$ in the subgraph, let $P^+$ denote the path between $u$ and $v$ in $T^+$. If $v$ is the only node on $P^+$ that ranks higher than $u$, then AH adds a shortcut $\langle u, v\rangle$ into the node hierarchy $\H$, such that the length of $\langle u, v\rangle$ equals that of $P^+$. In addition, AH associates $\langle u, v\rangle$ with the node $w$ that immediately follows $u$ on $P^+$, i.e., $w$ is the neighbor of $u$ that lies on $P^+$. This is to indicate that, when answering shortest path queries, AH may replace $\langle u, v\rangle$ with a two-hop path $\langle u, w, v\rangle$. (Our algorithm guarantees that the shortcut $\langle w, v\rangle$ would also be added into $\H$, when AH examines the ($5 {\times} 5$)-cell region centered at $w$.) Similarly, for any node $v$ in $T^-$, AH creates a shortcut $\langle v, u\rangle$ if, on the path from $v$ to $u$ in $T^-$, $v$ is the only node that ranks higher than $u$.
%
%\subsection{Node Ranking and Selection}\label{sec:ah-optimize}
%
%As mentioned, the shortcut construction algorithm of AH assumes that there is strict total order on the nodes in the same level. While any strict total order can be used without affecting the asymptotic bounds of AH, we have found a heuristic ordering approach that leads to high practical performance. Specifically, for nodes in the $0$-th level of the node hierarchy $\H$, we adopt a random order; for nodes in the $i$-th level ($i \in [1, h]$) of $\H$, we derive their ordering based on information from the preprocessing procedure of AH. To explain, recall that AH decides node levels by recursively applying a reduction procedure on the road network $G$. During the $i$-th reduction step ($i \in [1, h-1]$), AH examines a graph that contains level-($i{-}1$) cores; It identifies a set $S_i$ of pseudo-arterial edges in the graph, marks the endpoints of those edges as level-$i$ cores, and then assigns all unmarked level-($i{-}1$) cores to the ($i{-}1$)-th level of $\H$.
%
%We observe that the edges in $S_i$ are connected to some extend, and there are some level-$i$ cores that serve as hub nodes for the connections (i.e., they are adjacent to a relatively large number of edges in $S_i$). Intuitively, those hub nodes are more important than the rest of the level-$i$ cores. Motivated by this, we order the the level-$i$ cores using a {\em vertex cover} approach: we inspect the graph formed by the edges in $S_i$, and we compute a vertex cover of the graph using the standard linear-time $2$-approximation algorithm \cite{IntroAlgo}. The output of the algorithm is a sequence $\xi$ of nodes in the graph, such that the $i$-th node $v$ in $\xi$ is adjacent to the largest number of edges that are disjoint from the first $i-1$ nodes. Based on $\xi$, we order the level-$i$ cores as follows: The $i$-th node in $\xi$ is given the $i$-th highest rank, and the level-$i$ cores not in $\xi$ are given the lowest ranks arbitrarily.
%
%Interestingly, we find that if a level-$i$ core (i) does not appear in $\xi$ and (ii) is adjacent to a tiny number of edges in $\xi$, then we can downgrade it to a level-($i{-}1$) core without affecting correctness or asymptotic performance of AH. Such downgrading reduces the number of high-level nodes in the node hierarchy $\H$, which in turn improves query efficiency, since the high levels of $\H$ are frequently traversed during query processing. In our implementation of AH, we downgrade a level-$i$ core to level $i{-}1$ if it is absent from $\xi$ and is adjacent to less than $10$ edges in $S_i$.
%
%\subsection{Query Processing}\label{sec:ah-query}
%
%The query processing algorithm of AH is similar to that of FC. In particular, for distance query from a node $s$ to another node $t$, AH also answers the query with two traversals of the node hierarchy $\H$ starting from $s$ and $t$, respectively. As with the case of FC, each traversal of AH is performed with a constrained version of Dijkstra's algorithm. However, the constraints adopted by AH are slightly different: It adopts the proximity constraint (see Section~\ref{sec:fc-query}) and a {\em rank constraint} as follows:
%\begin{itemize}
%\item {\em Rank Constraint:} When the traversal from $s$ (resp.\ $t$) visits a node $u$, it ignores any neighbor of $u$ that ranks lower than $v$.
%\end{itemize}
%Intuitively, the rank constraint is a refined version of the level constraint, in that it takes into account not only the levels of nodes but also the strict total order defined on each level of $\H$. It leads to higher query efficiency as it helps AH bypass a larger number of relative unimportant nodes during query processing.
%
%In addition, AH also exploits the elevating edges in $\H$ (see Section~\ref{sec:ah-pre}) to reduce query cost, based on the following lemma:
%\begin{lemma} \label{lmm:ah-longpath}
%For any two nodes $u,v \in G$, if no ($4 {\times} 4$)-cell region in $R_i$ ($i \in [1, h]$) can cover $u$ and $v$ simultaneously, then the shortest path from $u$ to $v$ must go through a node at level $i$ or above.
%\end{lemma}
%Let $R_j$ ($j \in [1, h]$) be the coarsest grid where no ($4 {\times} 4$)-cell region contains both $s$ and $t$. By Lemma~\ref{lmm:ah-longpath}, the shortest path from $s$ to $t$ should pass through at least one node with a level at least $j$. This indicates that AH's traversal from $s$ would meet its traversal from $t$ at level $j$ or above. Therefore, if $s$ is a border node in $R_j$ (in which case $s$ has elevating edges to level $j$), then when we start the traversal from $s$, we can follow the elevating edges of $s$ to move directly to level $j$, ignoring any edge that connects $s$ to a node at a level lower than $j$. After that, we can continue the traversal from level $j$ under the rank and proximity constraints.
%
%More generally, for any level-$i$ ($i < j$) node $v$ visited in the traversal from $s$, if $v$ is a border node in $R_j$, then we move along the elevating edges of $v$ to level $j$ or above, omitting any other edges of $v$. On the other hand, if $v$ is a border node in $R_{j'}$ instead of $R_j$  ($j' < j$), then we follow the elevating edges $v$ to level $j'$ or higher, i.e., we traverse as close to level $j$ as possible. Meanwhile, if $v$ does not have any elevating edges or $v$ is at a level at least $j$, then we traverse the edges of $v$ that satisfies the rank and proximity constraints. The same strategy is used when AH traverses from $t$. This traversal strategy reduces query time, since it enables AH to avoid visiting the low levels of node hierarchy $\H$.
%
%So far we have only discussed distance queries. For any shortest path query from $s$ to $t$, AH first treats it as a distance query and computes the shortest path $P'$ from $s$ to $t$ in $\H$. After that, AH recursively replaces each shortcut in $P'$ with its corresponding two-hop path, which converts $P'$ to the actual shortest path from $s$ to $t$ in $G$, as explained in Section~\ref{sec:ah-overview}.
%
%\subsection{Space and Time Complexities}\label{sec:ah-complexity}
%
%To establish the space and time complexities of AH, we first introduce a lemma that quantifies the densities of nodes in each level of AH's node hierarchy $\H$:
%\begin{lemma} \label{lmm:ah-node-num}
%Any ($\alpha {\times} \alpha$)-cell region in $R_i$ contains $O(\alpha^2 \a^2)$ nodes whose level in $\H$ are no lower than $i$, where $\a$ is the arterial dimension of $G$.
%\end{lemma}
%Our proof of Lemma~\ref{lmm:ah-node-num} is similar to that of Lemma~\ref{lmm:fc-node-num}. We first show that any ($\alpha {\times} \alpha$)-cell region $B$ in $R_i$ contains the endpoints of $O(\alpha^2 \a)$ arterial edges in $G$. After that, we prove that there exists a one-to-many mapping from the arterial edges in $G$ to the nodes in $B$ with levels at least $i$, such that each edge is mapped to $O(\lambda)$ nodes. Based on this, we show that any ($\alpha {\times} \alpha$)-cell region in $R_i$ contains only $O(\alpha^2 \a^2)$ level-$i$ nodes of $\H$.
%
%\vspace{-0.3mm} \header
%{\bf Space Overhead.} Given Lemma~\ref{lmm:ah-node-num}, we can prove that each node in $\H$ has $O(h\a)$ elevating edges and $O(\a)$ non-elevating edges. This is because, by the preprocessing algorithm of AH, there is an elevating edge from a node $u$ to a level-$i$ node $v$, only if $u$ and $v$ are contained in the same ($4 {\times} 4$)-cell region in $R_i$. By Lemma~\ref{lmm:ah-node-num}, there exist $O(\a)$ such level-$i$ nodes. Since $\H$ contains only $h+1$ levels, that the total number of elevating edges adjacent to $u$ is $O(h\a)$. Similarly, we can prove that, when excluding elevating edges, each node in $\H$ has $O(\a)$ edges. Therefore, the space overhead of AH is $O(hn\a)$, which reduces to $O(hn)$ when $\a$ is a constant.
%
%\begin{table*}
%\centering
%\tblcapup \vspace{-1mm}
%\begin{small}
%\caption{Asymptotic performance of the state of the art.} \tblcapdown \vspace{-1mm}
%\label{tbl:related-complexity}
%\setlength{\extrarowheight}{0.1mm}
%%\begin{tabular}{|@{\hspace{1mm}}c@{\hspace{1mm}}|@{\hspace{1mm}}c@{\hspace{1mm}}|@{\hspace{1mm}}c@{\hspace{1mm}}|@{\hspace{1mm}}c@{\hspace{1mm}}|@{\hspace{1mm}}c@{\hspace{1mm}}|@{\hspace{1mm}}c@{\hspace{1mm}}|}
%%\begin{tabular}{|@{\hspace{1mm}}c@{\hspace{1mm}}|@{\hspace{1mm}}c@{\hspace{1mm}}|c|c|c|c|} \hline
%\begin{tabular}{|c|c|c|c|c|c|} \hline
%{\bf Reference} & {\bf  Space } &
%{\bf  Preprocessing} & {\bf Distance Query } & {\bf Shortest Path Query} & {\bf Remark}\\ \hline
%\multirow{2}{*}{\cite{ms12}} &  \gape[t]{$O(n)$} & $O(n \log n)$ & $O(n^{0.5+\epsilon})$ & N/A & \multirow{2}{5.4cm}{$S$ is a user-defined parameter in the range of $[n\log\log n, n^2]$. $D = l_{max}/l_{min}$, where $l_{max}$ (resp.\ $l_{min}$) is the largest (resp.\ smallest) road network distance between two nodes. $k$ is the number of edges in the shortest path between the source and destination of the query. $h$ is as defined in Section~\ref{sec:fc-pre}.}\\ \cline{2-5}
% & $O(S)$ & \gape[t]{$\tilde{O}(S)$} & $\tilde{O}(n/\sqrt{S})$ & N/A & \\ \cline{1-5}
%\multirow{2}{*}{\cite{afg10}} &  \gape[t]{$O(n \log n \log D)$} & $O(n^2 \log n)$ & $O(\log^2 n \log^2 D)$ & $O(k + \log^2 n \log^2 D)$ & \\ \cline{2-5}
% & \gape[t]{$O(n \log n \log D)$} & $O(n^2 \log n)$ & $O(\log n \log D)$ & N/A & \\ \cline{1-5}
%\cite{ssa08} &  $O(n\sqrt{n})$ & $O(n^2 \log n)$ & \gape[t]{$O(k\log n)$} & $O(k\log n)$ & \\ \cline{1-5}
%this paper &  $O(hn)$ & $O(hn^2)$ & \gape[t]{$O(h \log h)$} & $O(k + h\log h)$ &  \\ \hline
%%\cite{ssa09}  & road networks & $O(n)$ & $O(n^2 \log n)$ & $O(\log n \log D)$ & N/A \\ \cline{3-6}
%\end{tabular}
%\end{small}
%\tbldown \vspace{-1mm}
%\end{table*}
%
%\vspace{-0.3mm} \header
%{\bf Query Time.} AH answers any distance query with two traversals of $\H$, starting from the source $s$ and destination $t$ of the query, respectively. Due to the proximity constraint, in the $i$-th level of $\H$ ($i \in [0, h]$), each traversal of AH only visits the nodes in a ($5 {\times} 5$)-cell region in $R_{i+1}$. By Lemma~\ref{lmm:ah-node-num}, such a region only contains $O(\a^2)$ nodes at level $i$. Hence, the totally number of nodes traversed by AH is $O(h \a^2)$. Furthermore, for each node $v$ visited during a traversal, AH either follows the elevating edges of $v$ to a certain level of $\H$, or moves along the non-elevating edges of $v$ that satisfy the rank and proximity constraints. As previously discussed, $v$ has $O(\a)$ elevating edges to each level of $\H$, and has $O(\a)$ non-elevating edges. Therefore, the total number of edges visited by AH is $O(h \a^3)$. Since each traversal is performed using Dijkstra's algorithm, its overall time complexity is $O(h\a^2 \log(h\a^2) + h\a^3)$. Consequently, when $\a$ is a constant, the time complexity of AH for a distance query is $O(h \log h)$.
%
%To answer a shortest path query from $s$ to $t$, AH first processes its corresponding distance query to retrieve the shortest path $P'$ from $s$ to $t$ in $\H$, and then its transforms $P'$ into the actual shortest path $P$ from $s$ to $t$ in $G$. The transformation from $P'$ to $P$ takes $O(k)$ time, where $k$ is the number of edges in $P$. Therefore, AH requires $O(k + h \log h)$ time to answer a shortest path query.
%
%\header
%{\bf Preprocessing Cost.} The preprocessing algorithm of AH consists of three steps: (i) assigning nodes to each level of $\H$, (ii) deriving the strict total order on nodes at the same level, and (iii) creating shortcuts in $\H$. When assigning nodes to the $i$-th level of $\H$ ($i \in [0, h-1]$), AH inspects each non-empty ($4 {\times} 4$)-cell region in $R_{i+1}$, and construct a subgraph that consists of the level-$i$ cores and border nodes in the region. For each node in the subgraph, AH needs to apply Dijkstra's algorithm to traverse the subgraph a constant number of times. Given Lemma~\ref{lmm:ah-node-num} and the fact that each level-$i$ core or border node (i) has $O(\lambda^2)$ edges and (ii) is contained in a constant number of ($4 {\times} 4$)-cell region in $R_{i+1}$, it can be proved that AH require $O(n^2 \lambda^2)$ time to assign nodes to level $i$ of $\H$. Meanwhile, AH takes only $O(n)$ time to derive the strict total order at level $i$ of $\H$, since the derivation is based on a linear time algorithm for vertex cover.
%
%To construct shortcuts at the $i$-th level of $\H$, AH needs to inspect a graph $G^*_i$ reduced from $G$. For each node $u$ in $G^*_i$, AH examines the a ($5 {\times} 5$)-cell region in $R_{i+1}$ that is centered at $u$, and it creates shortcuts for $u$ by traversing the nodes in the region at level $i$ or above. Based on Lemma~\ref{lmm:ah-node-num}, it can be proved that total cost of generating shortcuts for $u$ is $O(\a^2)$. As such, the time required to create shortcut at level $i$ of $\H$ is $O(n \a^2)$.
%
%Summing up the above analysis, we have the following theorem:
%\begin{theorem} \label{thrm:ah-complexity}
%Given a road network with a constant arterial dimension, AH takes $O(h n^2)$ time to construct an index that requires $O(h n)$ space. With the index, AH answers any distance query in $(h \log h)$ time and any shortest path query in $O(k + h \log h)$ time, where $k$ is the number of edges in the shortest path.
%\end{theorem}
%
%%\section{Algorithm details}\label{sec:details}
%%This section covers two parts: (i) How to partition the nodes into levels and (ii) how to add shortcuts to the graph and further how the query works. For clarity in the following we describe the method on undirected graphs, and we will further modify the method to directed graphs in section \ref{subsec:directedgraphs}.
%%
%%\subsection{Layer Construction}\label{subsec:levelconstruct}
%%In this section we focus on the details of computing transit nodes at a certain iteration. Suppose we have already computed $V_0, V_1, \ldots, V_{i-1}$, now we are at the \textbf{$i$th} iteration and intend to compute $V_i$.
%%
%%\begin{algorithm}
%%\caption{$Compute\_Level(i, V_{i-1}, G=(V, E\cup E_{i-1}))$}\label{alg:main}
%%\begin{algorithmic}[1]
%%\State Impose a grid on $G$, the size of a cell is $g_i\times g_i$
%%\State $\mathcal{P} \leftarrow \varnothing$
%%\ForAll {cell $C$}
%%\ForAll {$s\in I_C$}
%%\State $G'\leftarrow SubGraph(C,E_{i-1})$
%%\State $T_s\leftarrow BuildDijkstraTree(s,G')$
%%\ForAll {$t\in O_C$}
%%\State $p\leftarrow FindLocalShortestPath(s,t,T_s)$
%%\State $\mathcal{P} \leftarrow \mathcal{P} \bigcup \{ p\}$
%%\EndFor
%%\EndFor
%%\EndFor
%%\State $V_i \leftarrow SelectTransitNode(\mathcal{P},i)$
%%\State $E_i \leftarrow AddShortcuts(V_i)$
%%\State return $V_i, E_i$
%%\end{algorithmic}
%%\end{algorithm}
%%
%%Algorithm \ref{alg:main} presents the computation process. $Compute\_Level$ has three input parameters. $i$ and $V_{i-1}$ are as described before. $E_{i-1}$ is a set of edges (original edges and generated shortcuts) computed at the ${i-1}$th iteration. The edges in $E_{i-1}$ are used to accelerate the local shortest path search in the current iteration. At the first iteration $E_0=E$. $Compute\_Level$ works in the following way: for each cell $C$, find out the local shortest path from an inner border node to an outer border node in the sub-graph centred at $C$. After the path is found, we can select the transit nodes by examining every shortest path, from which the endpoint of an edge that intersects the inner shell of $C$ is picked out. This selection could be done at between line 11 and line 12. Here we do a lazy-selection for optimization purpose, i.e. we select all the transit nodes after $\mathcal{P}$, the set of all the local shortest paths are obtained (line 13). In addition, to avoid redundant computation, we use a Dijkstra spanning tree to preserve the shortest paths information. Firstly we construct a Dijkstra spanning tree from an inner border node $s$ (line 6), then the shortest paths from $s$ to the outer border nodes could be quickly computed from that tree(line 8). After all the local shortest paths are computed, $SelectTransitNode$ selects the transit nodes $V_i$ (line 13), and subsequently we generate some shortcuts (line 14) and compute a set of edges $E_i$ in order to accelerate the local shortest paths search in the next iteration. Finally $V_i$ and $E_i$ are returned which serve as the input for the ${i+1}$th iteration.
%%
%%In the following we first describe how to add shortcuts, and then $BuildDijkstraTree$, $FindLocalShortestPath$ and $SelectTransitNode$ because the latter procedures use the shortcuts $E_i$ generated in the last iteration.
%%
%%\input{4-Algorithm_Details/adding_shortcuts.tex}
%%
%%\subsubsection{Finding local shortest path} \label{subsec:findshortestpath}
%%Now we focus on the local shortest path search of a cell $C$ (line 3 to line 9 in Algorithm ~\ref{alg:main}). A the beginning we impose a grid on the graph, then we compute the local shortest paths for each cell. As for a cell $C$, with the help of $E_{i-1}$, we can only consider the sub-graph composed of the edges in $E_{i-1}$ which overlap the outer shell of $C$. The first type of edges in $E_{i-1}$ guarantees that the inner border nodes (resp. outer border nodes) are connected with the nodes in $V_{i-1}$, and the second type of edges in $E_{i-1}$ guarantees the connectivity among the nodes in $V_{i-1}$. Then we run $BuildDijkstraTree$ to build a Dijkstra spanning tree $T_s$ rooted from an inner border node by traversing the sub-graph. In $FindLocalShortestPath$ we match the outer border node $t$ and the leaves of $T_s$ to form a local shortest path. Consider the example in Figure ~\ref{fig:level2}, we are about to compute the shortest path from the inner border node $v_2$. The sub-graph considered is the one $G'=(V',E')$ where $V'=\{v_1,v_2,v_3,v_4,v_7,v_8,v_9,v_{10},v_{11}\}$ and $E'=\{(v_2,v_9),(v_9,v_{11}),(v_1,v_{11}),(v_9,v_{10}),(v_{10},v_8),(v_8,v_7),(v_{11},v_7),(v_7,v_4),(v_3,v_8)\}$. At first we build a Dijkstra tree as shown in Figure ~\ref{fig:DijkstraTree}. Then we match $v_3$ and the leaves of the tree $v_7$, $v_8$ to form the local shortest path. Finally the local shortest path $(v_2,v_9,v_{10},v_8,v_3)$ is found.
%%
%%%\begin{figure}
%%%\centering
%%%\includegraphics[width = 3in]{figures/DijkstraTree.eps}
%%%\caption{Dijjkstra tree} \label{fig:DijkstraTree}
%%%\end{figure}
%%
%%\input{4-Algorithm_Details/select_transit_nodes.tex}
%%
%%\subsection{Query indexing and queries}\label{subsec:query}
%%After the preprocessing above, we have tagged each node with its highest layer flag. Then we partition the nodes into $h$ levels, with $h$ the height of the hierarchy: for $0\le i\le h$, $L_i$ is the set of nodes tagged with an layer $i$ flag. Consistent with the preprocessing, the nodes in $L_i$ are more important than the ones in $L_{i-1}$. Then we get an total order on $V$: each node is associated with a rank. The nodes in $L_{i+1}$ have a higher rank than those in $L_i$. As for the ones in the same level, say $L_i$, we rank them according to the number of paths in $\mathcal{P}$ they cover, i.e. the more local shortest paths in $\mathcal{P}$ a node covers, the higher the rank. We use $Rank_v$ to represent the rank associated to $v$.
%%
%%%\begin{figure}
%%%\centering
%%%\includegraphics[width = 3in]{figures/CH.eps}
%%%\caption{Contraction Order} \label{fig:CH}
%%%\end{figure}
%%
%%Now we remove all the shortcuts added before and rebuild the query index in the following way: we in turn contract the nodes in increasing order of their ranks. When contracting a node, say $v$, we remove $v$ after adding some shortcuts, which are called "\emph{query shortcuts}". For all pairs $u,w$ such that $u$ and $w$ are both adjacent to $v$, and both of them have higher ranks than $v$, we add a shortcut $(u,w)$ weighted $weight((u,v))+weight((v,w))$ if and only if the shortest path connecting $u$ and $w$ is the $(u,v,w)$ one. As shown in Figure ~\ref{fig:CH}, the numbering of the nodes is consistent with their ranks. The algorithm in turns contracts $v_1$ to $v_{11}$. when contracting $v_5$, its adjacent nodes are $v_8$ and $v_9$. Since the shortest path between $v_8$ and $v_9$ is not the $(v_8,v_5,v_9)$ one, so we do not add a shortcut. When contracting $v_6$, since the shortest path between $v_9$ and $v_{10}$ is the $(v_9,v_6,v_{10})$ one, so we add a shortcut $(v_9,v_{10})$. Similarly we add $(v_{10},v_{11})$ when contracting $v_9$. At a first glance the shortcuts may increase the data volume of $G$ which further slow down the queries, however, as we will see later, actually the shortcuts efficiently speed up the queries.
%%
%%With the shortcuts we can execute the query using a bidirectional Dijkstra method with some modifications. For a given source $s$ and a target $t$, we instantiate two Dijkstra searches respectively starting from $s$ and $t$. When $v$ is extracted from the priority queue, the algorithm only visit those nodes $u$ such that $u$ is adjacent to $v$ (by original edges or added shortcuts) and $Rank_u>Rank_v$. As shown in Figure ~\ref{fig:CH}, when we query the shortest path from $v_1$ to $v_3$, we instantiate two instances of Dijkstra traversal respectively from $v_1$ and $v_3$. Then we can reach $v_{11}$ and $v_8$. The search stalks at $v_{11}$ since none of the adjacent nodes of $v_{11}$ has a higher rank. When expanding $v_8$ it relaxes the edge $(v_8,v_{11})$ and $(v_8,v_{10})$. Finally the shortest path is found.  With the $Rank_u>Rank_v$ condition, the query does not need to visit all the nodes adjacent to $v$, which can reduce a large amount of unnecessary visits. Because, as the search continues, the number of adjacent nodes with higher ranks is limited. The query method can be proven correct. Due to space limit we omit the proof here.
%%
%%The query indexing and query method is similar to the \emph{Contraction Hierarchy(CH)}, but our method can be proven bounded in both space complexity and query time complexity due to the construction of the node ordering. As a contrast, the node ordering in \emph{CH} is constructed by heuristic methods which result in unstable complexities in both space and time.
%%
%%\subsection{proof of complexity} \label{subsec:proof}
%%(1) space complexity
%%(2) time complexity


%\input{5-Optimizations/optimizations.tex}

%% [inlined 4.5-Extension/extension.tex]

\section{Extension for SSSP Queries} \label{sec:exten}

Given a source node $s$, an SSSP query differs from an SSD query only in that it asks for not only (i) the distance from $s$ to any other node $v$, but also (ii) the predecessor of $v$, i.e., the node that immediately precedes $v$ on the shortest path from $s$ to $v$. To extend HoD for SSSP queries, we associate each edge $\langle u, w\rangle$ in the augment graph $\G$ with a node $v$, such that $v$ immediately precedes $w$ on the shortest path from $u$ to $w$ in $G$. For example, given the augmented graph in Figure~\ref{fig:over-steps}e, we would associate the edge $\langle v_9, v_7\rangle$ with $v_6$, since (i) the shortest path from $v_9$ to $v_7$ in $G$ is $\langle v_9, v_6, v_7\rangle$, and (ii) $v_6$ immediately precedes $v_7$ in the path.

With the above extension, HoD processes any SSSP query from a node $s$ using the algorithm for SSD query with one modification: Whenever HoD traverses an edge $\langle u, w\rangle$ and finds that $\d(s, u) + l(\langle u, w \rangle) < \d(s, w)$, HoD would not only update $\d(s, w)$ but also record the node associated with $\langle u, w\rangle$. That is, for each node $w$ visited, HoD keeps track of the predecessor of $w$ in the shortest path from $s$ to $w$ that have been seen so far. As such, when the SSD query algorithm terminates, HoD can immediately return $\d(s, w)$ as well as the predecessor of $w$.

Finally, we clarify how the preprocessing algorithm of HoD can be extended to derive the node associated with each edge. First, for each edge $e$ in the original graph, HoD associates $e$ with the starting point of $e$. After that, whenever HoD generates a candidate edge $\langle u, w\rangle$ during the removal of a node $v$, HoD would associate $\langle u, w\rangle$ with the node that is associated with the edge $\langle v, w\rangle$. For example, in Figure~\ref{fig:over-steps}c, when HoD removes $v_7$ and creates a candidate edge $\langle v_9, v_{10}\rangle$, it associates the edge with $v_7$, which is the node associated with $\langle v_7, v_{10}\rangle$.




%% [inlined 7-Experiments/experiments.tex]
\section{Experiments}\label{sec:exp}

%This section experimentally compares HoD with two methods: (i) VC-Index \cite{ckcc12}, the state-of-the-art approach for distance queries on large graphs, and (ii) Dijkstra's algorithm, which requires that the dataset fits in the main memory. We implement HoD with C++, and obtain the C++ implementation of VC-Index from its authors. All of our experiments are conducted on a machine with a $2.4$GHz CPU and $48$GB memory, running $64$bit Ubuntu version $12.04$.

This section experimentally compares HoD with three methods: (i) VC-Index \cite{ckcc12}, the state-of-the-art approach for SSD and SSSP queries on disk-resident graphs; (ii) EM-BFS \cite{adm06}, an I/O efficient method for breadth first search; and (iii) EM-Dijk \cite{mo09}, an I/O efficient version of Dijkstra's algorithm. We include EM-BFS since, on unweighted graphs, any SSD query can be answered using breadth first search, which is generally more efficient than Dijkstra's algorithm. We obtain the C++ implementations of VC-Index, EM-BFS, and EM-Dijk from their inventors, and we implement HoD with C++. As the implementation of VC-Index only supports SSD queries, we will focus on SSD queries instead of SSSP queries. All of our experiments are conducted on a machine with a $2.4$GHz CPU and $32$GB memory.
%running $64$bit Ubuntu $12.04$

%\begin{figure*}[t]
%\begin{small}
%\begin{tabular}{cccc}
%\multicolumn{4}{c}{\hspace{-4mm} \includegraphics[height=2.2mm]{./figs/legend1.eps}} \vspace{1mm} \\
%\hspace{-2mm}\includegraphics[height=32mm]{./figs/usa_time.eps}
%&
%\hspace{-6mm}\includegraphics[height=32mm]{./figs/face_time.eps}
%&
%\hspace{-6mm}\includegraphics[height=32mm]{./figs/btc_time.eps}
%&
%\hspace{-6mm}\includegraphics[height=32mm]{./figs/web_time.eps} \\
%\hspace{0mm}(a) USRN & \hspace{0mm}(b) FB &
%\hspace{0mm}(c) u-BTC & \hspace{0mm} (d) u-UKWeb
%\end{tabular}
%\end{small}
%\figcapup \vspace{-2mm} \caption{Preprocessing time.} %\vspace{-1mm} %\figcapdown
%\label{fig:exp-undir-pre}
%\end{figure*}
%
%\begin{figure*}[t]
%\begin{small}
%\begin{tabular}{cccc}
%\multicolumn{4}{c}{\hspace{-4mm} \includegraphics[height=2.2mm]{./figs/legend1.eps}} \vspace{1mm} \\
%\hspace{0mm}\includegraphics[height=33mm]{./figs/usa_space.eps}
%&
%\hspace{-4mm}\includegraphics[height=33mm]{./figs/face_space.eps}
%&
%\hspace{-4mm}\includegraphics[height=33mm]{./figs/btc_space.eps}
%&
%\hspace{-4mm}\includegraphics[height=33mm]{./figs/web_space.eps} \\
%\hspace{0mm}(a) USRN & \hspace{0mm}(b) FB &
%\hspace{0mm}(c) u-BTC & \hspace{0mm} (d) u-UKWeb
%\end{tabular}
%\end{small}
%\figcapup \vspace{-2mm} \caption{Space consumption.} %\vspace{-1mm} %\figcapdown
%\label{fig:exp-undir-space}
%\end{figure*}

%\begin{table}[t]
%\vspace{-2mm}
%\centering
%\begin{small}
%\caption{Datasets.}
%\label{tbl:exp-data}
%\tblup
%\begin{tabular}{|c|c|c|@{\hspace{1mm}}c@{\hspace{1mm}}|@{\hspace{1mm}}c@{\hspace{1mm}}|c|}
%\hline
%{\bf Name} &   { $|\mathbf{V}|$} & { $|\mathbf{E}|$ } &  {\bf Weighted?} & {\bf Directed?} & {\bf Size} \\  \hline
%USRN &  $24.9$M & $28.9$M  & yes & no & 1.3GB \\ \hline
%FB & $58.8$M & $92.2$M  & no & no & 4.4GB \\ \hline
%u-BTC & $16.3$M & $95.7$M  & no & no & 4.6GB \\ \hline
%u-UKWeb & $6.9$M & $56.5$M  & yes & no & 2.7GB \\ \hline
%BTC & $16.3$M & $99.4$M  & no & yes & 2.3GB \\ \hline
%Meme & $53.6$M & $117.9$M  & no & yes & 2.8GB \\ \hline
%UKWeb & $104$M & $3708$M  & no & yes & 89GB \\ \hline
%%ClueWeb & $1592$M & $15418$M  & no \\ \hline %\hline \
%\end{tabular}
%%\tblcapup
%\tblcapdown
%\end{small}
%\end{table}

\begin{table}[t]
\vspace{-2mm}
\centering
\begin{small}
\caption{Datasets.}
\label{tbl:exp-data}
\tblup
\begin{tabular}{|c|c|c|@{\hspace{1mm}}c@{\hspace{1mm}}|@{\hspace{1mm}}c@{\hspace{1mm}}|c|}
\hline
{\bf Name} &   { $|\mathbf{V}|$} & { $|\mathbf{E}|$ } &  {\bf Weighted?} & {\bf Directed?} & {\bf Size} \\  \hline
USRN &  $24.9$M & $28.9$M  & yes & no & 0.86GB \\ \hline
FB & $58.8$M & $92.2$M  & no & no & 2.42GB \\ \hline
u-BTC & $16.3$M & $95.7$M  & no & no & 1.79GB \\ \hline
u-UKWeb & $6.9$M & $56.5$M  & yes & no & 1.02GB \\ \hline
BTC & $16.3$M & $99.4$M  & no & yes & 1.98GB \\ \hline
Meme & $53.6$M & $117.9$M  & no & yes & 3.17GB \\ \hline
UKWeb & $104$M & $3708$M  & no & yes & 61.8GB \\ \hline
%ClueWeb & $1592$M & $15418$M  & no \\ \hline %\hline \
\end{tabular}
%\tblcapup
\tblcapdown
\end{small}
\end{table}

\subsection{Datasets} \label{sec:exp-data}

We use five real graph datasets as follows: (i) USRN \cite{USRN}, which represents the road network in the US; (ii) FB \cite{gkbm10}, a subgraph of the Facebook friendship graph; (iii) BTC \cite{BTC}, a semantic graph; (iv) Meme \cite{lbkk09} and UKWeb \cite{UKWeb}, which are two web graphs. Among them, only USRN and FB are undirected. Since VC-Index, EM-BFS, and EM-Dijk are all designed for undirected graphs only, we are indeed of more undirected datasets for experiments. For this purpose, we transform BTC and UKWeb into undirected graphs, using the same approach as in previous work (see \cite{ckcc12} for details). After that, for each undirected (resp.\ directed) graph $G$, we compute its largest connected component (resp.\ weakly connected component) $C$, and we use $C$ for experiments. Table~\ref{tbl:exp-data} illustrates the details of the largest component obtained from each graph. In particular, u-BTC and u-UKWeb are obtained from the undirected versions of BTC and UKWeb, respectively.

%, and ClueWeb \cite{ClueWeb}



%On each graph, we generate SSD queries from $100$ randomly selected nodes, and we report the average running time of each method for the queries.

%In addition, for each SSD query from a node $s$, we inspect the results of the query to identify the nodes that can be reached from $s$. From those nodes, we randomly select a node $t$ and generate a PPD query from $s$ to $t$. This results in $100$ PPD queries, and we also report the average time required by each method to process the queries.






%\begin{table}
%	\centering
%		\begin{tabular}{l|l|l|1|1|1|1}
%			\hline \hline
%			Dataset &   $\abs{V}(\times 10^6)$ & \abs{E} (\times 10^6)& average-deg & disk size & directed & weighted \\  \hline\hline
%			USRN &  24 & 58 & 2.4 & 1.2 GB & no & yes\\ \hline
%			facebook & 59 & 92 & 3.1 & 3.1 GB & no & no \\ \hline
%			BTC & 16 & 99 & 6.1 & 1.6 GB & yes & no \\ \hline
%			u-BTC & 16 & 191 & 11.8 & 3.0 GB & no & no \\ \hline
%			UKWeb & 104 & 3708 & 35.6 & 62.0 GB & yes & no \\ \hline
%			u-UKWeb & 6 & 113 & 16.4 & 1.9 GB & no & yes \\ \hline
%			ClueWeb & 1592 & 15419 & 9.7 & 288.0 GB & yes & no \\ \hline
%			MemeTracker & 53 & 118 & 2.2 & 2.0 GB & yes & no \\
%			\hline %\hline
%		\end{tabular}
%	\caption{Datasets}
%	\label{tab:Dataset}
%\end{table}

\header
{\bf Remark.} The previous experimental study on VC-Index \cite{ckcc12} uses USRN, u-BTC, and u-UKWeb instead of their largest connected components (CC) for experiments. We do not follow this approach as it leads to less meaningful results. To explain, consider a massive undirected graph $G$ where each CC is small enough to fit in the main-memory. On such a graph, even if a disk-based method can efficiently answer SSD queries, it does not necessarily mean that it is more scalable than a main-memory algorithm. In particular, one can easily answer an SSD query from any node $s$ in $G$, by first reading into memory the CC that contains $s$, and then running a main-memory SSD algorithm on the CC. In general, given any graph $G$, one can use an I/O efficient algorithm \cite{mr99} to pre-compute the (weakly) connected components in $G$, and then handle SSD queries on each component separately.
%
% the scability of the method -- we could have used a main-memory
%
%suppose that we issue an SSD query from a node in a tiny connected component $C'$ in a massive graph. Even if a disk-based method can efficiently answer such a query, it is not significant, since one may process the same query equally (or more) efficient by first reading $C'$ into the memory and then performing the SSD query using Dijkstra's algorithm.

\subsection{Results on Undirected Graphs} \label{sec:exp-undirected}

In the first sets of experiments, we evaluate the performance of each method on four undirected graphs: USRN, FB, u-BTC, and u-UKWeb. For HoD, EM-BFS, and EM-Dijk, we limit the amount of memory available to them to $1$GB, which is smaller than the sizes of all datasets. For VC-Index, we test it with $2$GB memory as it cannot handle any of our datasets under a smaller memory size.

Table~\ref{tbl:exp-pre} shows the preprocessing time of HoD and VC-Index on the four graphs. (EM-BFS and EM-Dijk are omitted as they do not require any pre-computation.) In all cases, HoD incurs a significantly smaller overhead than VC-Index does. In particular, on FB, the preprocessing time of HoD is more than ten times smaller than that of VC-Index. Table~\ref{tbl:exp-space} compares the space consumptions of HoD and VC-Index. Except for the case of u-BTC, the space required by VC-Index is consistently larger than that by HoD.

\begin{table}[t]
\centering
\vspace{-2mm}
\begin{small}
%\tblcapup
\caption{Preprocessing time (in minutes).}
\label{tbl:exp-pre}
\tblup
\begin{tabular}{|c|c|c|c|c|}
\hline
{\bf Method} &   {\bf USRN} & {\bf FB} &  {\bf u-BTC} & {\bf u-UKWeb} \\  \hline
HoD & $\;\,$4.0 & $\;\,$22.4 & 34.4 & 105.5 \\ \hline
VC-Index & 20.3 & 281.8 & 78.1 & 768.2 \\ \hline
\end{tabular}
\tblcapdown
\end{small}
\end{table}

\begin{table}[t]
\centering
\begin{small}
%\tblcapup
\caption{Space Consumption (in GB).}
\label{tbl:exp-space}
\tblup
\begin{tabular}{|c|c|c|c|c|}
\hline
{\bf Method} &   {\bf USRN} & {\bf FB} &  {\bf u-BTC} & {\bf u-UKWeb} \\  \hline
HoD & 2.5 & 5.1 & 3.8 & $\;\,$3.3 \\ \hline
VC-Index & 4.3 & 8.3 & 1.2 & 14.0 \\ \hline
\end{tabular}
\tblcapdown
\end{small}
\end{table}

\begin{table}[!t]
\centering
\begin{small}
%\tblcapup
\caption{Average running time for SSD queries (in seconds).}
\label{tbl:exp-ssd}
\tblup
\begin{tabular}{|c|c|c|c|c|}
\hline
{\bf Method} &   {\bf USRN} & {\bf FB} &  {\bf u-BTC} & {\bf u-UKWeb} \\  \hline
HoD & $\;\,\;\,$1.8 & $\;\,\;\,\;\,$3.2 & $\;\,\;\,$1.6 & $\;\,\;\,$1.4 \\ \hline
VC-Index & $\;\,$27.2 & $\;\,\;\,$94.9 & $\;\,$10.1 & $\;\,$70.0 \\ \hline
EM-BFS & $-$ & $\;\,$465.3 & 395.4 & $-$ \\ \hline
EM-Dijk & 430.7 & 1597.4 & 844.1 & 553.8 \\ \hline
\end{tabular}
\tblcapdown
\end{small}
\end{table}

To evaluate the query efficiency of each method, we generate $100$ SSD queries for each dataset, such that the source node of each query is randomly selected. Table~\ref{tbl:exp-ssd} shows the average running time of each approach in answering an SSD query. The query time of HoD is at least an order of magnitude smaller than that of VC-Index. Meanwhile, VC-Index always outperforms EM-BFS, which is consistent with the experimental results reported in previous work \cite{ckcc12}. We omit EM-BFS on USRN and u-UKWeb, since those two graphs are weighted, for which EM-BFS cannot be used to answer SSD queries. Finally, EM-Dijk incurs a larger query overhead than all other methods.

In the next experiment, we demonstrate an application of HoD for efficient graph analysis. In particular, we consider the task of approximating the {\em closeness} for all nodes in a graph $G$, using the algorithm by Eppstein and Wang \cite{ew04}. The algorithm requires executing $k = \ln n/\epsilon^2$ SSD queries from randomly selected source nodes, where $n$ is the number of nodes in $G$ and $\epsilon$ is a parameter that controls the approximation error. Following previous work \cite{ckcc12}, we set $\epsilon = 0.1$.

Table~\ref{tbl:exp-close} shows an estimation of the time required by each method to complete the approximation task. Specifically, we estimate the total processing time of each method as (i) its query time in Table~\ref{tbl:exp-ssd} multiplied by $k$, plus (ii) its preprocessing time (if any). Observe that both EM-BFS and EM-Dijk incur prohibitive overheads -- they require more than a week to finish the approximation task. In contrast, HoD takes at most $2.4$ hours to complete the task, despite that it pays an initial cost for pre-computation. Meanwhile, VC-Index is significantly outperformed by HoD, and it needs around two days to accomplish the task on FB and u-UKWeb.


%\begin{table*}
%	\centering
%		\begin{tabularx}{\textwidth}{c|c|c|c|c}
%		\hline
%		    & USRN & FB & u-BTC & u-UKWeb\\
%			\hline
%			k & 1699 & 1788 & 1659 & 1560\\
%			\hline
%			HoD(M=1G) & 0.9 & 2.0 & 1.3 & 2.4\\
%			\hline
%			VC-index & 13.2 & 51.8 & 6.1 & 43.1\\
%			\hline
%			EM-BFS & $-$ & 231.1 & 182.2 & $-$\\
%            \hline
%            EM-Dijk & 203.3 & 793.3 &389.0 & 240.0 \\
%		\hline
%		\end{tabularx}
%	\caption{Efficiency of computing closeness}
%	\label{tab:Closeness}
%\end{table*}


\begin{table}[t]
\centering
\vspace{-2mm}
\begin{small}
%\tblcapup
\caption{Estimated time for {\em closeness} computation (in hours).}
\label{tbl:exp-close}
\tblup
\begin{tabular}{|c|c|c|c|c|}
\hline
{\bf Method} &   {\bf USRN} & {\bf FB} &  {\bf u-BTC} & {\bf u-UKWeb} \\  \hline
HoD & $\;\,\;\,$0.9 & $\;\,\;\,$2.0 & $\;\,\;\,$1.3 & $\;\,\;\,$2.4 \\ \hline
VC-Index & $\;\,$13.2 & $\;\,$51.8 & $\;\,\;\,$6.1 & $\;\,$43.1 \\ \hline
EM-BFS & $-$ & 231.1 & 182.2 & $-$ \\ \hline
EM-Dijk & 203.2 & 793.3 & 389.0 & 240.0 \\ \hline
\end{tabular}
\tblcapdown
\end{small}
\end{table}

%\begin{table} [t]
%\centering
%\begin{small}
%\caption{Performance of HoD on directed graphs.}
%\label{tbl:exp-directed}
%\tblup
%\begin{tabular}{| c | c | c | c | c | } \hline
%{\bf Dataset} & {\bf Pre-computation} & {\bf Index Size} & {\bf SSD} & {\bf Closeness} \\ \hline
%BTC & 11.4 minutes & $\;\,$2.1 GB & $\;\,$2.6 sec & $\;\,$1.1 hours \\ \hline
%Meme & 1.2 minutes & $\;\,$2.3 GB & $\;\,$1.8 sec & $\;\,$1.3 hours\\ \hline
%UKWeb & 9.2 hours & 72.6 GB & 53.7 sec &  39.7 hours\\ \hline
%\end{tabular}
%%\tblcapup
%%\caption{Performance of HoD on UKWeb and ClueWeb} \label{tbl:exp-web}
%\tblcapdown
%\end{small}
%\end{table}


\begin{table} [t]
\centering
\begin{small}
\caption{Performance of HoD on directed graphs.}
\label{tbl:exp-directed}
\tblup
\begin{tabular}{| c | c | c | c | } \hline
{\bf Dataset} & {\bf Preprocessing} & {\bf Index Size} & {\bf SSD Query Time} \\ \hline
BTC & 11.4 minutes & $\;\,$2.1 GB & $\;\,$2.6 sec \\ \hline
Meme & $\;\,$1.2 minutes & $\;\,$2.3 GB & $\;\,$1.8 sec\\ \hline
UKWeb & 9.2 hours$\;\,$ & 72.6 GB & 53.7 sec \\ \hline
\end{tabular}
%\tblcapup
%\caption{Performance of HoD on UKWeb and ClueWeb} \label{tbl:exp-web}
\tblcapdown
\end{small}
\end{table}

%Following previous work \cite{ckcc12}, we vary the memory available to VC-Index from $2$GB to $4$GB. For HoD, we test it by varying the memory limit from $1$GB to $4$GB. (We are unable to run VC-Index with $1$GB of memory as it reports an error.)
%
%Figure~\ref{fig:exp-undir-pre} shows the pre-computation time of HoD and VC-Index. The preprocessing cost of HoD is consistently smaller than that of VC-Index, by up to two orders of magnitudes. On u-UKWeb, the pre-computation cost of HoD significantly decreases with the increase of memory size. To explain this, recall that the preprocessing algorithm of HoD can only terminate when the size of the core graph fits in the main memory. When the amount of available memory is large, the core graph can more easily fit in the main memory, and hence, HoD stops the preprocessing algorithm early, resulting in a lower pre-computation time. The same phenomenon also occurs on the other three datasets, but the effect is not as pronounced. This is because, for those three datasets, HoD is able to quickly reduce the size of the core graph to a small value, and thus, it can finish preprocessing early even with a small memory size. Figure~\ref{fig:exp-undir-ssd} illustrates the space consumption of HoD and VC-Index. Observe that HoD has a smaller space consumption than VC-Index on all datasets except BTC.
%
%
%
%Figure~\ref{fig:exp-undir-ssd} shows the query time of HoD and VC-Index for SSD queries. For further comparison, we also include the average query time required by an in-memory Dijkstra's algorithm. (Note that all four undirected graphs fit in the $48$GB main memory of our machine.) The query time of HoD is consistently smaller than that of VC-Index by an order of magnitude, and {\em is even smaller than running time of the in-memory Dijkstra's algorithm in all cases}. The efficiency of HoD does not change significantly with the size of the memory. This is because, when HoD processes an SSD query, the major overhead is incurred by the reversed search, which sequentially scans the reversed graph $\GB$. This overhead is not related to the memory size, and hence, a larger memory does not significantly reduce HoD's query time.

%\begin{figure*}[t]
%\begin{small}
%\begin{tabular}{cccc}
%\multicolumn{4}{c}{\hspace{-4mm} \includegraphics[height=2.5mm]{./figs/legend3.eps}} \vspace{1mm} \\
%\hspace{0mm}\includegraphics[height=30mm]{./figs/usa_ssd.eps}
%&
%\hspace{-4mm}\includegraphics[height=30mm]{./figs/face_ssd.eps}
%&
%\hspace{-4mm}\includegraphics[height=30mm]{./figs/btc_ssd.eps}
%&
%\hspace{-4mm}\includegraphics[height=30mm]{./figs/web_ssd.eps} \\
%\hspace{0mm}(a) USRN & \hspace{0mm}(b) FB &
%\hspace{0mm}(c) u-BTC & \hspace{0mm} (d) u-UKWeb
%\end{tabular}
%\end{small}
%\figcapup \caption{Efficiency of SSD queries.} %\figcapdown %\vspace{-1mm} %
%\label{fig:exp-undir-ssd}
%\end{figure*}

%Figure~\ref{fig:exp-undir-ppd} illustrates the query efficiency of HoD, VC-Index, and a {\em bi-directional} version \cite{sc77} of the in-memory Dijkstra's algorithm that is optimized for PPD queries. Again, HoD outperforms VC-Index by a large margin in all cases. HoD also incurs a smaller overhead than the in-memory Dijkstra's algorithm, except for the case of u-UKWeb. Furthermore, the query cost of HoD is significantly lower than its cost for SSD queries. This is expected since the SSD query algorithm of HoD needs to sequentially scan every node in the revered graph, while the PPD query algorithm can avoid such an expensive scan. On USRN and u-UKWeb, the PPD query cost of HoD increases with the memory size. This is because, given sizable memory, HoD tends to generate a large core graph; in turn, a large core graph will result in higher cost for the core search performed by HoD, and hence, the query time of PPD queries increases.

%\begin{figure*}[t]
%\begin{small}
%\begin{tabular}{cccc}
%\multicolumn{4}{c}{\hspace{-4mm} \includegraphics[height=2.5mm]{./figs/legend3.eps}} \vspace{1mm} \\
%\hspace{-1mm}\includegraphics[height=33mm]{./figs/usa_ppd.eps}
%&
%\hspace{-6mm}\includegraphics[height=33mm]{./figs/face_ppd.eps}
%&
%\hspace{-6mm}\includegraphics[height=33mm]{./figs/btc_ppd.eps}
%&
%\hspace{-6mm}\includegraphics[height=33mm]{./figs/web_ppd.eps} \\
%\hspace{0mm}(a) USRN & \hspace{0mm}(b) FB &
%\hspace{0mm}(c) u-BTC & \hspace{0mm} (d) u-UKWeb
%\end{tabular}
%\end{small}
%\figcapup \vspace{-2mm} \caption{Efficiency of PPD queries.} %\figcapdown %\vspace{-1mm} %
%\label{fig:exp-undir-ppd}
%\end{figure*}


\subsection{Results on Directed Graphs} \label{sec:exp-directed}

Our last experiments focus on the three directed graphs: BTC, Meme, and UKWeb. We run HoD on BTC and Meme with 1GB memory, and on UKWeb with 16GB memory, as the enormous size of UKWeb leads to a higher memory requirement. Table~\ref{tbl:exp-directed} shows the preprocessing and space overheads of HoD, as well as its average query time in answering $100$ randomly generated SSD queries on each dataset. (We omit VC-Index, EM-BFS, and EM-Dijk as they do not support directed graphs.) On BTC and Meme, HoD only incurs small pre-computation costs and moderate space consumptions. On UKWeb, the preprocessing, space, and query overheads of HoD are considerably higher, but are still reasonable given that UKWeb contains $30$ times more edges than BTC and Meme do. To our knowledge, this is the first result in the literature that demonstrates practical support for SSD queries on a billion-edge graph. %In addition, we note that the query time of HoD on UKWeb is comparable to the query time of VC-Index on FB and u-UKWeb, which are smaller in size than UKWeb by more than an order of magnitude.




%Figure~\ref{fig:exp-dir-time} illustrates the pre-computation costs of HoD on BTC and Meme, while Figure~\ref{fig:exp-dir-space} shows HoD's space consumptions on the same datasets. (We omit VC-Index since it does not support directed graphs.) The preprocessing and space overheads of HoD do not change significantly with the memory size, as is the case with the experiments in Section~\ref{sec:exp-undirected}.

%Figures \ref{fig:exp-dir-ssd} and \ref{fig:exp-dir-ppd} illustrates the running time of each method for SSD and PPD queries on BTC and Meme. The query performance of HoD is not as good as, but still comparable to the in-memory Dijkstra's algorithm. The reason that Dijkstra's algorithm is so efficient is that, for each node in BTC and Meme, the number of other nodes that can be reached from $s$ tend to be small. Consequently, Dijkstra's algorithm can answer SSD and PPD without travelling a significant portion of the graph, which leads to high efficiency. Nevertheless, for large graphs that do not fit in the memory, HoD remains the only practical approach for distance queries. To demonstrate this, we conduct an experiment on our largest graphs, UKWeb and ClueWeb, which are more than $60$GB and $180$GB, respectively (i.e., they do not fit in the main memory of our machine). We process the datasets with HoD, setting the available memory to $32GB$. (We use a larger memory size here to avoid pro-longed preprocessing time.) Table~\ref{tbl:exp-web} shows the preprocessing time, space consumption, and query time of HoD on UKWeb and ClueWeb. The query time of HoD is higher than in the previous experiments, but is still reasonable given the size of the input data. To the best of our knowledge, we are the first to demonstrate that distance queries on a billion-node graph can be practically supported, even with the memory available is relatively small.




%% [inlined 6-Related/related.tex]
\section{Related Work} \label{sec:related}

As mentioned in Section~\ref{sec:intro}, the existing techniques for I/O-efficient SSD and SSSP queries include VC-Index \cite{ckcc12} and a few methods that adopt Dijkstra's algorithm \cite{mo09,mz03,mz12,m09,ks96}. All of those techniques are exclusively designed for undirected graphs, and they incur significant query overheads, as is shown in our experiments. In contrast, HoD supports both directed and undirected graphs, and it offers high query efficiency along with small costs of pre-computation and space.

%and they do not efficiently support PPD queries. In contrast, our HoD approach works on both directed and undirected graphs, and it is optimized for both SSD and PPD queries.

Other than the aforementioned work, there exists large body of literature on in-memory algorithms for shortest path and distance queries (see \cite{s12,dss09,zmxl13,tsp11} for surveys). The majority of those algorithms focus on two types of queries: (i) {\em point-to-point shortest path (PPSP)} queries, which ask for the shortest path from one node to another, and (ii) {\em point-to-point distance (PPD)} queries, which ask for the length of the shortest path between two given nodes. These two types of queries are closely related to SSD and SSSP queries, in the sense that any SSD (resp.\ SSSP) query can be answered using the results of $n$ PPD (resp.\ PPSP) queries, where $n$ is the number of nodes in the graph. Therefore, it is possible to adopt a solution for PPD (resp.\ PPSP) queries to handle SSD (resp.\ SSSP) queries. Such an adoption, however, incurs significant overheads, especially when $n$ is large. For example, the state-of-the-art solution \cite{adgw11} for PPD queries requires $266$ns to answer a PPD query on the USRN dataset in Section~\ref{sec:exp}. (Note: the solution is not I/O efficient and it requires $25.4$GB memory to handle USRN.) If we use this solution to answer an SSD query on USRN, then we need to execute $24.5$ million PPD queries, which takes roughly $266\textrm{ns} \times 24.5 \times 10^6 = 6.52\textrm{s}$. In contrast, HoD requires only $1.8$s to process an SSD query on USRN, using only $1$GB memory.

Furthermore, almost all existing solutions for PPD and PPSP queries require that the dataset fits in the main memory during pre-computation and/or query processing. This renders them inapplicable for the massive disk-resident graphs considered in this paper. The only exception that we are aware of is a concurrent work by Fu et al.\ \cite{fwcc13}, who propose an I/O-efficient method called {\em IS-Label}. HoD and IS-Label's preprocessing algorithms are similar in spirit, but their index structures and query algorithms are drastically different, as they are designed for different types of queries. In particular, IS-Label focuses on PPD and PPSP queries, and does not efficiently support SSD or SSSP queries.

Finally, we note that previous work \cite{ss05,gss08,dss09} has exploited the idea of augmenting graphs with shortcuts to accelerate PPD and PPSP queries. Our adoption of shortcuts is inspired by previous work \cite{ss05,gss08,dss09}, but it is rather non-trivial due to the facts that (i) we address I/O efficiency under memory-constrained environments, while previous work \cite{ss05,gss08,dss09} focuses on main memory algorithms; (ii) we tackle SSD and SSSP queries instead of PPD and PPSP queries; (iii) we focus on general graphs, while pervious work \cite{ss05,gss08,dss09} considers only road networks (where each node is degree-bounded).



 %with a focus on {\em approximate} PPD queries \cite{w12}, PPD queries on planar graphs \cite{w12},  and PPD queries on road networks \cite{rt10,ss10,gw05,gss08}. Among these algorithms, the one most related to ours is by Goldberg and Werneck \cite{gw05}. Specifically, Goldberg and Werneck devise a disk-based approach for PPD queries on hand-held devices, with an aim to reduce memory consumption and disk I/O during query processing. Nevertheless, the preprocessing algorithm of the approach requires that the dataset can fit in the main memory, which renders the approach inapplicable for large graphs. Instead, Goldberg and Werneck's approach is more suitable for small or medium road networks, which can be (i) preprocessed on a machine with a large memory size and then (ii) transferred to hand-held devices for answering queries.


%
%
% The approach is optimized to reduce memory consumption and disk I/O during query processing, but its preprocessing algorithm requires that the dataset fits in the main memory. This renders the approach inapplicable for large graphs.

%Among these algorithms, the ones most related to ours are by Goldberg and Werneck \cite{gw05} and by Geisberger et al.\ \cite{gss08}. Specifically, Goldberg and Werneck devise a disk-based approach for PPD queries on hand-held devices. The approach is optimized to reduce memory consumption and disk I/O during query processing, but its preprocessing algorithm requires that the dataset fits in the main memory. This renders the approach inapplicable for large graphs.

%Geisberger et al.\ \cite{gss08} propose an indexing structure
%
%
%
%
%
%Compared with HoD, those works are designed only undirected graphs, and they
%
%Talk about the existing solution, with a focus on James' method and CH.
%
%Goodrich's method: Preprocessing requires that the input graph fixes in main memory.
%
%HoD is inspired by CH. Differences: CH cannot handle disk-resident graphs, while HoD can; CH is designed for road networks, while HoD is designed for more general graphs; CH is designed for PPD query, and there is no known algorithm for computing SSD query with CH. Why is it difficult to externalize CH: (i) The preprocessing algorithm of CH requires a heap. (ii) When constructing shortcuts from $s$ to $t$, CH requires computing the shortest path from $s$ to $t$, which is difficult when $G$ is disk-resident. (iii) During query processing, CH adopts a modified version of Dijkstra's algorithm, and the order of node traversal changes with the starting point of the query.
%
%



%% [inlined 8-Conclusions/conclusions.tex]
\section{Conclusions}\label{sec:conclu}

%This paper presents {\em HoD}, a practically efficient index structure for distance queries on disk-resident large graphs. In particular, HoD supports both directed and undirected graphs, and it efficiently handles {\em single source distance (SSD) queries} and {\em point-to-point distance (PPD) queries}. This contrasts the existing methods, which (i) are applicable only on undirected graphs and (ii) do not optimize the efficiency of PPD queries. With extensive experiments on a variety of large graphs, we demonstrate that HoD significantly outperforms the state of the art in terms of query efficiency, space consumption, and pre-computation costs. For future work, we plan to investigate how HoD can be extended to handle dynamic graphs that change with time.

This paper presents HoD, a practically efficient index structure for distance queries on massive disk-resident graphs. In particular, HoD supports both directed and undirected graphs, and it efficiently handles {\em single-source shortest path (SSSP) queries} and {\em single-source distance (SSD) queries} under memory-constrained environments. This contrasts the existing methods, which either (i) require that the dataset fits in the main memory during pre-computation and/or query processing, or (ii) support only undirected graphs. With extensive experiments on a variety of real-world graphs, we demonstrate that HoD significantly outperforms the state of the art in terms of query efficiency, space consumption, and pre-computation costs. For future work, we plan to investigate how HoD can be extended to (i) support point-to-point shortest path and distance queries and (ii) handle dynamic graphs that change with time.



%Some bullshit conclusion and future work.
\balance




\bibliographystyle{abbrv}
%\begin{small}
\bibliography{ref}
%\bibliography{ref}
%\end{small}




%% [inlined Appendix/appendix.tex]
\appendix

\section{Proof of Theorem 1} \label{sec:proofs}

%\begin{definition}[Arch Paths]
%Let $\P$ be a path in the augmented graph $\G$ of HoD. $\P$ is an {\bf arch path}, if and only if...
%\end{definition}
%
%Let $s$ and $v$ be two nodes in the original graph $G$, and let $P$ be the shortest path from $s$ to $t$ in $G$. In what follows, we will show that the augmented graph $\G$ contains an arch path $\P$ from $s$ to $t$, such that $l(\P) = l(P)$. In particular, we derive $\P$ from $P$ using an iterative procedure as follows.
%
%At each iteration, we consider the lowest unchecked vertex $w$ on $P(s,t)$, i.e. $rank(w)\leq rank(u),u\in P(s,t)$. If $w$ satisfies either of the following termination conditions, we set $P(s,t)$ as the canonical path $P_c(s,t)$ and terminate. The conditions are: a) $w\in G_c$, b)$w$ is the only unchecked vertex in $P(s,t)$. Otherwise, we check the neighbors $u$ and $v$ of $w$ on $P(s,t)$, suppose $u$ is closer to $s$ while $v$ is closer to $t$.
%
%Case 1: $rank(u)>rank(w)$ and $rank(v)>rank(w)$, recall the preprocessing algorithm, when we remove $w$, we would create a shortcut between $u$ and $v$ unless there exists another path $S(u,v)$ in $G_{rank(w)}$ with equal length (here $S(u,v)$ cannot be shorter, as $P(s,t)$ is the shortest path). If $S(u,v)$ exists, we replace the segment between $u$ and $v$ of $P(s,t)$ by $S(u,v)$. The vertices on $S(u,v)$ must have higher rank than $w$ as $S(u,v)\subseteq  G_{rank(w)}$. Otherwise, we replace the segment between $u$ and $v$ of $P(s,t)$ by the shortcut $\langle u,v\rangle$. Here we have two observations here: a) neither of these replacements will change the length of $P(s,t)$, so $P(s,t)$ is still a shortest path between $s$ and $t$. b) Once $w$ is removed from $P(s,t)$, it won't be brought back to $P(s,t)$ again.
%
%Case 2: $rank(u)<rank(w)$ or $w = s$, then we have $rank(w)<rank(v)$, otherwise, $w$ is the only unchecked vertex in $P(s,t)$, which contradicts to termination condition b). Recall the preprocessing algorithm, when we remove $w$, we would push shortcut $\langle w,v \rangle$ into $\G_f$. We mark $w$ as checked, and $P(s,t)$ remains the same.
%
%Case 3: $rank(w)>rank(v)$ or $w=t$, similar to case 2, we have $rank(u)>rank(w)$, and shortcut $\langle u,w \rangle$ is also in $\G_b$. Similarly we mark $w$ as checked, and $P(s,t)$ remains the same.
%
%In any of the three cases, a) the length of $P(s,t)$ remains the same, b) $w$ is removed or marked, so our process will always terminate and generate $P_c(s,t)$ which is still a shortest path between $s$ and $t$. In general, $P_c(s,t)$ can be divided into three segments: $S_f$, shortcuts connecting $s$ to the first vertex in $G_c$, namely $u$; $S_c$, shortcuts connecting $u$ to the last vertex in $G_c$, namely $v$; and $S_b$, shortcuts connecting $v$ to $t$. For certain $s,t$, some segments may not exist, e.g. if $s\in G_c$, $S_f=\emptyset$; if $t\in G_c$, $S_b=\emptyset$; if none of the vertices in $P_c(s,t)$ is in $G_c$, we set $u,v$ as the vertex in $P_c(s,t)$ with the highest rank, in this case, $S_c=\emptyset$.
%
%For any $P_c(s,t)$, if $S_f\neq \emptyset$, our forward search ensures $\kappa_f(u)\leq |S_f|$; if $S_c\neq \emptyset$, our core search ensures $\kappa_f(v)-\kappa_f(u)\leq |S_c|$; and finally if $S_b\neq \emptyset$, our backward search ensures $\kappa_f(t)-\kappa_f(v)\leq |S_b|$. By sum up the in-equations above, we have $\kappa_f(t)\leq |P_c(s,t)|$, at the same time, as $P_c(s,t)$ is a shortest path between $s$ and $t$, we proved the correctness of our query process.

%To prove the correctness of Theorem~\ref{thrm:sssp-correct}, it suffices to show that for node $t$ in $G$, the query algorithm of HoD correctly computes the distance from the source node $s$ to $t$.

Our proof of Theorem~\ref{thrm:sssp-correct} utilizes the concepts of {\em rank sequences} and {\em arch paths}, defined as follows.
\begin{definition}[Rank Sequence and Arch Path] \label{def:proof-rank}
Let $\P$ be a path in the augmented graph $\G$ of HoD, such that $\P$ contains $k$ nodes. The {\bf rank sequence} of $\P$ is a sequence of integers $\langle r_1, r_2, \ldots, r_k\rangle$, such that $r_i$ ($i \in [1, k]$) equals the rank of the $i$-th node in $\P$. $\P$ is an {\bf arch path}, if its rank sequence can be divided into three subsequences $\langle r_1, \ldots, r_x\rangle$, $\langle r_x, \ldots, r_y\rangle$, and $\langle r_y, \ldots, r_k\rangle$, such that
\begin{enumerate} [itemsep = 0.5mm, topsep = 2mm]
\item $r_1 < r_2 < \ldots < r_x$, and

\item $r_x = r_{x+1} = \ldots = r_y$, and

\item $r_y > r_{y+1} > \ldots > r_k$. \done
\end{enumerate}
\end{definition}
For instance, consider the path $\P = \langle v_1, v_9, v_{10}, v_8, v_4\rangle$ in Figure~\ref{fig:over-steps}e. The path's rank sequence is $\langle 1, 4, 4, 3, 2\rangle$. This rank sequence can be divided into three subsequences $\langle 1, 4\rangle$, $\langle 4, 4\rangle$, and $\langle 4, 3, 2\rangle$. By Definition~\ref{def:proof-rank}, $\P$ is an arch path.

Let $s$ and $t$ be any two nodes in the original graph $G$, and $P$ be the shortest path from $s$ to $t$ in $G$. If there exist multiple shortest paths from $s$ to $t$, we choose $P$ to be a path where the highest-rank node ranks no lower than any other node on any other shortest path from $s$ to $t$. In the following, we will prove three propositions:
\begin{itemize}
\item {\bf Proposition 1:} For any path $\P$ from $s$ to $t$ in the augmented graph $\G$, its length is no shorter than $P$'s.

\item {\bf Proposition 2:} There exists an arch path $\PA$ in $\G$, such that $l(\PA) = l(P)$, i.e., $\PA$ and $P$ have the same length.

\item {\bf Proposition 3:} When HoD's answers an SSD query from $s$, it will traverse a path no longer than $\PA$.
\end{itemize}
The combination of the above three propositions will establish Theorem~\ref{thrm:sssp-correct}.

\header
{\bf Proof of Proposition 1.} If all edges in $\P$ appear in the original graph $G$, the proposition trivially holds since any path from $s$ to $t$ in $G$ should be no shorter than $P$. In the following, we consider only the case when $P'$ contains at least one shortcut, and we will prove that there exists a path $P'$ in $G$, such that $l(\P) = l(P') \ge l(P)$.

Assume without loss of generality that $\P$ consists of a sequence of $k$ nodes $\langle v_1, v_2, \ldots, v_k\rangle$, where $v_1 = s$ and $v_k = t$. Further assume that $\langle v_i, v_{i+1}\rangle$ ($i \in [1, k]$) is a shortcut. By the preprocessing algorithm of HoD, this shortcut must be constructed when HoD removes a certain node $v$ from the reduced graph, such that $v_i$ and $v_{i+1}$ are incoming and outgoing neighbors of $v$, respectively. This indicates that the augmented graph $\G$ must contain two edges $\langle v_i, v \rangle$ and $\langle v, v_{i+1} \rangle$, such that their total length equals $l(\langle v_i, v_{i+1}\rangle)$. Given those two edges, we change $\P$ by replacing $\langle v_i, v_{i+1}\rangle$ with a two-hop path $\langle v_i, v, v_{i+1}\rangle$. This results in a modified path from $s$ to $t$ in $\G$ that has the same length with the original $\P$.

If the modified $\P$ does not contain any shortcut, then it is a path in $G$, and hence, $l(\P) \ge l(P)$. On the other hand, if the modified $\P$ contains any shortcut, then we can replace the shortcut with a two-hop path in $\G$, as in the previous case for the shortcut $\langle v_i, v_{i+1}\rangle$. By recursively applying this replacement procedure on the shortcuts in $\P$, we can transform $\P$ into a path $P'$ that (i) goes from $s$ to $t$, (ii) contains only edges in $G$, and (iii) has the same length with $\P$. Given that $P$ is the shortest path from $s$ to $t$ in $G$, we have $l(P) \le l(P') = l(\P)$, which completes the proof.

\header
{\bf Proof of Proposition 2.} The proposition trivially holds if $P$ itself is an arch path in $\G$. In the following, we assume that $P$ is not an arch path, and we show that $P$ can be transformed into a arch path $\PA$ in $\G$ with the same length.

Let $r_{max}$ be the highest rank of the nodes in $P$. Let $v_x$ and $v_y$ be the first and last nodes in $P$ whose ranks equal $r_{max}$. We divide $P$ into three subsequences $P_1$, $P_2$, and $P_3$, such that
\begin{enumerate}
\item $P_1$ is the sequence of nodes in $P$ before $v_x$ (including $v_x$).
\item $P_2$ is the sequence of nodes in $P$ between $v_x$ and $v_y$ (including $v_x$ and $v_y$).
\item $P_3$ is the sequence of nodes in $P$ after $v_y$ (including $v_y$).
\end{enumerate}

Let us first consider $P_1$. We say that a node $v$ in $P_1$ is {\em pit}, if $v$ ranks no higher than the node that immediately precedes $v$ in $P_1$. Among the pits in $P_1$, let $v'$ be one with the lowest rank. Let $u$ (resp.\ $w$) be the node that immediately precedes (resp.\ follows) $v'$ in $P_1$. By the preprocessing algorithm of HoD, when HoD removes $v'$ from the reduce graph, it would generate a candidate edge $e_1 = \langle u, w\rangle$, such that the edge has the same length with the two-hop path $\langle u, v', w\rangle$. If $e_1$ is in $\G$ (i.e., it is retained by HoD during preprocessing), then we transform $P_1$ into another path in $\G$, by using $e_1$ to replace the two-hop path $\langle u, v', w\rangle$ in $P_1$. The resulting path has the same length with $P_1$, and it has one less pits than $P_1$. 

On the other hand, if $e_1$ is not in $\G$, there are two possibilities:
\begin{enumerate} 
\item There already exists an edge $e_2 = \langle u, w\rangle$ in the reduced graph, such that $l(e_2) = l(e_1)$. (Note that $l(e_2) < l(e_1)$ is impossible given Proposition 1 and the fact that $P_1$ is the shortest path from $s$ to $v_x$ in $G$.) In that case, $e_2$ must appear in $\G$. Therefore, if we modify $P_1$ by using $e_2$ to replace the two-hop path $\langle u, v', w\rangle$ in $P_1$, we can still obtain a modified path that retains the length of $P_1$ but contains one less pits.
    
\item There exists a two-hop path $\langle u, v^\diamond, w\rangle$ in the reduced graph, such that (i) $l(e_3) = l(e_1)$, and (ii) $v^\diamond$ has a higher rank than $v'$. (Note that $l(e_3) < l(e_1)$ cannot occur due to Proposition 1 and the fact that $P_1$ is the shortest path from $s$ to $v_x$ in $G$.) In that case, we transform $P_1$ by replacing $v'$ with $v^\diamond$. This may not decrease the number of pits in $P_1$, but it retains the length of $P_1$ and replaces a node in $P_1$ with a higher rank node. 
\end{enumerate}
In summary, the above transformation procedure preserves the length of $P_1$, and it either (i) reduces the number of pits in $P_1$ or (ii) substitute a node in $P_1$ with a higher-rank node. Given that the ranks of nodes are bounded, if we recursively apply the procedure on the lowest-rank pit in $P_1$, eventually we should obtain a path $P_1^*$ in $\G$ without any pit, such that (i) $l(P_1^*) = l(P_1)$, and (ii) $P_1^*$ starts at $s$ and ends at $v_x$. In that case, the rank sequence of $P_1^*$ must be an ascending sequence.

Now consider $P_2$ and $P_3$. We say that a node in $P_2$ is a pit if its rank is smaller than $r_{max}$, and we define a node $v$ in $P_3$ as a pit if $v$ ranks no higher than the node that immediately follows $v$ in $P_3$. By applying the same transformation procedure as in the case of $P_1$, we can convert $P_2$ and $P_3$ into two paths $P_2^*$ and $P_3^*$ in $\G$, such that (i) $P_2^*$ and $P_3^*$ have no pit, (ii) $l(P_2^*) = l(P_2)$ and $l(P_3^*) = l(P_3)$, (iii) $P_2^*$ starts at $v_x$ and ends at $v_y$, and (iii) $P_3^*$ starts at $v_y$ and ends at $t$. It can be verified that all nodes in $P_2^*$ should have a rank $r_{max}$, and $P_3^*$'s rank sequence should be a descending sequence.

Let $\PA$ be a path in $\G$ obtained by concatenating $P_1^*$, $P_2^*$, and $P_3^*$. By Definition~\ref{def:proof-rank}, $\PA$ is an arch path, which completes the proof.


\header
{\bf Proof of Proposition 3.} Let $v_x$, $v_y$, $\PA$, $P_1^*$, $P_2^*$, and $P_3^*$ be as defined in the proof of Proposition 2. Without loss of generality, assume that each of $P_1^*$, $P_2^*$, and $P_3^*$ contains at least two nodes. We will prove the proposition by showing that, given an SSD query from $s$, the query processing algorithm of HoD will traverse (i) a path from $s$ to $v_x$ that is no longer than $P_1^*$, (ii) a path from $v_x$ to $v_y$ that is no longer than $P_2^*$, and (iii) a path from $v_y$ to $t$ that is no longer than $P_3^*$.

Recall that HoD's query algorithm consists of three phases: a forward search in the forward graph $\GF$, followed by a core search in the core graph $\GC$, and finally a backward search in the backward graph $\GB$. The forward search is a variant of Dijkstra's algorithm that follows only the outgoing edges whose endpoints rank higher than the starting points. Since the rank sequence of $P_1^*$ is in ascending order, $P_1^*$ should be in the search space of the forward search. Furthermore, by Proposition 1 and the construction of $P_1^*$, $\GF$ does not contain any path from $s$ to $v_x$ that is shorter than $P_1^*$. Therefore, when the forward search terminates, HoD should either identify $P_1^*$ as the shortest path from $s$ to $v_x$ in $\GF$, or identify another path from $s$ to $v_x$ that is no longer than $P_1^*$. In either case, HoD will correctly derive $dist(s, v_x)$, i.e., the distance from $s$ to $v_x$ in $G$.

Now consider $P_2^*$, where each node has the same rank. By the preprocessing algorithm of HoD, all nodes in $P_2^*$ must be in the core graph $\GC$, since any node not in the core graph can only have outgoing edges to higher-rank nodes (see Section~\ref{sec:pre}). Meanwhile, recall that the core search of HoD is a continuation of the forward search in $\GC$, using Dijkstra's algorithm. By the correctness of Dijkstra's algorithm and the fact that $P_2^*$ is the shortest path from $v_x$ to $v_y$ in $\GC$, the core search of HoD should traverse a path from $v_x$ to $v_y$ that is no longer than $P_2^*$.

It remains to prove that HoD's backward search will traverse $P_3^*$. Assume without loss of generality that $P_3^*$ contains a sequence of $k$ nodes $\langle v_1, v_2, \ldots, v_k$, where $v_1 = v_x$ and $v_k = t$. Recall that the backward search of HoD examines nodes in descending order of their ranks. Since $P_3^*$ has a descending rank sequence, the backward search of HoD should examine $v_i$ before $v_{i+1}$, for any $i \in [1, k-1]$.

We will prove by induction that, for any $v_i$ ($i \in [1, k]$), the backward search can correctly derive $dist(s, v_i)$. First, given that $P_1^*$ and $P_2^*$ have been traversed by HoD before the backward search, HoD should be able to compute the precise value of $\d(s, v_1)$ when after it visits $v_1$. Now assume that, after examining $v_j$ ($j \in [1, k-1]$), HoD correctly calculates $\d(s, v_j)$. Then, when HoD inspects $v_{j+1}$, it would identify $dist(s, v_j) + l(\langle v_j, v_{j+1}\rangle)$ as the length of a path from $s$ to $v_{j+1}$. Given that $P_1^*$ is a shortest path from $s$ to $t$ in $\G$ and that $v_{j+1}$ immediately follows $v_j$ on $P_1^*$, we have $dist(s, v_{j+1}) = dist(s, v_j) + l(\langle v_j, v_{j+1}\rangle)$. This indicates that HoD will correctly derive $dist(s, v_{j+1})$, which completes the proof.



%\begin{definition}[Arch Paths] \label{def:proof-arch}
%Let $\P$ be a path in the augmented graph $\G$ of HoD, such that . $\P$ is an {\bf arch path}, if and only if $\P$
%
%
%can be sequentially divided into three parts, namely $\P_f,\P_c,\P_b$, where the ranks of the nodes in $\P_f$ are in the ascending order, nodes in $\P_c$ all have top ranks, and the ranks of the nodes in $\P_b$ are in the descending order. The division of $\P$ is consistent with the division of $\G$, where $\G_f,\G_c,\G_b$ are the subgraphs that are used in the three steps, namely forward search, core search and backward search of the query algorithm of HoD respectively. Therefore, we have $\P_f\subseteq \G_f,\P_c\subseteq \G_c,\P_b\subseteq \G_b$.
%\end{definition}



%Let $t$ be an arbitrary node in $G$. In what follows, we will prove Theorem~\ref{thrm:sssp-correct} by showing that the query algorithm of HoD can correctly compute the distance from the source node $s$ to $t$. In particular, we will show that HoD can identify a path from $s$ to $t$ in the augmented graph $\G$, such that the path has the same length with the shortest path $P$ from $s$ to $t$ in the original graph $G$. Towards this end, we first introduce the concept of {\em arch paths} as follows.
%
%
%
%We will first prove that the augmented graph $\G$ contains an arch path $\P$ from $s$ to $t$, such that $l(\P) = l(P)$. In particular, we derive $\P$ from $P$ using an iterative procedure as follows.
%
%\begin{algorithm}
%  \caption{Procedure for generating $\P$}
%
%  \begin{algorithmic}
%    \State $\P_f\leftarrow \emptyset, \P_c\leftarrow \emptyset, \P_b\leftarrow\emptyset$
%    \State $P_0\leftarrow P,k\leftarrow 0$
%
%    \While{ $P_k\neq\emptyset$ }
%
%    \State $w\leftarrow \operatorname{argmin}\{ rank(u)|u\in P_k\}$
%
%    \If {$w$ has top rank}
%    \State break
%    \EndIf
%
%    \If {$w$ is the starting node of $P_k$}
%        \State $v\leftarrow$ neighbor of $w$ on $P_k$
%        \State $P_{k+1}\leftarrow P_k-\langle w,v \rangle$
%        \State $\P_f\leftarrow \P_f\cup\langle w,v \rangle$
%    \Else \If {$w$ is the ending node of $P$}
%                \State $u\leftarrow$ neighbor of $w$ in $P_k$
%                \State $P_{k+1}\leftarrow P_k-\langle u,w\rangle$
%                \State $\P_b\leftarrow \langle u,w\rangle\cup \P_b$
%          \Else
%                \State $u,v\leftarrow$ neighbors of $w$ on $P_k$
%                \If {$\langle u,v\rangle\in\G$}
%                \State $P_{k+1}\leftarrow (P_k-\langle u,w\rangle-\langle w,v\rangle)\cup\langle u,v\rangle$
%                \Else
%                \State there exists a path $S$ from $u$ to $v$ in $\G$
%                \State $P_{k+1}\leftarrow (P_k-\langle u,w\rangle-\langle w,v\rangle)\cup S$
%                \EndIf
%          \EndIf
%    \EndIf
%
%    \State $k\leftarrow k+1$
%
%    \EndWhile
%
%    \State $\P_c\leftarrow P_k$
%    \State $\P\leftarrow \P_f\cup \P_c\cup \P_b$
%  \end{algorithmic}
%\end{algorithm}
%
%%Initially, we set $\P_f=\P_c=\P_b=\emptyset$. During the procedure, $\P_f,\P_b$ collect those edges that are removed from $P$, and when the procedure terminates, we set $\P_c=P$.
%%
%%At each iteration, if $P=\emptyset$, the procedure terminates. Otherwise, we consider the node $w$ on $P$ with the lowest rank(i.e. $rank(w)=rank(P)=\min\{rank(u)|u\in P\}$), if $w\in \G_c$, the procedure terminates, otherwise, we discuss three cases based on the position of $w$ on $P$.
%%
%%Case 1: $w$ is the start node of $P$, consider the adjacent node of $w$ on $P$, denoted as $v$, we have $rank(v)>rank(w)$. In this case, we remove $\langle w,v\rangle$ from $P$ and connect it to the end of $\P_f$. Recall the preprocessing algorithm, we know that $\langle w,v \rangle \in \G_f$.
%%
%%Case 2: $w$ is the end node of $P$, consider the adjacent node of $w$ on $P$, denoted as $u$, we have $rank(u)>rank(w)$. In this case, we remove $\langle u,w\rangle$ from $P$ and connect it to the start of $\P_b$. Recall the preprocessing algorithm, we know that $\langle u,w \rangle \in \G_b$.
%%
%%Case 3: In this case, $w$ has two adjacent nodes $u$(closer to the start of $P$) and $v$(closer to the end of $P$ resp.), and we have $rank(u)>rank(w)$ and $rank(v)>rank(w)$, recall the preprocessing algorithm, when we remove $w$, we would create a shortcut $\langle u,v\rangle$ with the length of $l(\langle u,w\rangle)+l(\langle w,v\rangle)$ in $G_{rank(w)}$, unless there exists another path $S(u,v)$ in $G_{rank(w)}$ with smaller or equal length (here $S(u,v)$ cannot be shorter, as $P$ is the shortest path). If $S(u,v)$ exists, we replace the segment between $u$ and $v$ of $P$ by $S(u,v)$. The nodes on $S(u,v)$ must have higher rank than $w$ as $S(u,v)\subseteq  G_{rank(w)}$. Otherwise, we replace the segment between $u$ and $v$ of $P$ by the shortcut $\langle u,v\rangle$. Noticed that neither of the replacements changes the length of $P$, or brings in nodes with rank lower than equal to $rank(P)$ which means once $w$ and its corresponding edges are replaced in $P$, we won't bring $w$ back to $P$ any more.
%
%According to the procedure above, we have the following observations:
%
%a) $\P_f\subseteq \G_f$($\P_b\subseteq \G_b$ respectively). $\P_f$ collects edges when $w$ is the starting node of $P$, by the property of HoD, we have $rank(v)>rank(w)$ which means $\langle w,v\rangle\in \G_f$, therefore, $\P_f\subseteq \G_f$. Similarly, we have $\P_b\subseteq \G_b$.
%
%b) $\P_c\subseteq \G_c$. If our procedure terminates when $P_k=\emptyset$, it is trivial to show that $\P_c=P_k=\emptyset\subseteq \G_c$. Otherwise, we know that all the nodes on $P_k$ have top ranks, which means $\P_c=P_k\subseteq \G_c$.
%
%c) $l(\P)=l(P)$. In the beginning, we have $l(\P_f)+l(\P_b)+l(P_0)=0+0+l(P)=l(P)$, so the statement of $l(\P_f)+l(\P_b)+l(P_k)=l(P)$ is true when $k=0$, now we prove if the statement is true when $k=i$, it's still true when $k=i+1$. Firstly, when $w$ is the starting node(the ending node) of $P_k$, we cut one edge from $P_k$, and attach it to $\P_f$($\P_b$ resp.), it's easy to see that the statement is still true when $k=i+1$. Secondly, when $w$ is in the middle of $P_k$, if $\langle u,v\rangle\in\G$, HoD ensures that $l(\langle u,v\rangle)=l(\langle u,w\rangle)+l(\langle w,v\rangle)$, so in the case, $l(P_{k+1})=l(P_k)$. Otherwise, there exists a path $S$ from $u$ to $v$ in $\G$ with equal or less length than $l(\langle u,w\rangle)+l(\langle w,v\rangle)$, however, $P$ is the shortest path which means $S$ cannot be shorter, therefore, we have $l(S)=l(\langle u,w\rangle)+l(\langle w,v\rangle))$, so in this case, we also have $l(P_{k+1})=l(P_k)$. As neither $\P_f$ nor $\P_b$ changes, the statement is still true when $k=i+1$. So no matter which position $w$ has on $P_k$, we always have $l(\P_f)+l(\P_b)+l(P_k)=l(P)$. In the end of the procedure, we set $\P_c$ as $P_k$ which gives us $l(\P)=l(\P_f)+l(\P_c)+l(\P_b)=l(\P_f)+l(P_k)+l(\P_b)=l(P)$.
%
%For any given nodes $s$ and $t$ with the shortest path $P$ from $s$ to $t$ in $G$, the procedure above create a path $\P$ in $\G$ from $s$ to $t$. By observations a) and b), we know $\P$ is an arch path. As the forward search, core search, and backward search of the query algorithm of HoD identify the three parts of $\P$ respectively, we know that the distance computed by the query algorithm of HoD equals to the length of $\P$, at the same time, by observation c), we know that $\P$ has the same length of $P$ which is the shortest path from $s$ to $t$ in the original graph $G$, so the distance computed by the query algorithm of HoD equals to the length of $P$ which means the query algorithm of HoD is correct.



%============================================



\end{sloppy}
\end{document}

